% P005: exact first-order regret for stochastic bandits with bounded kurtosis.
% Written from templates/paper/paper.tex (AGENTS.md, section 5; docs/writing.md, sections 2-2a, the
% owner's house style of 2026-11-16). Two builds of one source:
%   reader:    latexmk -pdf -interaction=nonstopmode -halt-on-error -cd papers/P005-bounded-kurtosis/paper.tex
%   blueprint: latexmk -pdf -interaction=nonstopmode -halt-on-error -jobname=paper-blueprint
%                -usepretex='\def\IncognitaBlueprint{}' -cd papers/P005-bounded-kurtosis/paper.tex
% Blueprint conventions: every node has a \label (def:, lem:, prop:, thm:, cor:, cite:); statements
% start with \uses when they use definitions; every proof starts with an exhaustive \uses; a proof away
% from its statement has \proves. Never \lean, \leanok or \mathlibok in this file.
\documentclass[11pt]{article}
\usepackage[T1]{fontenc}
\usepackage[utf8]{inputenc}
\usepackage{lmodern}
\usepackage[a4paper,margin=2.5cm]{geometry}
\usepackage{amsmath,amssymb,mathtools}
\usepackage{booktabs}
\usepackage{algorithm}
\usepackage{algpseudocode}
\usepackage{enumitem}
\usepackage{xcolor}
\usepackage[colorlinks=true,linkcolor=blue!50!black,citecolor=blue!50!black,urlcolor=blue!50!black]{hyperref}
\usepackage[round]{natbib}
\usepackage{../../templates/paper/blueprint-macros}

% ---------------------------------------------------------------- macros (\newcommand only)
\newcommand{\E}{\mathbb{E}}
\newcommand{\Prob}{\mathbb{P}}
\newcommand{\R}{\mathbb{R}}
\newcommand{\ind}{\mathbf{1}}
\newcommand{\Pspace}{\term{w:law}{\mathcal{P}(\mathbb{R})}}
% Navigation (house style): an anchor at each definition, a macro per symbol linking to it.
\newcommand{\dt}[1]{\hypertarget{#1}{}}
\newcommand{\sy}[2]{{\hyperlink{#1}{#2}}}
\newcommand{\term}[2]{\hyperlink{#1}{#2}}
\newcommand{\mean}{\sy{s:mean}{m}}
\newcommand{\sd}{\sy{s:var}{\sigma}}
\newcommand{\kurt}{\sy{s:kurt}{\mathrm{kurt}}}
\newcommand{\KL}{\sy{s:KL}{\mathrm{KL}}}
\newcommand{\kl}{\sy{s:kl}{\mathrm{kl}}}
\newcommand{\Karms}{\sy{s:Karms}{K}}
\newcommand{\marm}{\sy{s:ma}{m}}
\newcommand{\mstar}{\sy{s:mstar}{m^{\star}}}
\newcommand{\gap}{\sy{s:gap}{\Delta}}
\newcommand{\sda}{\sy{s:sa}{\sigma}}
\newcommand{\Xr}{\sy{s:X}{X}}
\newcommand{\Arm}{\sy{s:At}{A}}
\newcommand{\Kinf}{\sy{s:KinfD}{\mathcal{K}_{\inf}}}
\newcommand{\Hk}{\sy{s:Hk}{\mathcal{H}_{\kappa}}}
\newcommand{\Hlt}{\sy{s:Hlt}{\mathcal{H}^{<}_{\kappa}}}
\newcommand{\Kc}{\sy{s:K}{\mathcal{K}}}
\newcommand{\Kge}{\sy{s:Kge}{\mathcal{K}^{\ge}}}
\newcommand{\phik}{\sy{s:phi}{\varphi}}
\newcommand{\Gk}{\sy{s:G}{\mathcal{G}}}
\newcommand{\KG}{\sy{s:KG}{\mathcal{K}_{\mathcal{G}}}}
\newcommand{\tauc}{\sy{s:tau}{\tau}}
\newcommand{\hl}{\sy{s:h}{h}}
\newcommand{\gam}{\sy{s:gamma}{\gamma}}
\newcommand{\Lam}{\sy{s:Lambda}{\Lambda}}
\newcommand{\Lw}{\sy{s:L}{L}}
\newcommand{\qk}{\sy{s:q}{q_{\kappa}}}
\newcommand{\Kinfty}{\sy{s:Kinfty}{\mathcal{K}_{\infty}}}
\newcommand{\Np}{\sy{s:N}{N}}
\newcommand{\Pe}{\sy{s:Pe}{\hat{P}}}
\newcommand{\Reg}{\sy{s:R}{R}}
\newcommand{\Cst}{\sy{s:Cstar}{C^{\star}}}
\newcommand{\thr}{\sy{s:g}{g}}
\newcommand{\Ix}{\sy{s:U}{U}}
\newcommand{\Ka}{\sy{s:Ka}{\mathcal{K}_{a}}}
\newcommand{\LT}{\sy{s:lT}{\ell_{T}}}
\newcommand{\nT}{\sy{s:nT}{n_{T}}}
\newcommand{\kep}{\sy{s:keps}{k_{\varepsilon}}}
\newcommand{\Ca}{\sy{s:Ca}{C_{a}}}
\newcommand{\algo}{\hyperref[alg:kucb]{\texttt{KL-UCB-Kurt}}}

\title{Exact first-order regret for stochastic bandits with bounded kurtosis}
\author{claude-opus-5-5, Incognita agents\thanks{Written by AI agents of the Incognita project
  (\url{https://github.com/prevostad/Incognita}) and refereed by two AI agents; its certification
  is under the date.}}
\date{2026-11-16}

\begin{document}
\maketitle

\begin{abstract}
We consider stochastic multi-armed bandits whose reward laws have a finite fourth moment, a positive
variance and a kurtosis at most a known number $\kappa>1$. The class is invariant under changes of
location and scale, and the best known regret bounds on it held only up to universal constant factors.
We show that KL-UCB run with the information function $\mathcal{K}_{\inf}$ of this class, an anytime index
policy, satisfies $\lim_{T}R_T/\log T=\sum_a\Delta_a/\mathcal{K}_{\inf}(\nu_a,m^\star)$ on every instance. Since this is
the constant of the Lai--Robbins and Burnetas--Katehakis lower bound for consistent algorithms, it is
the exact first-order constant of the class. We describe $\mathcal{K}_{\inf}$ by a dual formula valid for every
law, prove its lower semicontinuity and positivity, and find its largest value over laws with a finite
fourth moment, $\mathcal{K}_{\infty}(\kappa)=\log\bigl(2/(1+\sqrt{(\kappa-1)/(\kappa+3)})\bigr)$: the logarithm of the inverse of the largest atom
of a law of kurtosis at most $\kappa$, reached in the limit of large gaps. As
$\kappa<1/\mathcal{K}_{\infty}(\kappa)<\kappa+2$, each suboptimal arm costs more than $\kappa\Delta_a$ per $\log T$. The proof
convexifies the kurtosis constraint by optimising over location and scale, and bounds the escape of
near-optimal alternatives along the affine group by the largest atom. The second-order term is
existential, and the index is not given a certified computation.
\end{abstract}

% ================================================================ main text
\section{Introduction}\label{sec:intro}

A stochastic bandit with finitely many arms is played over rounds $t=1,2,\dots$: the learner pulls an
arm and observes a reward drawn from the law of that arm, independently of the past. Its regret after
$T$ rounds is $R_T=\sum_a\Delta_a\E[N_a(T)]$, where $\Delta_a$ is the gap between the best mean and the
mean of arm $a$ and $N_a(T)$ the number of pulls of $a$. When the learner knows a model $\mathcal D$
that contains every arm's law, the Lai--Robbins and Burnetas--Katehakis lower bound says that every
algorithm consistent over $\mathcal D$ satisfies
\begin{equation}\label{eq:intro-lb}
  \liminf_{T\to\infty}\frac{R_T}{\log T}\ \ge\ \sum_{a:\Delta_a>0}\frac{\Delta_a}{\mathcal{K}_{\inf}(\nu_a,m^\star;\mathcal D)},
  \qquad
  \mathcal{K}_{\inf}(P,\mu;\mathcal D):=\inf\{\KL(P,Q):Q\in\mathcal D,\ m(Q)>\mu\}.
\end{equation}
Algorithms attaining this constant are known for one-parameter exponential families, for laws with a
bounded or semi-bounded support, for Gaussian laws with unknown variance, and for laws with a bounded
moment $\E|X|^{1+\beta}\le B$.

All these models carry a scale: a support, a variance or a moment bound fixed in advance. \citet{L17}
proposed a model with no scale at all: laws whose kurtosis
$\E[(X-m)^4]/\sigma^4$ is at most a known number $\kappa$. The class is invariant under every affine map
$x\mapsto c+sx$ with $s>0$, the regret of a good algorithm should be too, and a bound on the kurtosis is
what is left of the Gaussian assumption when the shape is unknown. Lattimore gave a scale-free algorithm
with $\limsup_TR_T/\log T\le C\sum_a\Delta_a(\kappa-1+\sigma_a^2/\Delta_a^2)$ for a universal constant $C$,
and a lower bound of the same form for $\kappa\ge7/2$, also up to constants. He left the optimal form of
the regret open, suggesting an implicit analysis that would show whether the lower bound
\eqref{eq:intro-lb} can be attained on this class \citep[Section~4]{L17}; the textbook of
\citet[Section~16.4]{LS} records the class as understood only up to constant factors.

We answer this question at first order. Let $\mathcal H_\kappa$ be the laws with a finite fourth moment, a
positive variance and kurtosis at most $\kappa$, and let $\mathcal K(P,\mu)$ be the information function
$\mathcal{K}_{\inf}(P,\mu;\mathcal H_\kappa)$. We show that KL-UCB run with $\mathcal K$ satisfies, on every
instance of $\mathcal H_\kappa$,
\begin{equation}\label{eq:intro-main}
  \lim_{T\to\infty}\frac{R_T}{\log T}=\sum_{a:\Delta_a>0}\frac{\Delta_a}{\mathcal K(\nu_a,m^\star)},
\end{equation}
so the constant of \eqref{eq:intro-lb} is the exact first-order constant of the class (Theorem~\ref{thm:main} and
Corollary~\ref{cor:exact}). The constant is
not in closed form, but we locate it: $0<\mathcal K(\nu_a,m^\star)\le\mathcal K_\infty(\kappa)$, with
\begin{equation}\label{eq:intro-kinfty}
  \mathcal K_\infty(\kappa):=\log\frac{2}{1+\sqrt{(\kappa-1)/(\kappa+3)}},\qquad
  \kappa<\frac1{\mathcal K_\infty(\kappa)}<\kappa+2,\qquad
  \frac1{\mathcal K_\infty(\kappa)}=\kappa+\frac32+O\Bigl(\frac1\kappa\Bigr),
\end{equation}
and $\mathcal K(\nu_a,m^\star)\to\mathcal K_\infty(\kappa)$ when $\Delta_a/\sigma_a\to\infty$ for a fixed standardised law
(the law of $(X-m_a)/\sigma_a$ with $X\sim\nu_a$, where $m_a$ and $\sigma_a$ are the mean and the standard deviation of
$\nu_a$). So every suboptimal
arm costs more than $\kappa\Delta_a$ per $\log T$, and $(\kappa+3/2+o(1))\Delta_a$ when its gap is large
compared to its standard deviation: in that regime, this is the exact coefficient of the term
$(\kappa-1)\Delta_a$ of Lattimore's bound. The number $e^{-\mathcal K_\infty(\kappa)}$ is the largest mass that
a law of kurtosis at most $\kappa$ can put on one point.

The method of \citet{AJK21} for a bounded moment relies on the convexity of the class, and the kurtosis
constraint is not convex in the law. We remove the non-convexity by optimising over the location and
the scale of the alternative, which leaves a convex problem with two moment constraints. Because the
class is scale free, near-optimal alternatives can escape along the affine group; the new step is to
show that this escape costs exactly $\mathcal K_\infty(\kappa)$, which gives the lower semicontinuity of
$\mathcal K$ and the lower deviation inequality that the regret analysis needs.

\paragraph{Our contributions} We summarise our main contributions below.
\begin{enumerate}[label=(\roman*),leftmargin=2.2em]
  \item \emph{The information function of the class.} A dual formula for $\mathcal K(P,\mu)$, valid for every law
    $P$, as an infimum over location and scale of a two-dimensional concave maximisation
    (Proposition~\ref{prop:dual}); the largest atom of a law of kurtosis at most $\kappa$ and the value
    $\mathcal K_\infty(\kappa)$ (Proposition~\ref{prop:kinfty}), which improves the large-gap branch of the
    previous upper bound on $\mathcal{K}_{\inf}$ \citep[Theorem~9]{L17}; lower semicontinuity, continuity in the
    level, positivity and the limit of large gaps (Proposition~\ref{prop:regularity},
    Corollary~\ref{cor:largegap}).
  \item \emph{Algorithm.} The anytime index policy \algo{}, KL-UCB with $\mathcal K$ and the threshold of
    \citet{AJK21}; see Section~\ref{sec:algorithm}.
  \item \emph{Regret upper bound and exact constant.} A finite-time bound on $\E[N_a(T)]$ with an
    existential constant and the first-order limit \eqref{eq:intro-main} (Theorem~\ref{thm:main},
    Corollary~\ref{cor:exact}).
  \item \emph{Concentration.} An anytime upper deviation inequality for $n\mathcal K(\hat P_n,m(P))$ with explicit
    constants, a lower deviation inequality, and a scale-free anytime confidence sequence for the mean
    (Theorem~\ref{thm:deviation}, Corollary~\ref{cor:confseq}).
  \item \emph{Numerical illustration.} Values of $\mathcal K_\infty(\kappa)$ and of $\mathcal K$ on discretised laws,
    and a check of the dual formula against a direct search (Section~\ref{sec:examples}).
\end{enumerate}

\paragraph{Outline} Section~\ref{sec:related} discusses related work. Section~\ref{sec:setting} formalises
the setting. Section~\ref{sec:lower} recalls the lower bound and studies $\mathcal K$.
Section~\ref{sec:algorithm} recalls KL-UCB and defines \algo{}. Section~\ref{sec:upper} states the regret
bound, and Section~\ref{sec:concentration} the deviation inequalities. Section~\ref{sec:overview}
explains how the proofs work, Section~\ref{sec:examples} gives numerical examples, and
Section~\ref{sec:conclusion} concludes and Section~\ref{sec:open} states open questions. The proofs are in the
appendices.

% ----------------------------------------------------------------
\section{Related work and context}\label{sec:related}

Table~\ref{tab:comparison} compares the first-order results closest to ours. We discuss each line of
work below.

\paragraph{Asymptotic optimality in stochastic bandits} The lower bound \eqref{eq:intro-lb} is due to Lai
and Robbins for one-parameter families and to Burnetas and Katehakis for general models; we use the
form of \citet[Theorem~16.2]{LS}, recalled as Theorem~\ref{cite:LS-16.2}. KL-UCB attains it for
one-parameter exponential families and finitely supported laws \citep{CGMMS13}, and the dual formula of
\citet[Theorem~8]{HT10} gives the same for laws on $[0,1]$. For Gaussian laws with
unknown mean and variance, $\mathcal{K}_{\inf}=\frac12\log(1+\Delta^2/\sigma^2)$ \citep[equations~(2)--(3)]{HT14}, attained by
Thompson sampling with a suitable prior \citep[Theorem~1]{HT14} and by an index policy
\citep[Theorem~5]{CHK18}. For the class $\{P:\E_P|X|^{1+\beta}\le B\}$,
\citet{AJK21} introduced KL$_{\inf}$-UCB, whose index is an inverse of $\mathcal{K}_{\inf}$ computed through a dual
formula, and proved it asymptotically optimal \citep[Corollary~2, recalled as Lemma~\ref{cite:AJK21-cor2}]{AJK21}.

\paragraph{Scale-free bandits} \citet{L17} introduced the bounded-kurtosis class to obtain
logarithmic regret without any scale parameter. His algorithm uses median-of-means estimates of the
mean and of the variance; its regret is at most $C\sum_a\Delta_a(\kappa-1+\sigma_a^2/\Delta_a^2)\log T$
asymptotically (Lemma~\ref{cite:L17-thm2}). His Theorem~9 bounds $\mathcal{K}_{\inf}$ from above by explicit
perturbations of the law, a Bernoulli ``outlier'' and a smooth density change, for $\kappa\ge7/2$
(Lemma~\ref{cite:L17-thm9}); with \eqref{eq:intro-lb} it gives a lower bound matching his upper bound up
to constants. He also showed that $\kappa$ must be known: no algorithm with logarithmic regret on some
non-trivial instance is consistent over the union of the classes \citep[Theorem~7]{L17}.

\paragraph{Non-convex information functions} In their concluding remarks, \citet[Section~5]{AJK21}
observe that a bound on a \emph{centred} moment is not a convex constraint on the law, so that their
method does not apply, and that handling this non-convexity would need new techniques.
For best-arm identification under a CVaR objective, \citet{AKJ21} met the same obstacle for the lower information projection: they used the variational
formula of CVaR to write it as an outer minimisation over one scalar, which ranges over a compact
interval, of a convex problem treated by Lagrangian duality, and for the concentration they fixed that
scalar at its true value before mixing supermartingales over the dual variables \citep[see also][Chapter~6, Theorem~6.5 and Appendix~6.D.1]{A22}. In the discussion of her thesis,
\citet[Section~6.10]{A22} rewrites $\mathcal{K}_{\inf}$ for the class $\{P:\E_P|X-m(P)|^{1+\beta}\le B\}$ as an
outer minimisation over the location $y\ge x$ of a convex problem, states that a concentration result
similar to hers can be proven for it, and leaves the continuity of the new projection functions and the
extreme distributions of the class to separate research. The thesis does not treat the kurtosis or a
scale parameter. \citet[Section~4.4.2 and Appendix~D.5]{BSH23} carry out this reduction for the centred class
$\{P:\E_P h(|X-m(P)|)<B\}$, with $h$ convex and $B$ known, under a known lower bound on the means, which they
impose to use compactness arguments: they write $\mathcal{K}_{\inf}$ as an infimum over the location $\mu\ge x$ of a convex
problem at fixed location, with a two-dimensional dual, observe that the infimum is attained at $\mu=x$, and prove
that MED is asymptotically optimal on this class \citep[Lemma~9 and Corollary~6]{BSH23}. The level $B$ is fixed
and only the location is profiled, so a translation reduces their problem to the uncentred one; no scale
parameter appears.

\paragraph{Generic optimality of KL$_{\inf}$-UCB} \citet{PA26} prove that KL$_{\inf}$-UCB is
asymptotically optimal on every class whose laws satisfy two assumptions: an anytime upper deviation
inequality for $n\mathcal{K}_{\inf}(\hat P_n,m(P))$ and an exponential lower deviation inequality at every fixed
level above the mean (Lemma~\ref{cite:PA26-thm3.1}). Their examples are bounded supports and the
uncentred moment class; the kurtosis is not mentioned.

\paragraph{Closest related work} Five works are close to ours. \citet{L17} studies the same class and the
same criterion, but its upper and lower bounds hold up to universal constants, so it does not give the
constant \eqref{eq:intro-main}. \citet{AJK21} have the same algorithmic template and threshold, but its
analysis uses the convexity of the class $\{\E|X|^{1+\beta}\le B\}$ and its compact dual region at a fixed
mean; neither holds here. \citet{PA26} reduce optimality to two deviation inequalities, but checking them
is the content of the problem: the upper one requires a convexification, and the lower one, in the form
that \citet{PA26} ask (a rate $e^{-ncd^2}$ for all small $d$), is not proved here; we prove a weaker form
and give a self-contained regret analysis. \citet{AKJ21} and \citet{A22} contain the device of an outer scalar
minimisation over a convex problem, with the outer variable fixed at its true value for the upper
deviation; we use it with two outer variables, and their outer variable stays in a compact set, while
ours does not. \citet{BSH23} prove first-order optimality on a non-convex centred class by an outer
minimisation over the location, but their constraint has a fixed level and their means a known lower bound,
so the outer variable ranges over a half-line and the infimum is attained at its end; the kurtosis constraint
is a ratio of moments, invariant under changes of scale, and gives no bound on the location or the scale of
near-optimal alternatives.

\paragraph{Gap with the literature} On the bounded-kurtosis class, no algorithm was known to attain the
lower bound \eqref{eq:intro-lb}, and no description of $\mathcal{K}_{\inf}$ beyond upper bounds up to constants.
The methods that attain \eqref{eq:intro-lb} on non-parametric classes use the convexity of the
class \citep{AJK21}, assume the two deviation inequalities \citep{PA26}, or, for a centred moment class, profile
the location under a known lower bound on the means \citep{BSH23}; the convexification by an outer
minimisation was also used for CVaR, where the outer variable is confined to a compact interval
\citep{AKJ21,A22}. Outside bandits, treating the mean and the variance as nuisance parameters of moment
equations, in order to make inference on the skewness or the kurtosis, is classical in empirical likelihood
\citep[see also][Sections~1.1 and 3.5]{O01}. What is new here, as far as we found, is the convexification of the
inequality $\mathrm{kurt}\le\kappa$ by a quartic with a free location and a free scale inside an information function
$\mathcal{K}_{\inf}$, the analysis of the mass that escapes along the affine group, with no bound on the means or
the scales of the alternatives and with the exact boundary value $\mathcal K_\infty(\kappa)$ obtained from the
largest atom, and the resulting first-order optimality on the bounded-kurtosis class.

\begin{table}[t]
\centering
\small
\caption{First-order results for KL-type algorithms beyond one-parameter families. ``Exact'' means
that the upper bound on $\limsup_TR_T/\log T$ equals the constant of the lower bound
\eqref{eq:intro-lb} for the same class.}\label{tab:comparison}
\begin{tabular}{@{}p{0.24\textwidth}p{0.33\textwidth}p{0.1\textwidth}p{0.25\textwidth}@{}}
\toprule
Class of laws & First-order upper bound & Exact & Reference \\
\midrule
Gaussian, unknown mean and variance & $\sum_a\Delta_a/\bigl(\frac12\log(1+\Delta_a^2/\sigma_a^2)\bigr)$ & yes & \citet[Theorem~1]{HT14}, \citet[Theorem~5]{CHK18} \\
$\E|X|^{1+\beta}\le B$ (convex) & $\sum_a\Delta_a/\Kinf(\nu_a,\mstar)$ & yes & \citet[Corollary~2]{AJK21} \\
Any class satisfying two deviation inequalities & $\sum_a\Delta_a/\Kinf(\nu_a,\mstar)$ & yes, if they hold & \citet[Theorem~3.1]{PA26} \\
$\E h(|X-m|)<B$ (centred), means bounded below by a known number & $\sum_a\Delta_a/\Kinf(\nu_a,\mstar)$, for MED & yes & \citet[Corollary~6]{BSH23} \\
Kurtosis at most $\kappa$ & $C\sum_a\Delta_a(\kappa-1+\sigma_a^2/\Delta_a^2)$, $C$ universal & no & \citet[Theorem~2]{L17} \\
Kurtosis at most $\kappa$ & $\sum_a\Delta_a/\Kinf(\nu_a,\mstar)$, with $\Kc(\nu_a,\mstar)\le\Kinfty(\kappa)$ & yes & Theorem~\ref{thm:main}, Corollary~\ref{cor:exact} \\
\bottomrule
\end{tabular}
\end{table}

% ----------------------------------------------------------------
\section{Problem formulation}\label{sec:setting}

In this section we introduce the setting, first the classical one, then the bounded-kurtosis class and
the normalised objects that the proofs use.

\subsection{Classical setting}

\dt{w:law}A \emph{law} is a Borel probability measure on $\R$. The set of \term{w:law}{laws} is $\Pspace$, with the
topology of weak convergence (convergence of $\int f\,dP$ for every bounded continuous $f$), which is
metrisable.

\begin{definition}[Laws, moments, divergence]\label{def:laws}
Let $P,Q$ be \term{w:law}{laws}. If $\int x^2\,dP<\infty$, the mean and the variance of $P$ are
\begin{equation*}
  \dt{s:mean}\mean(P):=\int x\,dP(x),\qquad\dt{s:var}\sd^2(P):=\int(x-\mean(P))^2\,dP(x),
\end{equation*}
and $\sd(P):=\sqrt{\sd^2(P)}$ is its standard deviation. If moreover $\int x^4\,dP<\infty$ and $\sd^2(P)>0$, the
kurtosis of $P$ is
\begin{equation*}
  \dt{s:kurt}\kurt(P):=\frac{\int(x-\mean(P))^4\,dP(x)}{\sd^4(P)}.
\end{equation*}
The Kullback--Leibler divergence is $\dt{s:KL}\KL(P,Q):=\int\log\frac{dP}{dQ}\,dP$ if $P\ll Q$, and
$\KL(P,Q):=+\infty$ otherwise.
\end{definition}

\begin{definition}[Bandit model, counts, regret, consistency]\label{def:bandit}
\uses{def:laws}
A finite set $\alpha$ of arms is given. \dt{s:Karms}The number of arms is $\Karms:=|\alpha|\ge2$.
\dt{w:instance}An \emph{instance} is $\nu=(\nu_a)_{a\in\alpha}$, each $\nu_a$ a \term{w:law}{law} with a finite mean.
\dt{s:ma}The mean of arm $a$ is $\marm_a:=\mean(\nu_a)$. \dt{s:sa}Its standard deviation is $\sda_a:=\sd(\nu_a)$ when it is
finite. \dt{s:mstar}The best mean is $\mstar:=\max_a\marm_a$. \dt{s:gap}The gap of arm $a$ is $\gap_a:=\mstar-\marm_a$. \dt{w:subopt}Arm $a$ is \emph{suboptimal} if $\gap_a>0$.
\dt{w:model}A \emph{model} is a set $\mathcal D$ of \term{w:law}{laws}, and the \term{w:instance}{instances} of the
\term{w:model}{model} are $\mathcal D^\alpha$. The rewards follow the stack model: \dt{s:X}independent random
variables $(\Xr_{a,j})_{a\in\alpha,\,j\ge1}$ with $\Xr_{a,j}\sim\nu_a$, the $j$-th pull of arm $a$ returning $\Xr_{a,j}$.
\dt{s:At}An algorithm chooses the arm $\Arm_t$ of round $t\in\{1,2,\dots\}$ as a measurable function of the
past arms, the past rewards and an independent source of randomness. \dt{w:anytime}It is \emph{anytime}: it
never sees a horizon. The number of pulls of $a$ after round $t$ is
\begin{equation*}
  \dt{s:N}\Np_a(t):=\sum_{u=1}^{t}\ind(\Arm_u=a).
\end{equation*}
The regret after $T$ rounds is
\begin{equation*}
  \dt{s:R}\Reg_T:=\sum_{a\in\alpha}\gap_a\,\E[\Np_a(T)].
\end{equation*}
For $s\ge1$, the empirical law of the first $s$ rewards of arm $a$ is
$\dt{s:Pe}\Pe_{a,s}:=\frac1s\sum_{j\le s}\delta_{\Xr_{a,j}}$, and $\Pe_a(t):=\Pe_{a,\Np_a(t)}$ when $\Np_a(t)\ge1$.
\dt{w:consistent}An algorithm is \emph{consistent over $\mathcal D^\alpha$} if $\Reg_T(\nu)=o(T^p)$ for every $p>0$ and
every $\nu\in\mathcal D^\alpha$.
\end{definition}

\begin{definition}[Information function]\label{def:kinfD}
\uses{def:laws}
For a \term{w:model}{model} $\mathcal D$, a \term{w:law}{law} $P$ and $\mu\in\R$,
\begin{equation*}
  \dt{s:KinfD}\Kinf(P,\mu;\mathcal D):=\inf\{\KL(P,Q):Q\in\mathcal D,\ \mean(Q)>\mu\}\in[0,\infty],\qquad\inf\emptyset:=+\infty,
\end{equation*}
with a \emph{strict} inequality on the mean.
\end{definition}

\paragraph{Notational convention} Rounds are numbered from $1$; $\Np_a(T)$ counts the pulls in rounds
$1,\dots,T$. Some sources define $\Kinf$ with $\mean(Q)\ge\mu$ \citep{AJK21,PA26}; Proposition~\ref{prop:dual}(i)
shows that the two versions coincide on our class. We write $\log$ for the natural logarithm, $x\vee y$
for $\max(x,y)$, and $\tauc_{c,s}P$ for an image law (Definition~\ref{def:normalised}).

\subsection{The bounded-kurtosis class}

\begin{definition}[Bounded kurtosis]\label{def:class}
\uses{def:laws}
Let $\kappa>1$. The class of \term{w:law}{laws} with kurtosis at most $\kappa$ is
\begin{equation*}
  \dt{s:Hk}\Hk:=\Bigl\{P\in\Pspace:\ \int x^4\,dP<\infty,\ \sd^2(P)>0,\ \kurt(P)\le\kappa\Bigr\}.
\end{equation*}
Its strict version is
\begin{equation*}
  \dt{s:Hlt}\Hlt:=\bigl\{P\in\Hk:\ \kurt(P)<\kappa\bigr\}.
\end{equation*}
\end{definition}

Every \term{w:law}{law} with a finite fourth moment has kurtosis at least $1$, with equality exactly for the
symmetric two-point laws, so $\kappa>1$ excludes only degenerate classes. The assumption of
\citet[Section~3]{L17} is a kurtosis at most $\kappa$, and his class in the lower bound is $\Hlt$; our results hold
for both (Corollary~\ref{cor:exact}(iv)). Both classes are invariant under the maps $x\mapsto c+sx$, $c\in\R$,
$s>0$.

\begin{definition}[Information function of the class]\label{def:kinf}
\uses{def:class, def:kinfD}
For a \term{w:law}{law} $P$ and $\mu\in\R$, the information function of the class is
\begin{equation*}
  \dt{s:K}\Kc(P,\mu):=\Kinf(P,\mu;\Hk).
\end{equation*}
Its non-strict version is
\begin{equation*}
  \dt{s:Kge}\Kge(P,\mu):=\inf\{\KL(P,Q):Q\in\Hk,\ \mean(Q)\ge\mu\}.
\end{equation*}
\end{definition}

\subsection{Normalised objects}

The kurtosis constraint is not convex in the \term{w:law}{law}, but it becomes convex once the location and
the scale are fixed. The following objects encode this.

\begin{definition}[Normalised class]\label{def:normalised}
\uses{def:laws, def:kinfD}
Let $\kappa>1$. The quartic of the constraint is
\begin{equation*}
  \dt{s:phi}\phik(y):=(y^2-\kappa)^2-\kappa(\kappa-1)=y^4-2\kappa y^2+\kappa,\qquad y\in\R.
\end{equation*}
The normalised class is
\begin{equation*}
  \dt{s:G}\Gk:=\Bigl\{Q\in\Pspace:\ \int y^4\,dQ<\infty,\ \int y\,dQ(y)=0,\ \int\phik\,dQ\le0\Bigr\}.
\end{equation*}
Its information function is, for a \term{w:law}{law} $R$,
\begin{equation*}
  \dt{s:KG}\KG(R):=\inf_{Q\in\Gk}\KL(R,Q).
\end{equation*}
For $c\in\R$ and $s>0$, the normalising map is $\dt{s:tau}\tauc_{c,s}(x):=(x-c)/s$, and $\tauc_{c,s}P$ is the
image of a \term{w:law}{law} $P$ under $\tauc_{c,s}$.
\end{definition}

The function $\phik$ has its minimum $-\kappa(\kappa-1)$ at $y=\pm\sqrt\kappa$. If $Y$ has mean $0$ and variance
$1$, then $\E\phik(Y)=\E Y^4-\kappa$, so $\E\phik(Y)\le0$ says exactly that the kurtosis is at most $\kappa$; the
point of $\Gk$ is that it imposes this without fixing the variance, through a constraint that is affine
in the \term{w:law}{law}.

\begin{definition}[Dual set]\label{def:dualset}
\uses{def:normalised}
For $\lambda=(\lambda_1,\lambda_2)\in\R^2$ and $y\in\R$, the dual integrand is
\begin{equation*}
  \dt{s:h}\hl_\lambda(y):=1+\lambda_1y+\lambda_2\phik(y).
\end{equation*}
The dual set is
\begin{equation*}
  \dt{s:Lambda}\Lam:=\bigl\{\lambda\in\R\times[0,\infty):\ \hl_\lambda(y)\ge0\ \text{for all }y\in\R\bigr\}.
\end{equation*}
Its half-width is, for $\lambda_2\ge0$,
\begin{equation*}
  \dt{s:L}\Lw(\lambda_2):=\inf_{y>0}\frac{1+\lambda_2\phik(y)}{y}.
\end{equation*}
\end{definition}

\begin{definition}[Largest atom and value at infinity]\label{def:kinfty}
For $\kappa>1$, the largest atom is
\begin{equation*}
  \dt{s:q}\qk:=\frac12\Bigl(1+\sqrt{\frac{\kappa-1}{\kappa+3}}\Bigr)\in\Bigl(\frac12,1\Bigr).
\end{equation*}
The value at infinity is
\begin{equation*}
  \dt{s:Kinfty}\Kinfty(\kappa):=\log\frac1{\qk}.
\end{equation*}
\end{definition}

The number $\qk$ is the larger root of $q(1-q)=1/(\kappa+3)$; a two-point \term{w:law}{law} with masses $\qk$ and
$1-\qk$ has kurtosis $1/(\qk(1-\qk))-3=\kappa$.
% ----------------------------------------------------------------
\section{Lower bound and the information function of the class}\label{sec:lower}

In this section we first recall the classical lower bound, then describe the information function
$\Kc$ that it involves on the bounded-kurtosis class.

\subsection{Lower bound in the classical setting}

The following theorem is classical (we refer to \citet{LS} for a proof; its original source is
Burnetas and Katehakis). We state it in our notation: \citet{LS} write $n$ for the horizon, $k$ for the
number of arms and $d_{\inf}(P,\mu^*,\mathcal M)$ for $\Kinf(P,\mu^*;\mathcal M)$.

\begin{theorem}[{From \citet[Theorem~16.2]{LS}}]\label{cite:LS-16.2}
\uses{def:bandit, def:kinfD}
Let $\mathcal D$ be a set of \term{w:law}{laws} with finite means and let an algorithm be \term{w:consistent}{consistent} over $\mathcal D^\alpha$. Then for
every $\nu\in\mathcal D^\alpha$,
\begin{equation*}
  \liminf_{T\to\infty}\frac{\Reg_T}{\log T}\ \ge\ \sum_{a:\gap_a>0}\frac{\gap_a}{\Kinf(\nu_a,\mstar;\mathcal D)}.
\end{equation*}
\end{theorem}

On the class $\Hk$, the constant of this bound is
\begin{equation}\label{eq:cstar}
  \dt{s:Cstar}\Cst(\nu):=\sum_{a:\gap_a>0}\frac{\gap_a}{\Kc(\nu_a,\mstar)},\qquad \nu\in\Hk^\alpha.
\end{equation}
The rest of this section shows that $\Cst(\nu)$ is finite and positive, and describes $\Kc$.

\subsection{A dual formula}

\begin{proposition}[Normalisation and dual formula]\label{prop:dual}
\uses{def:kinf, def:normalised, def:dualset}
Let $\kappa>1$. Then:
\begin{enumerate}[label=(\roman*)]
  \item for every \term{w:law}{law} $P$ and every $\mu\in\R$,
  \begin{equation*}
    \Kc(P,\mu)=\Kge(P,\mu)=\inf_{c>\mu,\ s>0}\KG(\tauc_{c,s}P)=\inf_{c\ge\mu,\ s>0}\KG(\tauc_{c,s}P),
  \end{equation*}
  and $\mu\mapsto\Kc(P,\mu)$ is non-decreasing and right-continuous;
  \item $\Gk$ is a convex and weakly compact subset of $\Hk$; $\KG$ is weakly lower semicontinuous on
  $\Pspace$, and when $\KG(R)<\infty$ the infimum defining it is attained;
  \item $\Lam$ is a compact convex subset of $[-\kappa^{-1/2},\kappa^{-1/2}]\times[0,\frac1{\kappa(\kappa-1)}]$ with
  non-empty interior, and
  \begin{equation*}
    \Lam=\Bigl\{(\lambda_1,\lambda_2):\ 0\le\lambda_2\le\tfrac1{\kappa(\kappa-1)},\ |\lambda_1|\le\Lw(\lambda_2)\Bigr\};
  \end{equation*}
  \item for every \term{w:law}{law} $R$ (the \term{w:law}{law} of a random variable $Y$): if $\E_R\log(1+|Y|)=\infty$, then
  $\KG(R)=\infty=\sup_{\lambda\in\mathrm{int}\,\Lam}\E_R\log\hl_\lambda(Y)$; otherwise $\E_R\log\hl_\lambda(Y)\in[-\infty,\infty)$
  is well defined for every $\lambda\in\Lam$, and
  \begin{equation*}
    \KG(R)=\max_{\lambda\in\Lam}\E_R\log\hl_\lambda(Y)=\sup_{\lambda\in\mathrm{int}\,\Lam}\E_R\log\hl_\lambda(Y)<\infty;
  \end{equation*}
  \item consequently, for every \term{w:law}{law} $P$ and every $\mu\in\R$ (with $X\sim P$),
  \begin{equation}\label{eq:dual}
    \Kc(P,\mu)=\inf_{c>\mu,\ s>0}\ \sup_{\lambda\in\mathrm{int}\,\Lam}\ \E_P\log\Bigl(1+\lambda_1\frac{X-c}{s}+\lambda_2\,\phik\Bigl(\frac{X-c}{s}\Bigr)\Bigr);
  \end{equation}
  \item for a finitely supported \term{w:law}{law} $R$, the infimum defining $\KG(R)$ is attained by a \term{w:law}{law} with at most
  two atoms outside the support of $R$.
\end{enumerate}
\end{proposition}

(See the proof in Appendix~\ref{sec:proof-dual}.)

\paragraph{Why location and scale} By Lemma~\ref{lem:G} in Appendix~\ref{sec:proof-dual}, a \term{w:law}{law} $Q$ belongs to
$\Hk$ with mean $c$ if and only if $\tauc_{c,s}Q\in\Gk$ for some $s>0$. The scale $s$ is a free variable,
not the standard deviation of $Q$: the constraint $\int\phik\,d(\tauc_{c,s}Q)\le0$ is a quadratic inequality
in $s^2$ whose discriminant is non-negative exactly when $\kurt(Q)\le\kappa$. Fixing $(c,s)$ leaves the convex
set $\Gk$, defined by one linear equality and one linear inequality, and Lagrangian duality gives the
inner maximisation over $\Lam$. The outer problem over $(c,s)$ is not convex, and its domain is not
compact.

\paragraph{Comparison with the convex case} For the class $\{\E|X|^{1+\beta}\le B\}$ the analogous formula is a
single maximisation over a compact dual region that depends on the level \citep[Lemma~7]{AJK21}. Here the
dual set $\Lam$ is fixed, it depends on $\kappa$ only, and the level enters through the location $c$; this is
what makes the mixture argument of Theorem~\ref{thm:deviation}(i) work with explicit constants.

\paragraph{Structure of the optimal alternatives} Item (vi) mirrors the structure known for the convex
moment class \citep[Lemma~7]{AJK21}: an optimal alternative adds to the support of $R$ at most two points,
the zeros of the optimal dual function $\hl_{\lambda^\star}$. In the numerical examples of
Section~\ref{sec:examples} the optimal location is $c=\mu$ and the optimal alternative adds one point far to
the right of the data.

\subsection{The largest atom and the value at infinity}

\begin{proposition}[Value at infinity]\label{prop:kinfty}
\uses{def:kinf, def:kinfty, def:class}
Let $\kappa>1$. Then:
\begin{enumerate}[label=(\roman*)]
  \item for every $Q\in\Hk$ and every $y\in\R$, $Q(\{y\})\le\qk$, with equality for the two-point \term{w:law}{law}
  $\qk\delta_y+(1-\qk)\delta_{y'}$, $y'\ne y$, whose kurtosis is $\kappa$;
  \item for every \term{w:law}{law} $P$ with $\int x^4\,dP<\infty$ and every $\mu\in\R$, $\Kc(P,\mu)\le\Kinfty(\kappa)$;
  \item for every $y\in\R$ and every $\mu\in\R$, $\Kc(\delta_y,\mu)=\Kinfty(\kappa)$;
  \item the value at infinity satisfies
  \begin{equation*}
    \kappa<\frac1{\Kinfty(\kappa)}<\kappa+2\qquad\text{and}\qquad\frac1{\Kinfty(\kappa)}=\kappa+\frac32+O\Bigl(\frac1\kappa\Bigr)\quad(\kappa\to\infty).
  \end{equation*}
\end{enumerate}
\end{proposition}

(See the proof in Appendix~\ref{sec:proof-kinfty}.)

\paragraph{Comparison with the bound of Lattimore} The previous upper bound on $\Kinf$ for this class is
that of \citet[Theorem~9]{L17}, recalled as Lemma~\ref{cite:L17-thm9}, for $\kappa\ge7/2$. For large gaps, $\Delta/\sd(P)\ge1/\kappa$ and
$\Delta/\sd(P)$ beyond the threshold of its first case, it reads
\begin{equation*}
  \Kinf(P,\mean(P)+\Delta;\Hlt)\le\log\frac{\kappa}{\kappa-1}.
\end{equation*}
By (ii) and Corollary~\ref{cor:exact}(iv), the same quantity is at most $\Kinfty(\kappa)$, for every $\kappa>1$, every \term{w:law}{law}
with a finite fourth moment and every $\Delta$. This improves the large-gap branch of \citet[Theorem~9]{L17}, since
\begin{equation*}
  \Kinfty(\kappa)<\frac1\kappa<\log\frac{\kappa}{\kappa-1},
\end{equation*}
and it cannot be improved further: by Corollary~\ref{cor:largegap}, $\Kc(P,\mean(P)+\Delta)\to\Kinfty(\kappa)$ as
$\Delta/\sd(P)\to\infty$. In terms of regret, the lower bound per unit gap and per $\log T$ for large gaps goes from
$1/\log\frac\kappa{\kappa-1}\in(\kappa-1,\kappa)$ to $1/\Kinfty(\kappa)\in(\kappa,\kappa+2)$, which is exact. The small-gap branch of
\citet[Theorem~9]{L17}, of order $\Delta^2/\sd^2(P)$, is not improved here.

\paragraph{Interpretation} The proof of (ii) mixes $P$ with a far atom: $(1-p)P+p\delta_x$ with $x\to\infty$ has
a kurtosis that tends to $1/(p(1-p))-3$, below $\kappa$ as soon as $p>1-\qk$, and a mean that tends to
infinity. This is Lattimore's ``outlier'' idea with the sharp mass. Conversely, (i) says that an
alternative cannot put more than $\qk$ of its mass on a point, so it cannot be close to a Dirac mass; this
is where the constant $\Kinfty(\kappa)$ comes from in the lower semicontinuity below. The forced
exploration of Section~\ref{sec:algorithm} has the same origin: an arm keeps the largest possible index until it has
been pulled about $\log t/\Kinfty(\kappa)>\kappa\log t$ times.

\subsection{Regularity}

\begin{proposition}[Lower semicontinuity, continuity, positivity]\label{prop:regularity}
\uses{def:kinf, def:kinfty, def:class}
Let $\kappa>1$, let $P$ be a \term{w:law}{law} with $\int x^4\,dP<\infty$ that is not a Dirac mass, and let $\mu\in\R$. Then:
\begin{enumerate}[label=(\roman*)]
  \item for every sequence $(P_n,\mu_n)$ with $P_n\to P$ weakly and $\mu_n\to\mu$,
  $\liminf_n\Kc(P_n,\mu_n)\ge\Kc(P,\mu)$;
  \item $\mu'\mapsto\Kc(P,\mu')$ is continuous on $\R$ and $\lim_{\mu'\to\infty}\Kc(P,\mu')=\Kinfty(\kappa)$;
  \item if $\mu>\mean(P)$, then $\Kc(P,\mu)>0$.
\end{enumerate}
\end{proposition}

(See the proof in Appendix~\ref{sec:proof-regularity}.)

\paragraph{Failure at Dirac masses} The hypothesis that $P$ is not a Dirac mass cannot be removed from (i).
For $\kappa\ge3$, the Gaussian \term{w:law}{laws} $P_n:=\mathcal N(0,n^{-2})$ belong to $\Hk$ and have mean $0$, so
$\Kc(P_n,0)=\Kge(P_n,0)=0$ by Proposition~\ref{prop:dual}(i), while $P_n\to\delta_0$ and $\Kc(\delta_0,0)=\Kinfty(\kappa)$ by
Proposition~\ref{prop:kinfty}(iii). Empirical \term{w:law}{laws} can be Dirac masses, which is why the regret analysis
uses (i) only at the true \term{w:law}{law} of an arm.

\begin{corollary}[Location, scale and large gaps]\label{cor:largegap}
\uses{def:kinf, def:kinfty, def:class}
Let $\kappa>1$ and $P\in\Hk$ with mean $m$, standard deviation $\sigma$ and standardised \term{w:law}{law}
$P_0:=\tauc_{m,\sigma}P$. Then $\Kc(P,m+\Delta)=\Kc(P_0,\Delta/\sigma)$ for every $\Delta\in\R$, and
$\Delta\mapsto\Kc(P,m+\Delta)$ is continuous and non-decreasing on $\R$, equal to $0$ on $(-\infty,0]$,
positive on $(0,\infty)$, with
\begin{equation*}
  \lim_{\Delta\to\infty}\Kc(P,m+\Delta)=\Kinfty(\kappa).
\end{equation*}
\end{corollary}

(See the proof in Appendix~\ref{sec:proof-regularity}.)

% ----------------------------------------------------------------
\section{The algorithm}\label{sec:algorithm}

In this section we recall the KL-UCB template for a non-parametric class, then define \algo{}.

\subsection{KL-UCB in the classical setting}

For a model $\mathcal D$, KL-UCB pulls at each round an arm whose index, an upper confidence bound on its mean,
is maximal; the index of an arm with empirical \term{w:law}{law} $\hat P$ after $n$ pulls is
$\sup\{\mu:n\Kinf(\hat P,\mu;\mathcal D)\le g\}$ for a threshold $g$ of order $\log t$. For the class
$\{\E|X|^{1+\beta}\le B\}$, KL$_{\inf}$-UCB of \citet{AJK21} uses the threshold $\log t+2\log\log t+2\log(1+n)+1$
and is asymptotically optimal (Lemma~\ref{cite:AJK21-cor2}). Its index is computed through the dual formula
of $\Kinf$ for that class \citep[Lemma~7]{AJK21}.

\subsection{\texorpdfstring{The algorithm \texttt{KL-UCB-Kurt}}{The algorithm KL-UCB-Kurt}}

We use the same template with the information function $\Kc$ of the class $\Hk$. \dt{s:g}The threshold is
\begin{equation}\label{eq:threshold}
  \thr(t,n):=\log(t\vee3)+2\log\log(t\vee3)+2\log(1+n)+1,\qquad t,n\ge1.
\end{equation}
\dt{s:U}At the end of round $t\ge\Karms$, the index of arm $a$ is
\begin{equation}\label{eq:index}
  \Ix_a(t):=\sup\bigl\{\mu\in\R:\ \Np_a(t)\,\Kc\bigl(\Pe_a(t),\mu\bigr)\le\thr\bigl(t,\Np_a(t)\bigr)\bigr\}\in[-\infty,+\infty],
  \qquad\sup\emptyset:=-\infty,
\end{equation}
and the next arm is $\Arm_{t+1}\in\arg\max_a\Ix_a(t)$. Intuitively, $\Ix_a(t)$ is the largest mean of a \term{w:law}{law} of
kurtosis at most $\kappa$ that is still compatible with the rewards of arm $a$ at the confidence level set
by the threshold.

\begin{algorithm}[t]
\caption{\texttt{KL-UCB-Kurt}}\label{alg:kucb}
\begin{algorithmic}[1]
\Require the number $\kappa>1$ (no horizon, no scale)
\For{$t=1,\dots,\Karms$}
  \State pull each arm once, in an arbitrary order
\EndFor
\For{$t=\Karms,\Karms+1,\dots$}
  \State compute $\Ix_a(t)$ by \eqref{eq:index} for every arm $a$
  \State pull $\Arm_{t+1}\in\arg\max_a\Ix_a(t)$ (ties broken by any fixed rule, for example uniformly at random)
\EndFor
\end{algorithmic}
\end{algorithm}

\begin{definition}[\texttt{KL-UCB-Kurt}]\label{def:algorithm}
\uses{def:bandit, def:kinf}
For $\kappa>1$, \algo{} is the anytime algorithm of Algorithm~\ref{alg:kucb}: in rounds $1,\dots,\Karms$ it pulls
every arm once; for $t\ge\Karms$ it pulls at round $t+1$ an arm maximising the index \eqref{eq:index}, with the
threshold \eqref{eq:threshold}. The order of the first $\Karms$ pulls and the tie rule are arbitrary
(measurable); the results hold for all of them.
\end{definition}

\begin{lemma}[The index]\label{lem:index}
\uses{def:algorithm, def:kinf, def:kinfty}
Let $\kappa>1$, $n\ge1$ and $\mu\in\R$. Then:
\begin{enumerate}[label=(\roman*)]
  \item the map $(x_1,\dots,x_n)\mapsto\Kc\bigl(\frac1n\sum_{i\le n}\delta_{x_i},\mu\bigr)$ is upper semicontinuous on
  $\R^n$, hence Borel; the index $\Ix_a(t)$ is a Borel function of the rewards observed up to round $t$,
  so \algo{} is a well-defined algorithm;
  \item if $\Np_a(t)\Kinfty(\kappa)\le\thr(t,\Np_a(t))$, then $\Ix_a(t)=+\infty$;
  \item if the $\Np_a(t)$ rewards of arm $a$ are all equal, then $\Ix_a(t)=+\infty$ when
  $\Np_a(t)\Kinfty(\kappa)\le\thr(t,\Np_a(t))$, and $\Ix_a(t)=-\infty$ otherwise.
\end{enumerate}
\end{lemma}

(See the proof in Appendix~\ref{sec:proof-index}.)

\paragraph{Forced exploration} By (ii), an arm has the largest possible index, $+\infty$, as long as
$\Np_a(t)\Kinfty(\kappa)\le\thr(t,\Np_a(t))$, that is, as long as its count is below about
\begin{equation*}
  \frac{\log t}{\Kinfty(\kappa)}>\kappa\log t .
\end{equation*}
This exploration is intrinsic to the class: the lower bound charges every \term{w:subopt}{suboptimal} arm at least
$\gap_a/\Kinfty(\kappa)$ per $\log T$ (Corollary~\ref{cor:exact}(iii)).

\paragraph{Arms with a Dirac empirical law} By (iii), when all the observed rewards of an arm are equal (a
discrete arm, for example a Bernoulli \term{w:law}{law} of parameter $\theta\in[1-\qk,\qk]$, can do this with probability
at most $\qk^{n-1}$ after $n$ pulls), its index is $+\infty$ while $n\Kinfty(\kappa)\le\thr(t,n)$ and
$-\infty$ afterwards. The algorithm then leaves this arm aside until $\thr(t,n)$, which grows like $\log t$,
reaches $n\Kinfty(\kappa)$ again; the definition $\sup\emptyset:=-\infty$ is what makes this happen, and the
upper deviation inequality (Theorem~\ref{thm:deviation}(i)) covers this event for an optimal arm.

\paragraph{Computation} By \eqref{eq:dual}, evaluating $\Kc(\hat P,\mu)$ for an empirical \term{w:law}{law} with $n$ atoms means
minimising over $(c,s)$ a function whose value is a two-dimensional concave maximisation over the fixed set
$\Lam$, each evaluation costing $O(n)$; the index is then found by bisection in $\mu$, since $\Kc$ is
non-decreasing in $\mu$. The inner problem is well behaved, but the outer problem is not convex, its domain is
not compact, and its infimum may be approached at infinity. We give no procedure with a certified accuracy:
Theorem~\ref{thm:main} assumes that the index is computed exactly, and an efficient algorithm with the same
guarantee is an open question (Section~\ref{sec:open}).

% ----------------------------------------------------------------
\section{Regret upper bound}\label{sec:upper}

We obtain the following guarantee on the regret of \algo{}. For a \term{w:subopt}{suboptimal} arm $a$ and $T\ge3$ we use
the information distance of arm $a$ to the best mean,
\begin{equation}\label{eq:Ka}
  \dt{s:Ka}\Ka:=\Kc(\nu_a,\mstar),
\end{equation}
the deterministic part of the threshold at the horizon,
\begin{equation*}
  \dt{s:lT}\LT:=\log T+2\log\log T+1,
\end{equation*}
and, for $\varepsilon\in(0,\Ka)$, the largest sample size at which the threshold still exceeds $n(\Ka-\varepsilon)$,
\begin{equation*}
  \dt{s:nT}\nT(\varepsilon):=\max\bigl\{n\ge0:\ n(\Ka-\varepsilon)\le\thr(T,n)\bigr\}.
\end{equation*}
By Propositions~\ref{prop:kinfty}(ii) and \ref{prop:regularity}(iii), $0<\Ka\le\Kinfty(\kappa)<1$ for every
$\nu\in\Hk^\alpha$ and every \term{w:subopt}{suboptimal} $a$.

\begin{theorem}[Regret of \texttt{KL-UCB-Kurt}]\label{thm:main}
\uses{def:bandit, def:class, def:kinf, def:kinfty, def:algorithm}
For every $\kappa>1$, every \term{w:instance}{instance} $\nu\in\Hk^\alpha$, every \term{w:subopt}{suboptimal} arm $a$ and every
$\varepsilon\in(0,\Ka)$, there is a constant $\Ca(\varepsilon)\in[0,\infty)$, depending only on $\nu_a$, $\mstar$,
$\varepsilon$ and $\kappa$, such that for every $T\ge3$ the algorithm \algo{} satisfies:
\begin{enumerate}[label=(\roman*)]
  \item with $\dt{s:keps}\kep:=\Ka-\varepsilon\in(0,1)$,
  \begin{equation}\label{eq:main-finite}
    \E[\Np_a(T)]\ \le\ \nT(\varepsilon)+\frac52+\Ca(\varepsilon)
    \ \le\ \frac{\LT+2\log\Bigl(1+\frac{2(\LT+2\log(4/\kep))}{\kep}\Bigr)}{\kep}+\frac52+\Ca(\varepsilon);
  \end{equation}
  \item at first order,
  \begin{equation}\label{eq:main-first}
    \limsup_{T\to\infty}\frac{\E[\Np_a(T)]}{\log T}\le\frac1{\Ka},\qquad\text{hence}\qquad
    \limsup_{T\to\infty}\frac{\Reg_T}{\log T}\le\Cst(\nu).
  \end{equation}
\end{enumerate}
\end{theorem}

(See the proof in Appendix~\ref{sec:proof-main}.)

\begin{corollary}[Exact first-order constant]\label{cor:exact}
\uses{def:bandit, def:class, def:kinf, def:kinfty, def:algorithm, cite:LS-16.2}
Let $\kappa>1$ and $\nu\in\Hk^\alpha$ with at least one \term{w:subopt}{suboptimal} arm. Then:
\begin{enumerate}[label=(\roman*)]
  \item $0<\Cst(\nu)<\infty$, and every algorithm \term{w:consistent}{consistent} over $\Hk^\alpha$ satisfies
  $\liminf_T\Reg_T/\log T\ge\Cst(\nu)$;
  \item \algo{} is \term{w:consistent}{consistent} over $\Hk^\alpha$, and
  \begin{equation*}
    \lim_{T\to\infty}\frac{\Reg_T}{\log T}=\Cst(\nu)=\sum_{a:\gap_a>0}\frac{\gap_a}{\Kc(\nu_a,\mstar)};
  \end{equation*}
  \item for every \term{w:subopt}{suboptimal} arm $a$, $\gap_a/\Kc(\nu_a,\mstar)\ge\gap_a/\Kinfty(\kappa)>\kappa\gap_a$;
  \item $\Kinf(P,\mu;\Hlt)=\Kc(P,\mu)$ for every \term{w:law}{law} $P$ and $\mu\in\R$, and (i)--(ii) hold for every
  $\nu\in(\Hlt)^\alpha$ with $\Hlt$ in place of $\Hk$, for the same algorithm \algo{}.
\end{enumerate}
\end{corollary}

(See the proof in Appendix~\ref{sec:proof-main}.)

\paragraph{Comparison with the lower bound} Items (i)--(ii) pair the upper bound of Theorem~\ref{thm:main}
with the lower bound of Theorem~\ref{cite:LS-16.2} in the same class, the algorithms \term{w:consistent}{consistent} over
$\Hk^\alpha$: the first-order constant $\Cst(\nu)$ cannot be improved, and \algo{} attains it on every
\term{w:instance}{instance}. This answers, at first order, the question of \citet[Section~4]{L17}.

\paragraph{Comparison with the algorithm of Lattimore} Lattimore's bound
$C\sum_a\gap_a(\kappa-1+\sda_a^2/\gap_a^2)$ (Lemma~\ref{cite:L17-thm2}) has the right shape: by (iii) every
arm costs more than $\kappa\gap_a$, and the numerical values of Section~\ref{sec:examples} suggest that
$\gap_a/\Kc(\nu_a,\mstar)$ is within a factor between $1.4$ and $2.1$ of $\gap_a(\kappa-1+\sda_a^2/\gap_a^2)$
on the \term{w:law}{laws} we tried. For an arm with a fixed standardised \term{w:law}{law} and $\gap_a/\sda_a\to\infty$,
Corollary~\ref{cor:largegap} and Proposition~\ref{prop:kinfty}(iv) give the exact coefficient
\begin{equation*}
  \frac{\gap_a}{\Kc(\nu_a,\mstar)}\sim\frac{\gap_a}{\Kinfty(\kappa)}\quad\Bigl(\frac{\gap_a}{\sda_a}\to\infty\Bigr),
  \qquad\text{where}\qquad
  \frac{\gap_a}{\Kinfty(\kappa)}=\Bigl(\kappa+\frac32+O\Bigl(\frac1\kappa\Bigr)\Bigr)\gap_a ,
\end{equation*}
which is, in this regime, the exact form of the term $(\kappa-1)\gap_a$.

\paragraph{Comparison with KL$_{\inf}$-UCB and with the generic theorem} The algorithm and its threshold are
those of \citet{AJK21}, and the regret decomposition is the classical one; the class is different, and the
two deviation inequalities that the decomposition needs are proved in Theorem~\ref{thm:deviation}. The upper
one is the form that \citet{PA26} assume (Lemma~\ref{cite:PA26-thm3.1}); the lower one is weaker than theirs
(an existential rate, and a level shifted by $\delta$), so \citet[Theorem~3.1]{PA26} cannot be applied as such,
and Appendix~\ref{sec:proof-main} gives a self-contained proof.

\paragraph{Nature of the constants} The first-order constant is exact. The constant $\Ca(\varepsilon)$ of (i) is
existential: it comes from the lower deviation inequality, whose rate is obtained by a topological argument
(Theorem~\ref{thm:deviation}(ii)); it depends on $\nu_a$, $\mstar$, $\varepsilon$ and $\kappa$, and we have no
estimate of its size. Everything else in \eqref{eq:main-finite} is explicit and valid from $T\ge3$.

\paragraph{Scope} We do not give an explicit second-order term, a bound on the regret uniform over
\term{w:instance}{instances}, or a computationally efficient index. The number $\kappa$ must be known, which is necessary
\citep[Theorem~7]{L17}. Instances with several optimal arms are allowed.

% ----------------------------------------------------------------
\section{Concentration of the empirical information function}\label{sec:concentration}

The regret bound rests on two deviation inequalities for $\Kc$ evaluated at an empirical \term{w:law}{law}. They are of
independent use, for example for confidence sequences.

\begin{theorem}[Deviation inequalities]\label{thm:deviation}
\uses{def:class, def:kinf, def:laws}
Let $\kappa>1$, $P\in\Hk$, let $X_1,X_2,\dots$ be independent with \term{w:law}{law} $P$, and
$\hat P_n:=\frac1n\sum_{i\le n}\delta_{X_i}$. Then:
\begin{enumerate}[label=(\roman*)]
  \item (anytime upper deviation) for every $x\ge0$,
  \begin{equation}\label{eq:upper-dev}
    \Prob\bigl(\exists n\ge1:\ n\,\Kc(\hat P_n,\mean(P))\ge x+1+2\log(n+1)\bigr)\le e^{-x},
  \end{equation}
  and consequently $\Prob\bigl(n\Kc(\hat P_n,\mean(P))\ge x\bigr)\le e(n+1)^2e^{-x}$ for every $n\ge1$ and $x\in\R$;
  \item (lower deviation) for every $\mu\in\R$ and every $\varepsilon>0$ there are $\delta>0$, $J\in\mathbb N$ and
  $\eta>0$, depending only on $P$, $\mu$, $\varepsilon$ and $\kappa$, such that for every $n\ge1$
  \begin{equation}\label{eq:lower-dev}
    \Prob\bigl(\Kc(\hat P_n,\mu-\delta)\le\Kc(P,\mu)-\varepsilon\bigr)\le2J\,e^{-2n\eta^2}.
  \end{equation}
\end{enumerate}
The events are measurable by Lemma~\ref{lem:index}(i).
\end{theorem}

(See the proof in Appendix~\ref{sec:proof-deviation}.)

\begin{corollary}[Scale-free confidence sequence]\label{cor:confseq}
\uses{def:class, def:kinf}
Let $\kappa>1$, $P\in\Hk$, $X_1,X_2,\dots$ and $\hat P_n$ as in Theorem~\ref{thm:deviation}, $\hat P^-_n$ the
empirical \term{w:law}{law} of $-X_1,\dots,-X_n$, $x\ge0$ and $b_n:=x+1+2\log(n+1)$. Let
\begin{equation*}
  U^+_n:=\sup\{\mu:n\Kc(\hat P_n,\mu)\le b_n\},\qquad U^-_n:=\sup\{\mu:n\Kc(\hat P^-_n,\mu)\le b_n\}.
\end{equation*}
Then
\begin{equation*}
  \Prob\bigl(\forall n\ge1:\ -U^-_n\le\mean(P)\le U^+_n\bigr)\ge1-2e^{-x}.
\end{equation*}
\end{corollary}

(See the proof in Appendix~\ref{sec:proof-deviation}.)

\paragraph{Comparison with the heavy-tailed class} Inequality \eqref{eq:upper-dev} has the same form and
the same constants as \citet[Proposition~5]{AJK21} for $\{\E|X|^{1+\beta}\le B\}$. The lower deviation inequality of
\citet[Lemma~6]{AJK21} has the explicit rate $e^{-nc\delta^2}$; ours is existential. The confidence sequence
uses only the bound on the kurtosis: it adapts to the unknown scale.

\paragraph{Comparison with confidence intervals under a kurtosis bound} For a known bound on the kurtosis and an
unknown variance $v$, \citet[Section~4.2, Proposition~4.5]{Cat12} builds an estimator of the mean, from a
preliminary estimate of the variance \citep[Corollary~4.3]{Cat12}, whose error is at most an explicit quantity
equivalent, for large $n$, to $\sqrt{2v\log(1/\epsilon)/n}$, with probability at least $1-\epsilon$ split between the two
steps. That bound holds for a sample size $n$ and a level $\epsilon$ fixed in advance, under a condition
$n\gtrsim36(\kappa-1)\log(1/\epsilon)$ (stated for $\kappa\ge3$). Corollary~\ref{cor:confseq} holds simultaneously for all
$n\ge1$, for every $\kappa>1$, with the explicit threshold $b_n$, but we give no closed form for the width of the
interval $[-U^-_n,U^+_n]$.

\paragraph{Why the lower deviation is harder} The upper deviation only needs an upper bound on
$\Kc(\hat P_n,\mean(P))$, and fixing the location and the scale at the true mean and standard deviation in
\eqref{eq:dual} gives one. The lower deviation needs a lower bound on an infimum over all $(c,s)$, which may
be approached at infinity; this is where the escape of mass and the value $\Kinfty(\kappa)$ enter
(Section~\ref{sec:overview}).

% ----------------------------------------------------------------
\section{How the proofs work}\label{sec:overview}

The proof of Theorem~\ref{thm:main} has five steps; the first four concern the information function alone.

\paragraph{Step 1: convexification} A \term{w:law}{law} $Q$ has kurtosis at most $\kappa$ and mean $c$ if and only if
$\tauc_{c,s}Q\in\Gk$ for some $s>0$ (Lemma~\ref{lem:G}). Since $\KL$ is invariant under the bijection $\tauc_{c,s}$,
\begin{equation*}
  \Kc(P,\mu)=\inf_{c>\mu,\ s>0}\KG(\tauc_{c,s}P)
\end{equation*}
(Lemma~\ref{lem:repr}), and the strict and non-strict versions coincide (Lemma~\ref{lem:strict}).

\paragraph{Step 2: duality} For a finitely supported \term{w:law}{law} $R$, $\KG(R)=\max_{\Lam}\E_R\log\hl_\lambda$: weak duality
is Jensen's inequality (Lemma~\ref{lem:weak}); for the converse we build an optimal \term{w:law}{law} of $\Gk$ from a
maximiser of the dual and its normal cone (Lemma~\ref{lem:finite}). The extension to every \term{w:law}{law} uses the lower
semicontinuity of $\KG$ and Fatou's lemma (Lemma~\ref{lem:alllaws}).

\paragraph{Step 3: upper deviation} At the true mean $m$ and standard deviation $\sigma$ of $P\in\Hk$, the
choice $(c,s)=(m,\sigma)$ in Step 1 gives $n\Kc(\hat P_n,m)\le\max_{\Lam}\sum_{i\le n}\log\hl_\lambda(Y_i)$ with
$Y_i:=(X_i-m)/\sigma$ and $\E\hl_\lambda(Y_i)\le1$. A mixture of the supermartingales $\prod_i\hl_\lambda(Y_i)$ over
the uniform \term{w:law}{law} on the fixed set $\Lam$ and the maximal inequality give \eqref{eq:upper-dev}.

\paragraph{Step 4: the value at infinity and the escape of mass} This is the new step. To bound
$\Kc(P_n,\mu_n)$ from below when $P_n\to P$, take near-optimal $(c_n,s_n)$ in Step 1. If they stay in a
compact subset of $\R\times(0,\infty)$, the lower semicontinuity of $\KG$ concludes. Otherwise the
normalised \term{w:law}{laws} $\tauc_{c_n,s_n}P_n$ either spread their mass to infinity, which costs an unbounded amount
because every \term{w:law}{law} of $\Gk$ has a fourth moment at most $4\kappa^2$, or concentrate near a point, which costs
at least $\Kinfty(\kappa)$ because no \term{w:law}{law} of $\Hk$ has an atom larger than $\qk$. The following lemma states
this dichotomy.

\begin{lemma}[Compactness modulo the affine group]\label{lem:escape}
\uses{def:normalised, def:kinfty}
Let $\kappa>1$, let $P$ be a \term{w:law}{law} that is not a Dirac mass, $P_n\to P$ weakly, and
$(c_n,s_n)\in\R\times(0,\infty)$ with $\inf_nc_n>-\infty$; let $R_n:=\tauc_{c_n,s_n}P_n$. Then there is a
subsequence along which one of the following holds:
\begin{enumerate}[label=(\alph*)]
  \item $(c_n,s_n)\to(c,s)\in\R\times(0,\infty)$ and $\liminf_n\KG(R_n)\ge\KG(\tauc_{c,s}P)$;
  \item $\liminf_n\KG(R_n)\ge\Kinfty(\kappa)$.
\end{enumerate}
\end{lemma}

(See the proof in Appendix~\ref{sec:proof-regularity}.)

Since $\Kc(P,\mu)\le\Kinfty(\kappa)$ when $P$ has a finite fourth moment (Proposition~\ref{prop:kinfty}(ii)), both
alternatives give $\liminf\Kc(P_n,\mu_n)\ge\Kc(P,\mu)$: this is Proposition~\ref{prop:regularity}(i). The two
bounds meet exactly: the value $\Kinfty(\kappa)$ of the escape is the supremum of $\Kc$ over \term{w:law}{laws} with a finite fourth
moment, so no weaker
information on the escape would suffice.

\paragraph{Step 5: lower deviation and regret} Lower semicontinuity at $(P,\mu)$ in the metrisable weak
topology gives finitely many bounded continuous test functions whose empirical means control
$\Kc(\hat P_n,\mu-\delta)$ from below; Hoeffding's inequality gives \eqref{eq:lower-dev}. The regret analysis
then follows the classical decomposition of KL-UCB: rounds where the optimal arm's index falls below
$\mstar$, controlled by \eqref{eq:upper-dev}; rounds where a \term{w:subopt}{suboptimal} arm's index is above $\mstar$,
counted over the sample sizes and controlled by \eqref{eq:lower-dev}.

\paragraph{Where the difficulty lies} The convexification and the upper deviation are short once the
normalised class is found. The difficulty is the lower bound on $\Kc$ near a \term{w:law}{law}, because the class is
scale free: the infimum over $(c,s)$ need not be attained, and near-optimal parameters can run to infinity.
Lemma~\ref{lem:escape}, with the largest atom of Proposition~\ref{prop:kinfty}(i), is what controls this; without
it, neither the positivity of $\Kc$ nor the lower deviation would follow.

\paragraph{Map of the appendices} Appendix~\ref{sec:cited} collects the results from other papers.
Appendix~\ref{sec:proof-dual} proves Proposition~\ref{prop:dual} (Steps 1--2), Appendix~\ref{sec:proof-kinfty}
Proposition~\ref{prop:kinfty}, Appendix~\ref{sec:proof-regularity} Lemma~\ref{lem:escape},
Proposition~\ref{prop:regularity} and Corollary~\ref{cor:largegap} (Step 4), Appendix~\ref{sec:proof-index}
Lemma~\ref{lem:index}, Appendix~\ref{sec:proof-main} Theorem~\ref{thm:main} and Corollary~\ref{cor:exact}
(Step 5), and Appendix~\ref{sec:proof-deviation} Theorem~\ref{thm:deviation} and Corollary~\ref{cor:confseq}
(Steps 3 and 5). Appendix~\ref{sec:attribution} gives the sources of the techniques and
Appendix~\ref{sec:notation} the notation.

% ----------------------------------------------------------------
\section{Numerical examples}\label{sec:examples}

This section gives values of $\Kinfty(\kappa)$ and of $\Kc$, an \term{w:instance}{instance} that satisfies all the hypotheses, and the
degenerate cases.

\paragraph{Why no simulation of the algorithm} We did not simulate \algo{}. Its index \eqref{eq:index} requires
$\Kc(\hat P,\mu)$ for an empirical \term{w:law}{law} at every round, and by the paragraph ``Computation'' of
Section~\ref{sec:algorithm} the outer minimisation over $(c,s)$ in \eqref{eq:dual} is not convex, has a non-compact domain
and may be approached at infinity; we have no procedure that computes it to a certified accuracy. A simulation would
therefore run a heuristic approximation of the index, not the algorithm that Theorem~\ref{thm:main} analyses: an
uncontrolled error on $\Kc$ moves the index in either direction, and the regret observed would say nothing about the
theorem. Besides, a first-order constant is a limit as $T\to\infty$, which simulations at a finite horizon do not
show. We compute instead the quantities that the theorem involves:
$\Kinfty(\kappa)$ in closed form, and $\Kc$ on explicit \term{w:law}{laws}, where a direct search over alternatives gives an
independent upper bound.

\paragraph{Setup} The numbers below are computed from \eqref{eq:dual}. At fixed $(c,s)$, the maximisation over $\Lam$ uses
the description of Proposition~\ref{prop:dual}(iii): for $\lambda_2>0$ the infimum defining $\Lw(\lambda_2)$ is attained at
$y^2=(\kappa+\sqrt{\kappa^2+3\kappa+3/\lambda_2})/3$, the only positive critical point, and the problem becomes two nested
one-dimensional concave maximisations. The outer minimisation over $(c,s)$ uses a multi-start Nelder--Mead search in
$(\log(c-\mu),\log s)$ and a search on the boundary $c=\mu$. As an independent check, a direct search over \term{w:law}{laws}
$\sum_iw_i\delta_{x_i}+\beta_1\delta_{u_1}+\beta_2\delta_{u_2}$ with mean at least $\mu$ and kurtosis at most $\kappa$ (sequential quadratic
programming with random restarts) gives an upper bound on $\Kc$. These are heuristics without certified accuracy.
\emph{Reproduction.} The Python scripts (NumPy and SciPy) are distributed with the source of this paper, with their
outputs: \texttt{table1\_kinfty.py} gives Table~\ref{tab:kinfty} and the random search over four-atom \term{w:law}{laws} (a few
seconds), \texttt{table2\_duality.py} gives Table~\ref{tab:duality} and the check at fixed $(c,s)$ (a few minutes), and
\texttt{table3\_curves.py} gives Table~\ref{tab:curves} (about thirty minutes); the three share the routines of
\texttt{kinf\_tools.py}, and each runs without arguments and prints its table.

\begin{table}[t]
\centering
\small
\caption[The value at infinity]{The value at infinity $\Kinfty(\kappa)=\log(1/\qk)$, its inverse, and the large-gap branch
$\log\frac{\kappa}{\kappa-1}$ of \citet[Theorem~9]{L17} (valid there for $\kappa\ge7/2$).}\label{tab:kinfty}
\begin{tabular}{@{}lrrrrrrrr@{}}
\toprule
$\kappa$ & $1.5$ & $2$ & $3$ & $3.5$ & $5$ & $9$ & $20$ & $100$ \\
\midrule
$\qk$ & 0.66667 & 0.72361 & 0.78868 & 0.81009 & 0.85355 & 0.90825 & 0.95445 & 0.99020 \\
$\Kinfty(\kappa)$ & 0.40547 & 0.32351 & 0.23740 & 0.21061 & 0.15835 & 0.09624 & 0.04662 & 0.00985 \\
$1/\Kinfty(\kappa)$ & 2.466 & 3.091 & 4.212 & 4.748 & 6.315 & 10.391 & 21.448 & 101.489 \\
$\log\frac{\kappa}{\kappa-1}$ & 1.09861 & 0.69315 & 0.40547 & 0.33647 & 0.22314 & 0.11778 & 0.05129 & 0.01005 \\
\bottomrule
\end{tabular}
\end{table}

\paragraph{Interpretation of Table~\ref{tab:kinfty}} The inverse $1/\Kinfty(\kappa)$, the least cost per unit gap and per
$\log T$ of a \term{w:subopt}{suboptimal} arm, lies in $(\kappa,\kappa+2)$, and $1/\Kinfty(\kappa)-\kappa$ tends to $3/2$ (it is $0.966$ at $\kappa=1.5$, $1.212$ at
$\kappa=3$ and $1.4989$ at $\kappa=10^3$). A random search over \term{w:law}{laws} with four atoms and kurtosis at most $\kappa$ finds the
largest atom $\qk$ to six digits for $\kappa\in\{1.5,2.5,3,6\}$.

\begin{table}[t]
\centering
\small
\caption[Check of the dual formula]{Check of the dual formula. Left: $P$ uniform on $\{-1,0,1\}$ (kurtosis $3/2$), $\kappa=3$; the profiled dual
\eqref{eq:dual} and a direct search over alternatives. Right: the two-point law $q_3\delta_0+(1-q_3)\delta_1$ (kurtosis exactly
$3$, variance $1/6$) at the level $m+\Delta$, and the ratio of $\Kc$ to $\Delta^2/(2\sigma^2)$.}\label{tab:duality}
\begin{tabular}{@{}lrrrrr@{}}
\toprule
$\mu$ & $0.1$ & $0.3$ & $0.6$ & $1.0$ & $2.0$ \\
\midrule
dual \eqref{eq:dual} & 0.007492 & 0.066107 & 0.164736 & 0.205489 & 0.228409 \\
direct search & 0.007492 & 0.066107 & 0.164736 & 0.205489 & 0.228409 \\
\bottomrule
\end{tabular}
\quad
\begin{tabular}{@{}lrrrr@{}}
\toprule
$\Delta$ & $0.01$ & $0.2$ & $1$ & $5$ \\
\midrule
$\Kc(P,m+\Delta)$ & 0.000293 & 0.082477 & 0.227261 & 0.237014 \\
ratio & 0.978 & 0.687 & 0.076 & 0.003 \\
\bottomrule
\end{tabular}
\end{table}

\paragraph{Interpretation of Table~\ref{tab:duality}} The two computations agree to six digits at every level, on a \term{w:law}{law}
inside the class and on a \term{w:law}{law} on its boundary. In every case the optimal location is $c=\mu$, and the direct search
adds one atom far to the right of the data (mass $0.20$ at $9.78$ for $\mu=2$), as Proposition~\ref{prop:dual}(vi)
allows. At fixed $(c,s)$, on four random \term{w:law}{laws} with four atoms, the maximum over $\Lam$ and the minimum over $\Gk$
also agree to seven digits. As $\Delta$ grows, $\Kc(P,m+\Delta)$ increases to $\Kinfty(3)=0.23740$ (Corollary~\ref{cor:largegap}).

\begin{table}[t]
\centering
\small
\caption[The information function on two standardised laws]{$\Kc(P,\Delta)$ for two standardised laws (mean $0$, variance $1$): the $41$-point quantile discretisation
of $\mathcal N(0,1)$ (kurtosis $2.690$) and the $21$-point uniform grid (kurtosis $1.795$). The last columns compare the
exact regret coefficient $\Delta/\Kc$ of an arm with gap $\Delta$ and unit variance (Gaussian column) with the bracket
$\Delta(\kappa-1+1/\Delta^2)$ of \citet[Theorem~2]{L17}, without its universal constant.}\label{tab:curves}
\begin{tabular}{@{}rrrrrrrr@{}}
\toprule
 & $\Delta$ & $\Kc$ (Gauss.) & $\frac{\Kc}{\Delta^2/2}$ & $\Kc$ (unif.) & $\Delta/\Kc$ & $\Delta(\kappa-1+\Delta^{-2})$ & ratio \\
\midrule
$\kappa=3$ & 0.02 & 0.000200 & 0.9998 & 0.000200 & 100.016 & 50.040 & 1.999 \\
 & 0.05 & 0.001249 & 0.9990 & 0.001250 & 40.041 & 20.100 & 1.992 \\
 & 0.1 & 0.004979 & 0.9959 & 0.004992 & 20.083 & 10.200 & 1.969 \\
 & 0.25 & 0.030420 & 0.9734 & 0.030902 & 8.218 & 4.500 & 1.826 \\
 & 0.5 & 0.108498 & 0.8680 & 0.111946 & 4.608 & 3.000 & 1.536 \\
 & 1.0 & 0.193042 & 0.3861 & 0.192677 & 5.180 & 3.000 & 1.727 \\
 & 2.0 & 0.224303 & 0.1122 & 0.224233 & 8.917 & 4.500 & 1.981 \\
 & 4.0 & 0.233935 & 0.0292 & 0.233929 & 17.099 & 8.250 & 2.073 \\
 & 8.0 & 0.236521 & 0.0074 & 0.236520 & 33.824 & 16.125 & 2.098 \\
\midrule
$\kappa=5$ & 0.02 & 0.000200 & 0.9997 & 0.000200 & 100.033 & 50.080 & 1.997 \\
 & 0.05 & 0.001247 & 0.9979 & 0.001248 & 40.084 & 20.200 & 1.984 \\
 & 0.1 & 0.004958 & 0.9916 & 0.004970 & 20.170 & 10.400 & 1.939 \\
 & 0.25 & 0.029451 & 0.9424 & 0.029846 & 8.489 & 5.000 & 1.698 \\
 & 0.5 & 0.086715 & 0.6937 & 0.087192 & 5.766 & 4.000 & 1.441 \\
 & 1.0 & 0.131888 & 0.2638 & 0.131874 & 7.582 & 5.000 & 1.516 \\
 & 2.0 & 0.150472 & 0.0752 & 0.150469 & 13.291 & 8.500 & 1.564 \\
 & 4.0 & 0.156265 & 0.0195 & 0.156264 & 25.598 & 16.250 & 1.575 \\
 & 8.0 & 0.157819 & 0.0049 & 0.157819 & 50.691 & 32.125 & 1.578 \\
\bottomrule
\end{tabular}
\end{table}

\paragraph{Interpretation of Table~\ref{tab:curves}} Four observations, on these two \term{w:law}{laws} only. First, for small gaps
$\Kc\approx\Delta^2/(2\sd^2)$: the ratio tends to $1$, so the optimal regret per arm is close to $2\sda_a^2/\gap_a$, twice the term
$\sda_a^2/\gap_a$ of Lattimore's bracket (we do not prove this; see Open question~\ref{oq:smallgap}). Second, from
$\Delta\approx\sd$ on, $\Kc$ is within $20\%$ of $\Kinfty(\kappa)$, and $\Delta/\Kc$ approaches $\Delta/\Kinfty(\kappa)$, which is
$(\kappa+1.21)\Delta$ for $\kappa=3$ and $(\kappa+1.32)\Delta$ for $\kappa=5$. Third, over the whole range, the exact coefficient is between $1.4$ and $2.1$ times the bracket of
\citet[Theorem~2]{L17}, which suggests that this bracket has the right shape. Fourth, the shape of the \term{w:law}{law} matters
little once $\Delta\gtrsim\sd$: the two columns agree to three or four digits, as Corollary~\ref{cor:largegap} predicts
in the limit.

\paragraph{An instance satisfying all the hypotheses} Let $\kappa:=3$ and two arms $\nu_1:=\mathcal N(1,1)$, $\nu_2:=\mathcal N(0,1)$. Both
have kurtosis $3$, so $\nu\in\mathcal H_3^\alpha$, with $\gap_2=1$. By Corollary~\ref{cor:exact}, \algo{} run with $\kappa=3$ has
$\lim_TR_T/\log T=1/\Kc(\mathcal N(0,1),1)$, and no algorithm \term{w:consistent}{consistent} over $\mathcal H_3^\alpha$ does better. By
Proposition~\ref{prop:kinfty}(ii), $\Kc(\mathcal N(0,1),1)\le\Kinfty(3)=0.2374$, so this constant is at least $4.21$; the constant
of the smaller model of Gaussian \term{w:law}{laws} with unknown mean and variance is $1/(\frac12\log2)=2.885$ \citep[equations~(2)--(3)]{HT14}.
Knowing only that the kurtosis is at most $3$, instead of knowing that the arms are Gaussian, therefore costs at
least a factor $1.46$ on this \term{w:instance}{instance}; the $41$-point discretisation of Table~\ref{tab:curves} suggests a factor
of about $1.8$.

\paragraph{Degenerate and boundary cases}
\begin{itemize}
  \item \emph{Bernoulli arms.} A Bernoulli \term{w:law}{law} of parameter $\theta$ has kurtosis $1/(\theta(1-\theta))-3$; it belongs to $\Hk$ if and
  only if $\theta\in[1-\qk,\qk]$, and the extreme \term{w:law}{laws} carry the largest atom of Proposition~\ref{prop:kinfty}(i). Such an arm
  produces a Dirac empirical \term{w:law}{law} with probability at most $\qk^{n-1}$ after $n$ pulls, and its index is then $\pm\infty$
  (Lemma~\ref{lem:index}(iii)).
  \item \emph{Gaussian arms.} For $\kappa\ge3$ every Gaussian \term{w:law}{law} belongs to $\Hk$, and
  $\Kc(\mathcal N(m,\sigma^2),m+\Delta)\le\min\{\frac12\log(1+\Delta^2/\sigma^2),\Kinfty(\kappa)\}$: the first term is $\Kinf$ over the smaller
  Gaussian model, the second is Proposition~\ref{prop:kinfty}(ii).
  \item \emph{Laws on the boundary} $\kurt(P)=\kappa$, such as Gaussian arms with $\kappa=3$ or the extreme Bernoulli \term{w:law}{laws}, are
  covered by all the results: Theorem~\ref{thm:deviation}(i) uses only $\E\phik(Y)\le0$, and
  Proposition~\ref{prop:regularity} only a finite fourth moment and a \term{w:law}{law} that is not a Dirac mass.
  \item \emph{The parameter $\kappa$.} As $\kappa\downarrow1$, $\qk\to\frac12$ and $\Kinfty(\kappa)\to\log2$; for $\kappa=1$ the class would only
  contain symmetric two-point \term{w:law}{laws}, for which $\Kc=+\infty$ above the mean. As $\kappa\to\infty$, $\Kinfty(\kappa)\sim1/\kappa\to0$: the forced
  exploration grows linearly in $\kappa$, in line with the impossibility of logarithmic regret uniformly over all the
  classes \citep[Theorem~7]{L17}.
\end{itemize}

% ----------------------------------------------------------------
\section{Conclusion}\label{sec:conclusion}

On the class of \term{w:law}{laws} with kurtosis at most a known $\kappa$, the first-order regret constant of \term{w:consistent}{consistent} algorithms
is the Burnetas--Katehakis constant $\sum_a\gap_a/\Kc(\nu_a,\mstar)$, and KL-UCB with the information function of the
class attains it. The information function has a dual formula valid for every \term{w:law}{law}, its largest value over
\term{w:law}{laws} with a finite fourth moment is the explicit $\Kinfty(\kappa)$, which governs the large-gap regime, and it is lower semicontinuous away from Dirac
masses. The limitations are the existential second-order term, the absence of a certified computation of the
index, and the absence of a closed form for $\Kc$; the questions of Section~\ref{sec:open} start from them.

\section{Open questions}\label{sec:open}

Each question below is stated with its class and criterion, and says where the method of this paper stops.

\begin{openquestion}[Explicit finite-time bounds]\label{oq:finite}
\extends{D8, D7}
For every $\kappa>1$ and every $P\in\Hk$ with $\mu>\mean(P)$, are there $d_0>0$, $C<\infty$ and $c>0$, given by explicit
functions of $\kappa$, $\kurt(P)$ and $(\mu-\mean(P))/\sd(P)$, such that
$\Prob(\Kc(\hat P_n,\mu)\le\Kc(P,\mu)-d)\le Ce^{-cnd^2}$ for all $n\ge1$ and $d\in(0,d_0)$? This would give, for \algo{} on every
$\nu\in\Hk^\alpha$, $\E[\Np_a(T)]\le\log T/\Ka+O((\log T)^{2/3})$ with explicit constants, as \citet[Theorem~1]{AJK21} for a bounded
moment. Our method stops at Theorem~\ref{thm:deviation}(ii) because the lower bound on $\Kc$ comes from a
compactness argument (Lemma~\ref{lem:escape}) with no modulus: one needs a quantitative version, saying how
fast near-optimal locations and scales can escape for a \term{w:law}{law} close to $P$.
\end{openquestion}

\begin{openquestion}[Efficient computation]\label{oq:computation}
\extends{D8}
Is there an algorithm whose index can be computed to a certified accuracy in time polynomial in the number of
samples, or that recomputes its indices at only $o(T)$ rounds, and that still satisfies
$\lim_TR_T/\log T=\Cst(\nu)$ for every $\nu\in\Hk^\alpha$? The inner problem of \eqref{eq:dual} is a two-dimensional concave
program over a fixed set, but the outer problem over $(c,s)$ is not convex, its domain is not compact, and its
infimum may be approached at infinity. A first step would be a confinement lemma for empirical \term{w:law}{laws}: near-optimal
$(c,s)$ lie in an explicit compact set, or the value there is at least $\Kinfty(\kappa)-\zeta$ for a given $\zeta>0$; the error of a grid would then
be added to the threshold.
\end{openquestion}

\begin{openquestion}[The shape of the information function]\label{oq:smallgap}
\extends{D7}
Is it true that, for every $\kappa>1$ and every $P\in\Hk$ with $\kurt(P)<\kappa$,
$\Kc(P,\mean(P)+\Delta)=\frac{\Delta^2}{2\sd^2(P)}(1+o(1))$ as $\Delta\downarrow0$? What happens on the boundary $\kurt(P)=\kappa$, where
Table~\ref{tab:duality} suggests the same behaviour for one \term{w:law}{law}? More globally, one half of a two-sided bound
is known. Lattimore's algorithm is \term{w:consistent}{consistent} over $\Hk^\alpha$ by Lemma~\ref{cite:L17-thm2}; on the
\term{w:instance}{instance} with two arms, of \term{w:law}{laws} $P$ and its translate by $\Delta>0$, Theorem~\ref{cite:LS-16.2} and
Lemma~\ref{cite:L17-thm2} give $\Delta/\Kc(P,\mean(P)+\Delta)\le C\Delta(\kappa-1+\sd^2(P)/\Delta^2)$, hence
$\Kc(P,\mean(P)+\Delta)\ge(2C)^{-1}\min\{\Delta^2/\sd^2(P),1/\kappa\}$ for every $\kappa>1$, $P\in\Hk$ and $\Delta>0$, with the universal
constant $C$ of Lemma~\ref{cite:L17-thm2}. For $\Delta\ge\sd(P)/\sqrt\kappa$, Proposition~\ref{prop:kinfty} gives
$\Kc(P,\mean(P)+\Delta)\le\Kinfty(\kappa)<1/\kappa\le\Delta^2/\sd^2(P)$. Is there a universal constant $C'$ with
$\Kc(P,\mean(P)+\Delta)\le C'\Delta^2/\sd^2(P)$ for every $\kappa>1$, $P\in\Hk$ and $\Delta\in(0,\sd(P)/\sqrt\kappa)$? Lemma~\ref{cite:L17-thm9}
gives this only for $\kappa\ge7/2$, $\kurt(P)<\kappa$ and $\Delta/\sd(P)$ below a threshold depending on $\kappa$ and $\kurt(P)$ (its
first case, with its own constants). Our method gives no bound on $\Kc$ that is
quantitative in $\Delta$ for small gaps: Proposition~\ref{prop:regularity}(iii) is proved by contradiction, and the
far-atom construction of Proposition~\ref{prop:kinfty} costs a constant.
\end{openquestion}

\begin{openquestion}[Randomised and simpler index algorithms]\label{oq:randomised}
\extends{D10}
Is there a non-parametric Thompson-type algorithm, or an IMED-type index algorithm, pulling an arm that
minimises $\Np_a(t)\Kc(\Pe_a(t),\hat m^\star(t))+\log\Np_a(t)$ with $\hat m^\star(t)$ the largest empirical mean, that satisfies
$\limsup_TR_T/\log T\le\Cst(\nu)$ for every $\nu\in\Hk^\alpha$? The first
is the second open problem of \citet[Section~4]{L17}. Our lower deviation inequality is uniform over a
neighbourhood of the level, which an IMED analysis at the random level $\hat m^\star(t)$ needs, but the events specific
to these algorithms were not checked.
\end{openquestion}

\begin{openquestion}[Classes convex up to location and scale]\label{oq:profiled}
\extends{D7}
Let $p>1$, $B>0$ and $\mathcal C_{p,B}:=\{P\in\Pspace:\E_P|X-\mean(P)|^p\le B\}$. Does KL-UCB with the information function
$\Kinf(\cdot,\cdot;\mathcal C_{p,B})$ and the threshold \eqref{eq:threshold} satisfy
$\limsup_T\E[\Np_a(T)]/\log T\le1/\Kinf(\nu_a,\mstar;\mathcal C_{p,B})$ for every \term{w:instance}{instance} $\nu\in\mathcal C_{p,B}^\alpha$ and every
\term{w:subopt}{suboptimal} arm $a$, with no known lower bound on the means? For MED, \citet[Corollary~6]{BSH23} prove first-order
optimality on the centred class $\{\E h(|X-m|)<B\}$, with a further integrability condition and a known lower bound on
the means, which they use for compactness; the reduction to an outer minimisation over the location is also in
\citet[Section~6.10]{A22}. Our method would profile the location only, and the analogue of Lemma~\ref{lem:escape} must
show that the near-optimal locations $c>\mu$ stay bounded for \term{w:law}{laws} close to the arm's law: a location $c\to\infty$
moves the whole mass of the normalised \term{w:law}{law} to $-\infty$, which we expect to cost $+\infty$ by a bound of the type
of Lemma~\ref{lem:tail}, and the KL-UCB analysis needs this only at the arm's law, not uniformly over the class; we
do not write this proof. The same question, for KL-UCB with both location and scale
profiled, is open for the ratio classes $\{\E|X-m|^p\le\kappa\sd^p\}$, $p>2$, and for bounds on skewness and kurtosis
together.
\end{openquestion}


% ================================================================ appendices
\appendix

\section{Useful lemmas}\label{sec:cited}

In this section we present the results from other papers that we use: the lower bound of
\citet{LS} is recalled in the main text as Theorem~\ref{cite:LS-16.2}; the four results below serve the
comparisons of Sections~\ref{sec:related}, \ref{sec:lower} and \ref{sec:upper} and Open question~\ref{oq:smallgap}, and
none of our proofs uses them. We also use by name a few textbook facts: the maximal inequality for non-negative
supermartingales \citep[Theorem~3.9]{LS} (Ville's inequality), Hoeffding's inequality for bounded
variables, the strong \term{w:law}{law} of large numbers, Varadarajan's theorem (the empirical \term{w:law}{laws} of an i.i.d.
sequence converge weakly almost surely), Prokhorov's theorem, Skorokhod's representation theorem, the
portmanteau theorem, Slutsky's lemma, the data processing inequality for $\KL$, and the
Donsker--Varadhan formula $\KL(R,Q)=\sup_f\{\int f\,dR-\log\int e^f\,dQ\}$ over bounded continuous $f$, which
makes $\KL$ jointly weakly lower semicontinuous.

\citet{L17} writes $n$ for the horizon, $k$ for the number of arms and $\mathrm{Kurt}$ for the kurtosis; his
algorithm is the one of \citet[Section~2]{L17}, based on median-of-means estimators, which we do not
describe.

% Verbatim: literature/notes/L17.md, Theorem 2 (restated for instances of the class H_kappa only).
\begin{lemma}[{From \citet[Theorem~2]{L17}}]\label{cite:L17-thm2}
\uses{def:bandit, def:class}
There is a universal constant $C>0$ such that, for every $\kappa>1$ and every \term{w:instance}{instance} $\nu\in\Hk^\alpha$, the
algorithm of \citet[Section~2]{L17} run with $\kappa$ satisfies
\begin{equation*}
  \limsup_{T\to\infty}\frac{\Reg_T}{\log T}\le C\sum_{a:\gap_a>0}\gap_a\Bigl(\kappa-1+\frac{\sd^2(\nu_a)}{\gap_a^2}\Bigr).
\end{equation*}
\end{lemma}

In \citet[Theorem~9]{L17}, $\mathcal H_{\kappa_\circ}$ is the class of \term{w:law}{laws} with kurtosis less than $\kappa_\circ$, our
$\Hlt$ with $\kappa=\kappa_\circ$, and $\kappa$ denotes the kurtosis of the \term{w:law}{law} $\nu$; we write $\kappa$ for the class
parameter and $k_P$ for the kurtosis of $P$. The statement prints the class of the alternatives with the
subscript of the kurtosis of $\nu$; its proof \citep[Appendix~A]{L17} builds alternatives of kurtosis at
most $\kappa_\circ$, and we restate it with alternatives in $\Hlt$, which is the weaker reading.

% Verbatim: literature/notes/L17.md, Theorem 9 (and the remark there on the class of the alternatives).
\begin{lemma}[{From \citet[Theorem~9]{L17}}]\label{cite:L17-thm9}
\uses{def:class, def:kinfD}
There are universal constants $C,C'>0$ such that, for every $\kappa\ge7/2$, every $\Delta>0$ and every \term{w:law}{law}
$P\in\Hlt$ with mean $m$, variance $\sigma^2>0$ and kurtosis $k_P$, with $p:=\min\{\Delta/\sigma,1/\kappa\}$,
\begin{equation*}
  \Kinf(P,m+\Delta;\Hlt)\le
  \begin{cases}
    \min\bigl\{\log\frac1{1-p},\ C'\Delta^2/\sigma^2\bigr\} & \text{if } Ck_P^{1/2}(k_P+1)\Delta/\sigma<\kappa,\\[2pt]
    \log\frac1{1-p} & \text{otherwise.}
  \end{cases}
\end{equation*}
\end{lemma}

\citet{AJK21} write $\mu$ for the \term{w:instance}{instance}, $\epsilon$ for our $\beta$, and use the non-strict
$\mathrm{KL}_{\inf}^{\mathcal L}(\eta,x):=\inf\{\KL(\eta,Q'):Q'\in\mathcal L,\ \mean(Q')\ge x\}$, where they write
$\eta$ for a law and the infimum is over laws $Q'$ of the class $\mathcal L$; their
$N_a(T)$ counts the pulls in the first $T-1$ rounds. KL$_{\inf}$-UCB is their Algorithm~1, with batch factor
$\tilde\eta\ge0$.

% Verbatim: literature/notes/AJK21.md, Corollary 2.
\begin{lemma}[{From \citet[Corollary~2]{AJK21}}]\label{cite:AJK21-cor2}
\uses{def:bandit, def:laws}
Let $\beta,B>0$ and $\mathcal L_B:=\{P\in\Pspace:\int|x|^{1+\beta}\,dP\le B\}$. For every $\nu\in\mathcal L_B^\alpha$, every
$\tilde\eta\ge0$ and every \term{w:subopt}{suboptimal} arm $a$, KL$_{\inf}$-UCB with batch factor $\tilde\eta$ and threshold
$g_a(t)=\log t+2\log\log t+2\log(1+N_a(t))+1$ satisfies
\begin{equation*}
  \limsup_{T\to\infty}\frac{\E[N_a(T)]}{\log T}\le\frac{1+\tilde\eta}{\mathrm{KL}_{\inf}^{\mathcal L_B}(\nu_a,\mstar)}.
\end{equation*}
\end{lemma}

\citet{PA26} write $\mathcal L$ for the class, $\hat\nu_n$ for the empirical \term{w:law}{law} of $n$ independent samples of
$\nu$ and use the non-strict $\mathrm{KL}_{\inf}^{\mathcal L}$ of Lemma~\ref{cite:AJK21-cor2}; generalized KL$_{\inf}$-UCB is
KL$_{\inf}$-UCB with batch factor $0$ and the threshold $f_a$ below.

% Verbatim: literature/notes/PA26.md, Assumptions 2.3-2.4 and Theorem 3.1.
\begin{lemma}[{From \citet[Theorem~3.1 with Assumptions~2.3--2.4]{PA26}}]\label{cite:PA26-thm3.1}
\uses{def:bandit, def:laws}
Let $\mathcal L$ be a set of \term{w:law}{laws} and $g(n):=C_1\log(1+n)+C_2$ with $C_1,C_2>0$. Assume that every
$\nu\in\mathcal L$, with $\hat\nu_n$ the empirical \term{w:law}{law} of $n$ independent samples of $\nu$, satisfies
\begin{enumerate}[label=(\alph*)]
  \item for every $x\ge0$, $\Prob\bigl(\exists n\in\mathbb N:\ n\,\mathrm{KL}_{\inf}^{\mathcal L}(\hat\nu_n,\mean(\nu))-g(n)\ge x\bigr)\le e^{-x}$;
  \item for every $\delta>0$ there are $d_0>0$ and $c_\nu>0$ such that, for every $d<d_0$ and every $n$,
  $\Prob\bigl(\mathrm{KL}_{\inf}^{\mathcal L}(\hat\nu_n,\mean(\nu)+\delta)\le\mathrm{KL}_{\inf}^{\mathcal L}(\nu,\mean(\nu)+\delta)-d\bigr)\le e^{-nc_\nu d^2}$.
\end{enumerate}
Then for every $\nu\in\mathcal L^\alpha$ and every \term{w:subopt}{suboptimal} arm $a$, generalized KL$_{\inf}$-UCB with
$f_a(t)=\log t+2\log\log t+g(N_a(t))$ satisfies
$\limsup_{T\to\infty}\E[N_a(T)]/\log T\le1/\mathrm{KL}_{\inf}^{\mathcal L}(\nu_a,\mstar)$.
\end{lemma}

% ----------------------------------------------------------------
\section{\texorpdfstring{Proof of Proposition~\ref{prop:dual}}{Proof of Proposition 4.2}}\label{sec:proof-dual}

Throughout, $\kappa>1$ is fixed. For $c\in\R$ and $s>0$ we write $x\mapsto c+sx$ for the inverse of $\tauc_{c,s}$,
and $\dt{s:gamma}\gam(y):=(y,\phik(y))\in\R^2$, so that $\hl_\lambda(y)=1+\lambda\cdot\gam(y)$ is affine in $\lambda$. To prove
Proposition~\ref{prop:dual} we use the following eight lemmas; the first three are the convexification
(Step~1 of Section~\ref{sec:overview}), the last five the duality (Step~2).

\begin{lemma}[The normalised class]\label{lem:G}
\uses{def:class, def:kinf, def:normalised}
\begin{enumerate}[label=(\alph*)]
  \item Every $Q\in\Gk$ has $v_2:=\int y^2\,dQ\in[\kappa-\sqrt{\kappa^2-\kappa},\ \kappa+\sqrt{\kappa^2-\kappa}]$,
  $\int y^4\,dQ\le2\kappa v_2-\kappa<4\kappa^2$ and $\kurt(Q)\le\kappa$; so $\Gk\subset\Hk$, and every \term{w:law}{law} of $\Gk$ has mean $0$.
  \item For every \term{w:law}{law} $Q$ and every $c\in\R$: $Q\in\Hk$ and $\mean(Q)=c$ if and only if $\tauc_{c,s}Q\in\Gk$ for some $s>0$;
  and then $s=\sd(Q)$ is such an $s$.
  \item $\Gk$ is convex and weakly compact.
  \item For every \term{w:law}{law} $P$, $c\in\R$, $s>0$ and $\mu\in\R$, with $P'$ the image of $P$ under $x\mapsto c+sx$,
  $\Kc(P',c+s\mu)=\Kc(P,\mu)$.
\end{enumerate}
\end{lemma}

(See the proof in Subsection~\ref{sec:proof-lem-G}.) This lemma is elementary. For a centred moment, the
analogous reduction with the location only is in \citet[Section~6.10]{A22} and \citet[Appendix~D.5]{BSH23};
treating the location and the scale as nuisance parameters of moment equations is classical in empirical
likelihood \citep[see also][Sections~1.1 and 3.5]{O01}. The form used here, the inequality $\kurt\le\kappa$ written as
the linear constraint $\int\phik\,dQ\le0$ with a free scale, we did not find.

\begin{lemma}[Representation]\label{lem:repr}
\uses{def:kinf, def:normalised}
For every \term{w:law}{law} $P$ and every $\mu\in\R$,
\begin{equation*}
  \Kge(P,\mu)=\inf_{c\ge\mu,\ s>0}\KG(\tauc_{c,s}P),\qquad
  \Kc(P,\mu)=\inf_{c>\mu,\ s>0}\KG(\tauc_{c,s}P).
\end{equation*}
\end{lemma}

(See the proof in Subsection~\ref{sec:proof-lem-repr}.) This is Lemma~\ref{lem:G}(b) and the invariance of
$\KL$ under bijections; the outer minimisation over a parameter is the device of
\citet[Theorem~3.4]{AKJ21}, \citet[Section~6.10]{A22} and \citet[Appendix~D.5]{BSH23}, here with two parameters
whose domain $\{(c,s):c>\mu,\ s>0\}$ is not compact.

\begin{lemma}[Strict and non-strict versions]\label{lem:strict}
\uses{def:kinf, def:class}
\begin{enumerate}[label=(\alph*)]
  \item For every \term{w:law}{law} $P$ and $\mu\in\R$, $\Kc(P,\mu)=\Kge(P,\mu)$; hence $\mu\mapsto\Kc(P,\mu)$ is non-decreasing and
  right-continuous.
  \item For every \term{w:law}{law} $P$ and $\mu\in\R$, $\Kinf(P,\mu;\Hlt)=\Kc(P,\mu)$.
\end{enumerate}
\end{lemma}

(See the proof in Subsection~\ref{sec:proof-lem-strict}.) The mixing argument is classical; it replaces the
continuity in the level used for the same purpose by \citet[Lemma~2]{HT10}.

\begin{lemma}[Lower semicontinuity of the normalised problem]\label{lem:lsc}
\uses{def:normalised, lem:G}
$\KG:\Pspace\to[0,\infty]$ is weakly lower semicontinuous, and when $\KG(R)<\infty$ there is $Q\in\Gk$ with
$\KL(R,Q)=\KG(R)$.
\end{lemma}

(See the proof in Subsection~\ref{sec:proof-lem-lsc}.) Classical: the joint lower semicontinuity of $\KL$
and the compactness of $\Gk$.

\begin{lemma}[The dual set]\label{lem:lambda}
\uses{def:dualset}
\begin{enumerate}[label=(\alph*)]
  \item $\Lam$ is closed and convex, and it is contained in the rectangle
  $[-\kappa^{-1/2},\kappa^{-1/2}]\times[0,1/(\kappa(\kappa-1))]$; hence it is compact.
  \item $\Lam=\{(\lambda_1,\lambda_2):0\le\lambda_2\le\frac1{\kappa(\kappa-1)},\ |\lambda_1|\le\Lw(\lambda_2)\}$, and $\Lw$ is concave and
  positive on $(0,\frac1{\kappa(\kappa-1)})$; hence $\Lam$ has non-empty interior.
  \item For every $\lambda\in\mathrm{int}\,\Lam$: $\lambda_2>0$ and $\inf_{y\in\R}\hl_\lambda(y)>0$.
  \item $C_\Lam:=\sup_{\lambda\in\Lam,\ y\in\R}\hl_\lambda(y)/(1+y^4)<\infty$.
\end{enumerate}
\end{lemma}

(See the proof in Subsection~\ref{sec:proof-lem-lambda}.) Elementary.

\begin{lemma}[Weak duality]\label{lem:weak}
\uses{def:normalised, def:dualset}
Let $R$ be the \term{w:law}{law} of $Y$, $Q\in\Gk$ with $\KL(R,Q)<\infty$, and $\lambda\in\Lam$. Then the positive part of
$\log\hl_\lambda(Y)$ is $R$-integrable, and
\begin{equation*}
  \E_R\log\hl_\lambda(Y)\le\KL(R,Q).
\end{equation*}
\end{lemma}

(See the proof in Subsection~\ref{sec:proof-lem-weak}.) Classical (Jensen's inequality, as in
\citealp[Lemma~7]{AJK21}).

\begin{lemma}[Duality for finitely supported \term{w:law}{laws}]\label{lem:finite}
\uses{def:normalised, def:dualset, lem:lambda, lem:weak}
Let $R=\sum_{i=1}^kp_i\delta_{y_i}$ with distinct $y_i$ and $p_i>0$, and $\Phi(\lambda):=\sum_ip_i\log\hl_\lambda(y_i)$.
Then $\Phi$ attains its maximum on $\Lam$, and
\begin{equation*}
  \KG(R)=\max_{\lambda\in\Lam}\Phi(\lambda).
\end{equation*}
Moreover the infimum defining $\KG(R)$ is attained by a \term{w:law}{law} of the form
$\sum_iw_i\delta_{y_i}+\sum_{u\in A}\beta_u\delta_u$, where $A$ is the set of zeros of $\hl_{\lambda^\star}$ for any maximiser
$\lambda^\star$ of $\Phi$, and has at most two points.
\end{lemma}

(See the proof in Subsection~\ref{sec:proof-lem-finite}.) Strong duality of this kind is classical
(\citealp[Theorem~8]{HT10}; \citealp[Lemma~7]{AJK21}). The proof here is the standard first-order
(Karush--Kuhn--Tucker) argument for a concave maximisation over a convex set: the optimal \term{w:law}{law} is read
off from a maximiser of $\Phi$ and the multipliers of its active constraints; it does not invoke Slater's condition.

\begin{lemma}[Duality for every \term{w:law}{law}]\label{lem:alllaws}
\uses{def:normalised, def:dualset, lem:lsc, lem:lambda, lem:weak, lem:finite}
Let $R$ be the \term{w:law}{law} of $Y$. If $\E_R\log(1+|Y|)=\infty$, then $\KG(R)=\infty=\sup_{\lambda\in\mathrm{int}\,\Lam}\E_R\log\hl_\lambda(Y)$.
Otherwise $\E_R\log\hl_\lambda(Y)\in[-\infty,\infty)$ is well defined for every $\lambda\in\Lam$, the map
$\lambda\mapsto\E_R\log\hl_\lambda(Y)$ is upper semicontinuous on $\Lam$, and
\begin{equation*}
  \KG(R)=\max_{\lambda\in\Lam}\E_R\log\hl_\lambda(Y)=\sup_{\lambda\in\mathrm{int}\,\Lam}\E_R\log\hl_\lambda(Y)<\infty.
\end{equation*}
\end{lemma}

(See the proof in Subsection~\ref{sec:proof-lem-alllaws}.) The scheme, duality first for finitely supported
\term{w:law}{laws} and then for every \term{w:law}{law} by approximation, is that of the proof of
\citet[Theorem~8]{HT10} for \term{w:law}{laws} on a bounded interval. On the unbounded support here, the limit is taken
with the lower semicontinuity of $\KG$ and Fatou's lemma, under the growth bound $C_\Lam(1+y^4)$ of
Lemma~\ref{lem:lambda}(d).

\begin{proof}[Proof of Proposition~\ref{prop:dual}]
\proves{prop:dual}
\uses{def:kinf, def:normalised, def:dualset, lem:G, lem:repr, lem:strict, lem:lsc, lem:lambda, lem:weak, lem:finite, lem:alllaws}
Item (i) is Lemma~\ref{lem:repr} together with Lemma~\ref{lem:strict}(a), which also gives the monotonicity and
the right-continuity. Item (ii) is Lemma~\ref{lem:G}(a), (c) and Lemma~\ref{lem:lsc}. Item (iii) is
Lemma~\ref{lem:lambda}(a)--(b). Item (iv) is Lemma~\ref{lem:alllaws}. For item (v), apply (i), then (iv) to
$R:=\tauc_{c,s}P$, the \term{w:law}{law} of $Y:=(X-c)/s$:
\begin{equation*}
  \Kc(P,\mu)=\inf_{c>\mu,\ s>0}\KG(\tauc_{c,s}P)=\inf_{c>\mu,\ s>0}\ \sup_{\lambda\in\mathrm{int}\,\Lam}\E_P\log\hl_\lambda\Bigl(\frac{X-c}{s}\Bigr),
\end{equation*}
which is \eqref{eq:dual}. Item (vi) is the last statement of Lemma~\ref{lem:finite}.
\end{proof}

\subsection{Proof of Lemma~\ref{lem:G}}\label{sec:proof-lem-G}

\begin{proof}
\proves{lem:G}
\uses{def:class, def:kinf, def:normalised, def:laws}
\emph{(a)} Let $Q\in\Gk$ and $v_4:=\int y^4\,dQ$. The constraint $\int\phik\,dQ\le0$ reads
\begin{equation*}
  v_4-2\kappa v_2+\kappa\le0.
\end{equation*}
By Jensen's inequality $v_4\ge v_2^2$, hence $v_2^2-2\kappa v_2+\kappa\le0$, which is the interval of the
statement; in particular $v_2\ge\kappa-\sqrt{\kappa^2-\kappa}>0$. Then $v_4\le2\kappa v_2-\kappa\le2\kappa(\kappa+\sqrt{\kappa^2-\kappa})-\kappa<4\kappa^2$.
Since $Q$ has mean $0$, its variance is $v_2>0$ and
\begin{equation*}
  \kurt(Q)=\frac{v_4}{v_2^2}\le\frac{2\kappa v_2-\kappa}{v_2^2}=\kappa-\frac{\kappa(v_2-1)^2}{v_2^2}\le\kappa .
\end{equation*}
So $Q\in\Hk$.

\emph{(b)} The map $x\mapsto c+sx$ with $s>0$ preserves the finiteness of the fourth moment, multiplies the
variance by $s^2$ and preserves the kurtosis. If $\tauc_{c,s}Q\in\Gk$, then $\tauc_{c,s}Q\in\Hk$ with mean $0$ by
(a), and $Q$, its image under $x\mapsto c+sx$, belongs to $\Hk$ with mean $c$. Conversely, let $Q\in\Hk$ with
$\mean(Q)=c$, and $s:=\sd(Q)$. Then $Y:=(X-c)/s$ with $X\sim Q$ has mean $0$, variance $1$ and $\E Y^4=\kurt(Q)$, so
\begin{equation*}
  \int\phik\,d(\tauc_{c,s}Q)=\kurt(Q)-2\kappa+\kappa=\kurt(Q)-\kappa\le0,
\end{equation*}
and $\tauc_{c,s}Q\in\Gk$.

\emph{(c)} The three conditions defining $\Gk$ are preserved by mixtures, so $\Gk$ is convex. By (a) and
Markov's inequality, $Q([-M,M]^c)\le4\kappa^2/M^4$ for every $Q\in\Gk$ and $M>0$, so $\Gk$ is tight. It is
closed: let $Q_n\in\Gk$ converge weakly to $Q$. Since $\phik$ is continuous and bounded below,
\begin{equation*}
  \int\phik\,dQ\le\liminf_n\int\phik\,dQ_n\le0 .
\end{equation*}
Since $2\kappa y^2\le y^4/2+2\kappa^2$, we have $\phik(y)\ge y^4/2-2\kappa^2$, so $\int y^4\,dQ\le2\int\phik\,dQ+4\kappa^2<\infty$.
Finally $\sup_n\int y^4\,dQ_n<4\kappa^2$ makes $y\mapsto y$ uniformly integrable along $(Q_n)$, so
$\int y\,dQ=\lim_n\int y\,dQ_n=0$. Hence $Q\in\Gk$, and Prokhorov's theorem shows that $\Gk$ is compact.

\emph{(d)} The map $Q\mapsto Q'$ (image under $x\mapsto c+sx$) is a bijection of $\Pspace$ that maps $\Hk$ onto
$\Hk$ (by the first sentence of (b)), with $\mean(Q')=c+s\,\mean(Q)$ and $\KL(P',Q')=\KL(P,Q)$, since $\KL$ is
invariant under a measurable bijection with measurable inverse. Hence
\begin{align*}
  \Kc(P',c+s\mu)&=\inf\{\KL(P',Q'):Q\in\Hk,\ c+s\,\mean(Q)>c+s\mu\}\\
  &=\inf\{\KL(P,Q):Q\in\Hk,\ \mean(Q)>\mu\}=\Kc(P,\mu).
\end{align*}
\end{proof}

\subsection{Proof of Lemma~\ref{lem:repr}}\label{sec:proof-lem-repr}

\begin{proof}
\proves{lem:repr}
\uses{def:kinf, def:normalised, lem:G}
We prove the first identity; the second is the same with strict inequalities on $\mean(Q)$ and $c$.

\emph{Lower bound.} Let $Q\in\Hk$ with $\mean(Q)\ge\mu$, and $c:=\mean(Q)$, $s:=\sd(Q)$. By Lemma~\ref{lem:G}(b),
$\tauc_{c,s}Q\in\Gk$, so
\begin{equation*}
  \KL(P,Q)=\KL(\tauc_{c,s}P,\tauc_{c,s}Q)\ge\KG(\tauc_{c,s}P)\ge\inf_{c'\ge\mu,\ s'>0}\KG(\tauc_{c',s'}P).
\end{equation*}
Taking the infimum over $Q$ gives $\Kge(P,\mu)\ge\inf_{c\ge\mu,s>0}\KG(\tauc_{c,s}P)$.

\emph{Upper bound.} Let $c\ge\mu$, $s>0$ and $Q'\in\Gk$, and let $Q$ be the image of $Q'$ under
$x\mapsto c+sx$. By Lemma~\ref{lem:G}(b), $Q\in\Hk$ with mean $c\ge\mu$, so
\begin{equation*}
  \Kge(P,\mu)\le\KL(P,Q)=\KL(\tauc_{c,s}P,Q').
\end{equation*}
Taking the infimum over $Q'\in\Gk$, then over $(c,s)$, gives the converse inequality.
\end{proof}

\subsection{Proof of Lemma~\ref{lem:strict}}\label{sec:proof-lem-strict}

We use the following fact here and in the proof of Lemma~\ref{lem:mixture}.

\begin{lemma}[Divergence from a mixture]\label{lem:mixkl}
\uses{def:laws}
Let $P,Q,Q'$ be \term{w:law}{laws} and $\theta\in(0,1]$ with $Q'\ge\theta Q$ as measures. Then
\begin{equation}\label{eq:mixkl}
  \KL(P,Q')\le\KL(P,Q)+\log\frac1\theta .
\end{equation}
\end{lemma}

\begin{proof}
\proves{lem:mixkl}
\uses{def:laws}
If $\KL(P,Q)=\infty$ there is nothing to prove. Otherwise $P\ll Q\ll Q'$ and $dQ/dQ'\le1/\theta$ $Q'$-almost surely, so
$dP/dQ'=(dP/dQ)(dQ/dQ')\le\theta^{-1}dP/dQ$ $P$-almost surely; integrating the logarithm against $P$ gives \eqref{eq:mixkl}.
\end{proof}

\begin{proof}
\proves{lem:strict}
\uses{def:kinf, def:class, def:laws, lem:mixkl}
\emph{Idea.} An alternative of mean exactly $\mu$ and kurtosis at most $\kappa$ is first mixed with two symmetric
atoms, which lowers the kurtosis strictly below $\kappa$, then with a small atom to the right, which lifts the
mean above $\mu$; both mixtures cost little in $\KL$.

\emph{Step 1: a strictly smaller kurtosis.} Let $Q\in\Hk$ with mean $m$, variance $\sigma^2>0$ and fourth central
moment $\mu_4=\kurt(Q)\sigma^4\le\kappa\sigma^4$, and $\eta\in(0,1)$. Let
\begin{equation*}
  Q_1:=(1-\eta)Q+\frac\eta2\bigl(\delta_{m-\sqrt\kappa\sigma}+\delta_{m+\sqrt\kappa\sigma}\bigr).
\end{equation*}
It has mean $m$, variance $\sigma^2(1+\eta(\kappa-1))$ and fourth central moment
$(1-\eta)\mu_4+\eta\kappa^2\sigma^4\le\kappa\sigma^4(1+\eta(\kappa-1))$, so
\begin{equation*}
  \kurt(Q_1)\le\frac{\kappa}{1+\eta(\kappa-1)}<\kappa .
\end{equation*}
Since $Q_1\ge(1-\eta)Q$, Lemma~\ref{lem:mixkl} gives $\KL(P,Q_1)\le\KL(P,Q)+\log\frac1{1-\eta}$.

\emph{Step 2: proof of (b).} Since $\Hlt\subset\Hk$, $\Kinf(P,\mu;\Hlt)\ge\Kc(P,\mu)$. Conversely, for $Q\in\Hk$ with
$\mean(Q)>\mu$, the \term{w:law}{law} $Q_1$ of Step 1 belongs to $\Hlt$ and has the same mean, so
$\Kinf(P,\mu;\Hlt)\le\KL(P,Q)+\log\frac1{1-\eta}$. Taking the infimum over $Q$, then $\eta\to0$, gives (b).

\emph{Step 3: proof of (a).} Clearly $\Kc\ge\Kge$. Let $Q\in\Hk$ with $\mean(Q)\ge\mu$ and $\KL(P,Q)<\infty$; if
$\mean(Q)>\mu$ there is nothing to do, so let $\mean(Q)=\mu$. Take $Q_1$ from Step 1 and, for $\eta'\in(0,1)$,
\begin{equation*}
  Q_2:=(1-\eta')Q_1+\eta'\delta_{\mu+\sigma}.
\end{equation*}
Its mean is $\mu+\eta'\sigma>\mu$. Its variance and fourth central moment are polynomials in $\eta'$, and its
variance is positive at $\eta'=0$, so $\kurt(Q_2)\to\kurt(Q_1)<\kappa$ as $\eta'\to0$; for $\eta'$ small,
$Q_2\in\Hk$. Since $Q_2\ge(1-\eta)(1-\eta')Q$, Lemma~\ref{lem:mixkl} gives
\begin{equation*}
  \Kc(P,\mu)\le\KL(P,Q_2)\le\KL(P,Q)+\log\frac1{(1-\eta)(1-\eta')}.
\end{equation*}
Letting $\eta,\eta'\to0$ and taking the infimum over $Q$ gives $\Kc(P,\mu)\le\Kge(P,\mu)$.

\emph{Step 4: monotonicity and right-continuity.} $\Kc(P,\cdot)$ is non-decreasing because the set of admissible
alternatives shrinks as $\mu$ grows. Since $\{Q:\mean(Q)>\mu\}=\bigcup_{\mu'>\mu}\{Q:\mean(Q)\ge\mu'\}$,
\begin{equation*}
  \Kc(P,\mu)=\inf_{\mu'>\mu}\Kge(P,\mu')=\inf_{\mu'>\mu}\Kc(P,\mu')=\lim_{\mu'\downarrow\mu}\Kc(P,\mu'),
\end{equation*}
using (a) for the second equality and the monotonicity for the third.
\end{proof}

\subsection{Proof of Lemma~\ref{lem:lsc}}\label{sec:proof-lem-lsc}

\begin{proof}
\proves{lem:lsc}
\uses{def:normalised, lem:G}
By the Donsker--Varadhan formula, $\KL$ is a supremum of functions that are jointly continuous for the weak
topology, hence jointly weakly lower semicontinuous. Let $R_n\to R$ weakly and $\ell:=\liminf_n\KG(R_n)$; assume
$\ell<\infty$. Along a subsequence realising $\ell$, pick $Q_n\in\Gk$ with $\KL(R_n,Q_n)\le\KG(R_n)+1/n$. By
Lemma~\ref{lem:G}(c), a further subsequence has $Q_n\to Q\in\Gk$, and
\begin{equation*}
  \KG(R)\le\KL(R,Q)\le\liminf_n\KL(R_n,Q_n)=\ell .
\end{equation*}
With $R_n:=R$ for every $n$, the subsequence $(Q_n)$ chosen in this proof converges to some $Q\in\Gk$ with
$\KL(R,Q)\le\KG(R)$, so $Q$ attains $\KG(R)$.
\end{proof}

\subsection{Proof of Lemma~\ref{lem:lambda}}\label{sec:proof-lem-lambda}

\begin{proof}
\proves{lem:lambda}
\uses{def:dualset, def:normalised}
Write $b:=\kappa(\kappa-1)$, so that $\min\phik=-b$, attained at $\pm\sqrt\kappa$.

\emph{(a)} $\Lam$ is the intersection of the closed half-plane $\{\lambda_2\ge0\}$ with the closed half-planes
$\{\lambda:\hl_\lambda(y)\ge0\}$, $y\in\R$; it is closed and convex. At $y=\pm\sqrt\kappa$, $\hl_\lambda(\pm\sqrt\kappa)=1\pm\lambda_1\sqrt\kappa-\lambda_2b\ge0$,
hence $\lambda_2\le1/b$ and $|\lambda_1|\sqrt\kappa\le1-\lambda_2b\le1$.

\emph{(b)} Let $0\le\lambda_2\le1/b$. At $y=0$, $\hl_\lambda(0)=1+\lambda_2\kappa>0$. For $y>0$, $\hl_\lambda(y)\ge0$ means
$\lambda_1\ge-(1+\lambda_2\phik(y))/y$, and $\hl_\lambda(-y)\ge0$ means $\lambda_1\le(1+\lambda_2\phik(y))/y$, since $\phik$ is even.
So $\lambda\in\Lam$ if and only if $|\lambda_1|\le\Lw(\lambda_2)$ (and $\Lw(\lambda_2)\ge0$, which holds because
$1+\lambda_2\phik\ge1-\lambda_2b\ge0$). For $\lambda_2>1/b$, $1+\lambda_2\phik(\sqrt\kappa)<0$, so no $\lambda_1$ is admissible. The
function $\Lw$ is an infimum of affine functions of $\lambda_2$, hence concave. For $0<\lambda_2<1/b$, the function
$y\mapsto(1+\lambda_2\phik(y))/y$ is continuous and positive on $(0,\infty)$ (its numerator is at least $1-\lambda_2b>0$)
and tends to $+\infty$ at $0$ and at $\infty$, so its infimum $\Lw(\lambda_2)$ is positive. A concave function is
continuous on an open interval, so $\{(\lambda_1,\lambda_2):0<\lambda_2<1/b,\ |\lambda_1|<\Lw(\lambda_2)\}$ is a non-empty open subset of
$\Lam$.

\emph{(c)} Let $\lambda\in\mathrm{int}\,\Lam$. Points with $\lambda_2=0$ are not interior, since $\Lam\subset\{\lambda_2\ge0\}$; so
$\lambda_2>0$. Suppose $\hl_\lambda(y_0)=0$ for some $y_0$. Since $\gam(y_0)\neq0$ (if $y_0=0$, $\phik(0)=\kappa\neq0$),
$\lambda-t\gam(y_0)\in\Lam$ for $t>0$ small, while
\begin{equation*}
  \hl_{\lambda-t\gam(y_0)}(y_0)=\hl_\lambda(y_0)-t|\gam(y_0)|^2<0,
\end{equation*}
a contradiction. So $\hl_\lambda>0$ on $\R$; since $\lambda_2>0$, $\hl_\lambda(y)\to+\infty$ as $|y|\to\infty$, so the continuous
function $\hl_\lambda$ attains a positive minimum.

\emph{(d)} For $\lambda\in\Lam$, $|\lambda_1|\le1$ and $0\le\lambda_2\le1/b$, and $\phik(y)\le y^4+\kappa$, so
$\hl_\lambda(y)\le1+|y|+(y^4+\kappa)/b\le(2+(1+\kappa)/b)(1+y^4)$.
\end{proof}

\subsection{Proof of Lemma~\ref{lem:weak}}\label{sec:proof-lem-weak}

\begin{proof}
\proves{lem:weak}
\uses{def:normalised, def:dualset}
\emph{Setting.} $R\ll Q$ since $\KL(R,Q)<\infty$; let $\rho:=dR/dQ$, so $\rho(Y)>0$ $R$-almost surely. \emph{Goal.} Split
$\log\hl_\lambda(Y)$ into a part with integrable positive part and an integrable part, and apply Jensen's inequality
to the first.

Let $Z:=\hl_\lambda(Y)/\rho(Y)\ge0$, defined $R$-almost surely. Since $\hl_\lambda\ge0$, $Q$ has mean $0$ and a finite fourth
moment, and $\lambda_2\ge0$,
\begin{equation*}
  \E_R[Z]=\int_{\{\rho>0\}}\hl_\lambda\,dQ\le\int\hl_\lambda\,dQ=1+\lambda_2\int\phik\,dQ\le1 .
\end{equation*}
So $\E_R[(\log Z)^+]\le\E_R[Z]\le1$. The function $\log\rho(Y)$ is $R$-integrable: its integral is $\KL(R,Q)<\infty$
and its negative part has integral $\int_{\{\rho<1\}}\rho\log(1/\rho)\,dQ\le1/e$. Since
$\log\hl_\lambda(Y)=\log Z+\log\rho(Y)$ $R$-almost surely, the positive part of $\log\hl_\lambda(Y)$ is integrable and
\begin{align*}
  \E_R\log\hl_\lambda(Y)&=\E_R\log Z+\E_R\log\rho(Y)\\
  &\le\log\E_R[Z]+\KL(R,Q) && \text{(Jensen's inequality for the concave }\log\text{)}\\
  &\le\KL(R,Q) && (\E_R[Z]\le1).
\end{align*}
\end{proof}

\subsection{Proof of Lemma~\ref{lem:finite}}\label{sec:proof-lem-finite}

\begin{proof}
\proves{lem:finite}
\uses{def:normalised, def:dualset, lem:lambda, lem:weak}
\emph{Idea.} Weak duality gives $\max\Phi\le\KG(R)$. For the converse we take a maximiser $\lambda^\star$ of $\Phi$ and
write the first-order optimality condition: the gradient of $\Phi$ at $\lambda^\star$ is a non-positive combination
of the vectors $\gam(u)$ at the zeros $u$ of $\hl_{\lambda^\star}$. The coefficients of this combination, with the weights
$p_i/\hl_{\lambda^\star}(y_i)$, define a \term{w:law}{law} of $\Gk$ whose divergence from $R$ is $\Phi(\lambda^\star)$.

\emph{Step 1: a maximiser.} Each map $\lambda\mapsto\log\hl_\lambda(y_i)$ is upper semicontinuous with values in $[-\infty,\infty)$
($\hl_\lambda(y_i)$ is continuous in $\lambda$ and $\log0:=-\infty$), so $\Phi$ attains its maximum on the compact set
$\Lam$ (Lemma~\ref{lem:lambda}(a)). Let $\lambda^\star$ be any maximiser. Since $0\in\Lam$ and $\Phi(0)=0$, $\Phi(\lambda^\star)\ge0$; hence
$\hl_{\lambda^\star}(y_i)>0$ for every $i$, $\Phi$ is differentiable at $\lambda^\star$, and its gradient is
\begin{equation*}
  g:=\nabla\Phi(\lambda^\star)=\sum_{i=1}^k\frac{p_i\,\gam(y_i)}{\hl_{\lambda^\star}(y_i)} .
\end{equation*}
We use repeatedly that if $\lambda^\star+td\in\Lam$ for all small $t>0$, then $g\cdot d\le0$, because
$\Phi(\lambda^\star+td)\le\Phi(\lambda^\star)$ and the right derivative at $t=0$ is $g\cdot d$.

\emph{Step 2: the case $\lambda^\star_2=0$.} Then $1+\lambda^\star_1y\ge0$ for all $y$, so $\lambda^\star=0$ and $g=\sum_ip_i\gam(y_i)$.
For $d\in\R^2$ with $d_2>0$, the quartic $d_1y+d_2\phik(y)$ is bounded below, so $td\in\Lam$ for small $t>0$, and
$g\cdot d\le0$. With $d=(\pm1,\zeta)$, $\zeta\downarrow0$, and $d=(0,1)$ we get
\begin{equation*}
  \sum_ip_iy_i=0,\qquad \sum_ip_i\phik(y_i)\le0 .
\end{equation*}
So $R\in\Gk$, hence $\KG(R)=\KL(R,R)=0=\Phi(0)=\Phi(\lambda^\star)$, and $R$ itself attains $\KG(R)$ (here $A=\emptyset$).

\emph{Step 3: the case $\lambda^\star_2>0$; the zeros.} Then $\hl_{\lambda^\star}$ is a non-negative polynomial of degree $4$
with positive leading coefficient, so its real zeros have even multiplicity: the set $A$ of its zeros has at
most two points, none of them equal to a $y_i$. Let $C:=\{\sum_{u\in A}\beta_u\gam(u):\beta_u\ge0\}$, a finitely generated, hence
closed, convex cone ($C=\{0\}$ if $A=\emptyset$).

\emph{Step 4: the first-order condition, $-g\in C$.} Suppose not. Separating $-g$ from the closed convex cone $C$
gives $d\in\R^2$ with $d\cdot g>0$ and $d\cdot\gam(u)\ge0$ for every $u\in A$. For $\zeta>0$ let $d_\zeta:=d-\zeta\lambda^\star$.
For $u\in A$, $\lambda^\star\cdot\gam(u)=\hl_{\lambda^\star}(u)-1=-1$, so
\begin{equation*}
  d_\zeta\cdot\gam(u)=d\cdot\gam(u)+\zeta\ge\zeta>0,
\end{equation*}
and $d_\zeta\cdot g>0$ for $\zeta$ small; fix such a $\zeta$. We show that $\lambda^\star+td_\zeta\in\Lam$ for small $t>0$, which
contradicts Step 1. Let $\psi(y):=d_\zeta\cdot\gam(y)$, so that $\hl_{\lambda^\star+td_\zeta}=\hl_{\lambda^\star}+t\psi$, and let
$U:=\{\psi>0\}$, an open set containing $A$.
\begin{itemize}
  \item On $U$: $\hl_{\lambda^\star+td_\zeta}\ge t\psi>0$.
  \item On $F:=\R\setminus U$: $\hl_{\lambda^\star}>0$ on the closed set $F$ and $\hl_{\lambda^\star}(y)/(1+y^4)\to\lambda^\star_2>0$ as $|y|\to\infty$, so
  $c_F:=\inf_F\hl_{\lambda^\star}(y)/(1+y^4)>0$; and $|\psi(y)|\le C_\psi(1+y^4)$ for some $C_\psi<\infty$. Hence, for $t\le c_F/C_\psi$,
  $\hl_{\lambda^\star+td_\zeta}\ge(c_F-tC_\psi)(1+y^4)\ge0$ on $F$.
  \item The second coordinate $\lambda^\star_2+t(d_\zeta)_2$ is non-negative for $t$ small.
\end{itemize}
So $\lambda^\star+td_\zeta\in\Lam$ for small $t>0$, and $g\cdot d_\zeta\le0$ by Step 1: a contradiction. Hence
$-g=\sum_{u\in A}\beta_u\gam(u)$ with $\beta_u\ge0$.

\emph{Step 5: an optimal \term{w:law}{law}.} Let $w_i:=p_i/\hl_{\lambda^\star}(y_i)>0$ and $Q^\star:=\sum_iw_i\delta_{y_i}+\sum_{u\in A}\beta_u\delta_u$. By
definition of $g$ and Step 4,
\begin{equation*}
  \sum_iw_i\gam(y_i)=g=-\sum_{u\in A}\beta_u\gam(u).
\end{equation*}
\begin{itemize}
  \item Total mass: $\sum_iw_i\hl_{\lambda^\star}(y_i)=\sum_ip_i=1$, that is $\sum_iw_i+\lambda^\star\cdot g=1$; and
  $\lambda^\star\cdot g=-\sum_u\beta_u\lambda^\star\cdot\gam(u)=\sum_u\beta_u$. So $\sum_iw_i+\sum_u\beta_u=1$ and $Q^\star$ is a \term{w:law}{law}.
  \item Constraints: the first coordinate of the display gives $\int y\,dQ^\star=0$, the second
  $\int\phik\,dQ^\star=0$; $Q^\star$ has finite support. So $Q^\star\in\Gk$.
  \item Divergence: the $u\in A$ differ from the $y_i$, so $dR/dQ^\star(y_i)=p_i/w_i=\hl_{\lambda^\star}(y_i)$ and
  $\KL(R,Q^\star)=\sum_ip_i\log\hl_{\lambda^\star}(y_i)=\Phi(\lambda^\star)$.
\end{itemize}
Hence $\KG(R)\le\Phi(\lambda^\star)$. With Lemma~\ref{lem:weak}, $\KG(R)=\Phi(\lambda^\star)=\max_\Lam\Phi$, attained by $Q^\star$.
\end{proof}

\subsection{Proof of Lemma~\ref{lem:alllaws}}\label{sec:proof-lem-alllaws}

\begin{proof}
\proves{lem:alllaws}
\uses{def:normalised, def:dualset, lem:lsc, lem:lambda, lem:weak, lem:finite}
\emph{Idea.} Approximate $R$ by its empirical \term{w:law}{laws}, for which Lemma~\ref{lem:finite} holds, take limits of the
dual maximisers, and pass to the limit in the dual objective with Fatou's lemma, dominated by the growth
bound $C_\Lam(1+y^4)$.

\emph{Step 1: the case $\E_R\log(1+|Y|)=\infty$.} Let $b:=\kappa(\kappa-1)$ and $\lambda^0:=(0,\frac1{2b})$, an interior point of $\Lam$
(Lemma~\ref{lem:lambda}(b)). Then
\begin{equation*}
  \hl_{\lambda^0}(y)=\frac12+\frac{(y^2-\kappa)^2}{2b}\ \ge\ \max\Bigl\{\frac12,\ \frac{(y^2-\kappa)^2}{2b}\Bigr\},
\end{equation*}
so $\log\hl_{\lambda^0}(y)\ge4\log(1+|y|)-c_\kappa$ for a constant $c_\kappa$ (use the first bound for $|y|\le2\sqrt\kappa$, and
$y^2-\kappa\ge\frac34y^2\ge\frac3{16}(1+|y|)^2$ for $|y|\ge2\sqrt\kappa$); hence $\E_R\log\hl_{\lambda^0}(Y)=\infty$. By
Lemma~\ref{lem:weak}, $\KL(R,Q)=\infty$ for every $Q\in\Gk$, so $\KG(R)=\infty$.

\emph{Step 2: well-definedness and upper semicontinuity.} Assume $\E_R\log(1+|Y|)<\infty$. By
Lemma~\ref{lem:lambda}(d), $(\log\hl_\lambda(Y))^+\le\log C_\Lam+\log(1+Y^4)$, which is $R$-integrable; so
$\E_R\log\hl_\lambda(Y)\in[-\infty,\infty)$ is defined for every $\lambda\in\Lam$. Let $\lambda_n\to\lambda$ in $\Lam$, $f_n:=\log\hl_{\lambda_n}(Y)$ and
$G:=\log C_\Lam+\log(1+Y^4)$. Then $G-f_n\ge0$ and $\limsup_nf_n\le\log\hl_\lambda(Y)$ pointwise, by continuity of
$(\lambda,y)\mapsto\hl_\lambda(y)$ and upper semicontinuity of $\log$ on $[0,\infty)$. Fatou's lemma applied to $G-f_n$ gives
\begin{equation*}
  \limsup_n\E_Rf_n\le\E_R\log\hl_\lambda(Y).
\end{equation*}
So $\lambda\mapsto\E_R\log\hl_\lambda(Y)$ is upper semicontinuous, and it attains its maximum on the compact set $\Lam$.

\emph{Step 3: $\KG(R)\le\max_\Lam\E_R\log\hl_\lambda$.} Let $Y_1,Y_2,\dots$ be independent with \term{w:law}{law} $R$ and $R_n$ their empirical
\term{w:law}{laws}. Almost surely, $R_n\to R$ weakly (Varadarajan's theorem) and $\frac1n\sum_{i\le n}\log(1+Y_i^4)\to\E_R\log(1+Y^4)$ (strong
\term{w:law}{law} of large numbers); fix an outcome where both hold, so that $(R_n)$ is a deterministic sequence of finitely
supported \term{w:law}{laws}. By Lemma~\ref{lem:finite}, $\KG(R_n)=\E_{R_n}\log\hl_{\lambda_n}$ for some $\lambda_n\in\Lam$. Along a subsequence realising
$\liminf_n\KG(R_n)$, we may assume $\lambda_n\to\bar\lambda\in\Lam$. By Skorokhod's representation theorem, there are random
variables $Z_n\sim R_n$ and $Z\sim R$ with $Z_n\to Z$ almost surely. Let $f_n:=\log\hl_{\lambda_n}(Z_n)$,
$G_n:=\log C_\Lam+\log(1+Z_n^4)\ge f_n$ and $G:=\log C_\Lam+\log(1+Z^4)$. Then $G_n\to G$ almost surely, $\E G_n\to\E G<\infty$
(by the choice of the outcome), and $\limsup_nf_n\le\log\hl_{\bar\lambda}(Z)$ almost surely. Fatou's lemma applied to
$G_n-f_n\ge0$ gives
\begin{equation*}
  \E G-\E\log\hl_{\bar\lambda}(Z)\le\E\liminf_n(G_n-f_n)\le\liminf_n\E(G_n-f_n)=\E G-\limsup_n\E f_n .
\end{equation*}
So $\limsup_n\E f_n\le\E_R\log\hl_{\bar\lambda}(Y)$. With Lemma~\ref{lem:lsc},
\begin{equation*}
  \KG(R)\le\liminf_n\KG(R_n)=\lim_n\E f_n\le\E_R\log\hl_{\bar\lambda}(Y)\le\max_{\lambda\in\Lam}\E_R\log\hl_\lambda(Y).
\end{equation*}

\emph{Step 4: conclusion.} Lemma~\ref{lem:weak} gives the converse inequality, so $\KG(R)$ equals the maximum, which
is at most $\log C_\Lam+\E_R\log(1+Y^4)<\infty$. Finally, let $\lambda^\star$ be a maximiser and $\lambda^0\in\mathrm{int}\,\Lam$; for
$t\in(0,1]$, $\lambda_t:=(1-t)\lambda^\star+t\lambda^0\in\mathrm{int}\,\Lam$, and since $\hl$ is affine in $\lambda$ and $\log$ is concave,
\begin{equation*}
  \E_R\log\hl_{\lambda_t}(Y)\ge(1-t)\E_R\log\hl_{\lambda^\star}(Y)+t\,\E_R\log\hl_{\lambda^0}(Y),
\end{equation*}
where $\E_R\log\hl_{\lambda^0}(Y)$ is finite by Lemma~\ref{lem:lambda}(c)--(d). Letting $t\to0$ shows that the supremum over
$\mathrm{int}\,\Lam$ equals the maximum over $\Lam$.
\end{proof}

% ----------------------------------------------------------------
\section{\texorpdfstring{Proof of Proposition~\ref{prop:kinfty}}{Proof of Proposition 4.3}}\label{sec:proof-kinfty}

To prove Proposition~\ref{prop:kinfty} we use the following three lemmas.

\begin{lemma}[The largest atom]\label{lem:atom}
\uses{def:class, def:kinfty}
Let $\kappa>1$. For every $Q\in\Hk$ and $y\in\R$, $Q(\{y\})\le\qk$. For $y'\neq y$, the \term{w:law}{law}
$\qk\delta_y+(1-\qk)\delta_{y'}$ has kurtosis $\kappa$.
\end{lemma}

(See the proof in Subsection~\ref{sec:proof-lem-atom}.) An elementary extremal moment computation, an instance of
the classical question of the largest mass a \term{w:law}{law} can put on a point under moment constraints (the theory of
Chebyshev, Markov and Krein); we did not find this statement in the literature.

\begin{lemma}[Mixing with a far atom]\label{lem:mixture}
\uses{def:kinf, def:kinfty}
Let $\kappa>1$, $P$ a \term{w:law}{law} with $\int x^4\,dP<\infty$ and $\mu\in\R$. Then $\Kc(P,\mu)\le\Kinfty(\kappa)$.
\end{lemma}

(See the proof in Subsection~\ref{sec:proof-lem-mixture}.) The ``outlier'' perturbation is the idea of
\citet[Appendix~A]{L17}; the sharp mass $1-\qk$ is new.

\begin{lemma}[Size of the value at infinity]\label{lem:expansion}
\uses{def:kinfty}
For every $\kappa>1$, $\kappa<1/\Kinfty(\kappa)<\kappa+2$, and $1/\Kinfty(\kappa)=\kappa+\frac32+O(1/\kappa)$ as $\kappa\to\infty$, with an
absolute constant in the $O(\cdot)$.
\end{lemma}

(See the proof in Subsection~\ref{sec:proof-lem-expansion}.) Elementary.

\begin{proof}[Proof of Proposition~\ref{prop:kinfty}]
\proves{prop:kinfty}
\uses{def:kinf, def:kinfty, def:class, lem:atom, lem:mixture, lem:expansion, lem:G, lem:repr}
Items (i), (ii) and (iv) are Lemmas~\ref{lem:atom}, \ref{lem:mixture} and \ref{lem:expansion}. For (iii), let
$y,\mu\in\R$. For every $z\in\R$ and $Q\in\Gk$, $\KL(\delta_z,Q)=\log(1/Q(\{z\}))$ (or $+\infty$ if $Q(\{z\})=0$), and
$Q\in\Hk$ by Lemma~\ref{lem:G}(a), so Lemma~\ref{lem:atom} gives
\begin{equation*}
  \KG(\delta_z)\ge\log\frac1{\qk}=\Kinfty(\kappa).
\end{equation*}
Since $\tauc_{c,s}\delta_y=\delta_{(y-c)/s}$, Lemma~\ref{lem:repr} gives $\Kc(\delta_y,\mu)\ge\Kinfty(\kappa)$, and
Lemma~\ref{lem:mixture}, applied to $P=\delta_y$, gives the converse inequality.
\end{proof}

\subsection{Proof of Lemma~\ref{lem:atom}}\label{sec:proof-lem-atom}

\begin{proof}
\proves{lem:atom}
\uses{def:class, def:kinfty, def:laws}
\emph{Idea.} Conditionally on not being at $y$, the \term{w:law}{law} has a mean forced by the atom; Jensen's inequality
bounds the kurtosis from below by that of the two-point \term{w:law}{law} with the same atom, which is
$1/(q(1-q))-3$.

\emph{Setting.} Let $q:=Q(\{y\})$. If $q\le\frac12$, then $q<\qk$; so assume $q>\frac12$. Since $\sd^2(Q)>0$, $q<1$.
Let $m:=\mean(Q)$, $d:=y-m$, and $S$ the \term{w:law}{law} $Q$ conditioned on $\R\setminus\{y\}$. Since
$qd+(1-q)\int(x-m)\,dS=0$, the \term{w:law}{law} $S$ has mean $m+e$ with
\begin{equation*}
  e:=-\frac{qd}{1-q},
\end{equation*}
and $u:=\int(x-m)^2\,dS\ge e^2$. \emph{Goal.} Show $q(1-q)\ge1/(\kappa+3)$, which with $q>\frac12$ gives $q\le\qk$.

\emph{Step 1: two moments.} We have
\begin{equation*}
  \sd^2(Q)=qd^2+(1-q)u,\qquad \int(x-m)^4\,dQ=qd^4+(1-q)\int(x-m)^4\,dS\ge qd^4+(1-q)u^2,
\end{equation*}
the last inequality by Jensen's inequality.

\emph{Step 2: the case $d\neq0$.} By Step 1, $\kurt(Q)\ge f(u)$ with
\begin{equation*}
  f(v):=\frac{qd^4+(1-q)v^2}{(qd^2+(1-q)v)^2},\qquad f'(v)=\frac{2q(1-q)d^2(v-d^2)}{(qd^2+(1-q)v)^3},
\end{equation*}
so $f$ is non-decreasing on $[d^2,\infty)$. Since $q>\frac12$, $e^2=q^2d^2/(1-q)^2>d^2$, hence $u\ge e^2>d^2$ and,
using $qd^2+(1-q)e^2=qd^2/(1-q)$ and $q^3+(1-q)^3=1-3q(1-q)$,
\begin{equation*}
  \kappa\ge\kurt(Q)\ge f(e^2)=\frac{(1-q)^3+q^3}{q(1-q)}=\frac1{q(1-q)}-3 .
\end{equation*}
So $q(1-q)\ge1/(\kappa+3)$.

\emph{Step 3: the case $d=0$.} Then $\sd^2(Q)=(1-q)u$ with $u>0$, and $\int(x-m)^4\,dQ\ge(1-q)u^2$, so
$\kappa\ge\kurt(Q)\ge1/(1-q)$, that is $q\le1-1/\kappa$. If $\kappa\le2$ this contradicts $q>\frac12$. If $\kappa>2$, then
$t:=1-1/\kappa\in(\frac12,1)$ satisfies
\begin{equation*}
  t(1-t)=\frac{\kappa-1}{\kappa^2}\ge\frac1{\kappa+3}\quad\text{(equivalent to }2\kappa\ge3\text{)},
\end{equation*}
so $t\le\qk$ and $q\le\qk$.

\emph{Step 4: the two-point \term{w:law}{law}.} A \term{w:law}{law} with masses $q$ and $1-q$ at two points at distance $D$ has variance
$q(1-q)D^2$ and fourth central moment $q(1-q)D^4(q^3+(1-q)^3)$, hence kurtosis $1/(q(1-q))-3$, which equals
$\kappa$ for $q=\qk$.
\end{proof}

\subsection{Proof of Lemma~\ref{lem:mixture}}\label{sec:proof-lem-mixture}

\begin{proof}
\proves{lem:mixture}
\uses{def:kinf, def:kinfty, def:class, def:laws, lem:mixkl}
\emph{Setting.} Let $m:=\mean(P)$, $\sigma^2:=\sd^2(P)\ge0$, $\mu_3$ and $\mu_4$ the third and fourth central moments of $P$
(finite), $p\in(1-\qk,\frac12)$, and for $x\in\R$
\begin{equation*}
  Q_x:=(1-p)P+p\,\delta_x,\qquad D:=x-m .
\end{equation*}
\emph{Goal.} Show that $Q_x\in\Hk$ and $\mean(Q_x)>\mu$ for $x$ large, and $\KL(P,Q_x)\le\log\frac1{1-p}$.

\emph{Step 1: moments of $Q_x$.} The mean is $\mean(Q_x)=m+pD$, and the variance is
$\sd^2(Q_x)=(1-p)\sigma^2+p(1-p)D^2$. Since $\E_P(X-m)=0$, the fourth central moment is
\begin{align*}
  \int(y-m-pD)^4\,dQ_x&=(1-p)\E_P(X-m-pD)^4+p(1-p)^4D^4\\
  &=(1-p)\bigl(\mu_4-4pD\mu_3+6p^2D^2\sigma^2+p^4D^4\bigr)+p(1-p)^4D^4\\
  &=p(1-p)\bigl(1-3p(1-p)\bigr)D^4+O(D^2)\qquad(D\to\infty),
\end{align*}
using $(1-p)p^4+p(1-p)^4=p(1-p)(p^3+(1-p)^3)=p(1-p)(1-3p(1-p))$.

\emph{Step 2: the kurtosis.} Since $\sd^4(Q_x)=p^2(1-p)^2D^4+O(D^2)$,
\begin{equation*}
  \kurt(Q_x)\longrightarrow\frac{1-3p(1-p)}{p(1-p)}=\frac1{p(1-p)}-3\qquad(x\to\infty).
\end{equation*}
The map $t\mapsto t(1-t)$ is increasing on $(0,\frac12)$ and $p>1-\qk$, so $p(1-p)>\qk(1-\qk)=1/(\kappa+3)$ and the limit
is less than $\kappa$. Hence, for $x$ large, $Q_x$ has a finite fourth moment, a positive variance and kurtosis
less than $\kappa$, so $Q_x\in\Hk$, and $\mean(Q_x)=m+p(x-m)>\mu$.

\emph{Step 3: the divergence.} Since $Q_x\ge(1-p)P$, Lemma~\ref{lem:mixkl} gives $\KL(P,Q_x)\le\log\frac1{1-p}$. Hence
$\Kc(P,\mu)\le\log\frac1{1-p}$ for every $p\in(1-\qk,\frac12)$, and letting $p\downarrow1-\qk$ gives $\Kc(P,\mu)\le\log\frac1{\qk}$.
\end{proof}

\subsection{Proof of Lemma~\ref{lem:expansion}}\label{sec:proof-lem-expansion}

\begin{proof}
\proves{lem:expansion}
\uses{def:kinfty}
Let $p:=1-\qk\in(0,\frac12)$, so that $p(1-p)=1/(\kappa+3)$ and $\Kinfty(\kappa)=-\log(1-p)$.

\emph{Bounds.} For $p\in(0,1)$, $p<-\log(1-p)<p/(1-p)$, so
\begin{equation*}
  \frac{1-p}p<\frac1{\Kinfty(\kappa)}<\frac1p .
\end{equation*}
Now $(\kappa+3)p=1/(1-p)$, so
\begin{align*}
  \frac1p&=\frac{1-p}{p(1-p)}=(\kappa+3)(1-p)=\kappa+3-\frac1{1-p}<\kappa+2,\\
  \frac{1-p}p&=(\kappa+3)(1-p)^2=(\kappa+3)(1-p)-(\kappa+3)p(1-p)=\kappa+2-\frac1{1-p}>\kappa,
\end{align*}
since $1<1/(1-p)<2$ for $p\in(0,\frac12)$.

\emph{Limit.} As $\kappa\to\infty$, $p=(\kappa+3)^{-1}/(1-p)\le2/(\kappa+3)\to0$. By Taylor's formula,
$-\log(1-p)=p+\frac{p^2}2+O(p^3)$, so
\begin{equation*}
  \frac1{\Kinfty(\kappa)}=\frac1p\Bigl(1-\frac p2+O(p^2)\Bigr)=\frac1p-\frac12+O(p)
  =\kappa+2-\frac p{1-p}-\frac12+O(p)=\kappa+\frac32+O\Bigl(\frac1\kappa\Bigr).
\end{equation*}
\end{proof}

% ----------------------------------------------------------------
\section{\texorpdfstring{Proof of Lemma~\ref{lem:escape}, Proposition~\ref{prop:regularity} and Corollary~\ref{cor:largegap}}{Proof of Lemma 8.1, Proposition 4.4 and Corollary 4.5}}\label{sec:proof-regularity}

This appendix carries Step~4 of Section~\ref{sec:overview}. We first prove the escape lemma, which uses one
more elementary lemma, then derive the regularity of $\Kc$ from it.

\begin{lemma}[Mass far away is expensive]\label{lem:tail}
\uses{def:normalised, lem:G}
For every \term{w:law}{law} $R$ and every Borel set $A$,
\begin{equation*}
  \KG(R)\ge R(A)\log\frac1{\sup_{Q\in\Gk}Q(A)}-\log2,\qquad\text{and}\qquad\sup_{Q\in\Gk}Q([-M,M]^c)\le\frac{4\kappa^2}{M^4}\quad(M>0).
\end{equation*}
\end{lemma}

(See the proof in Subsection~\ref{sec:proof-lem-tail}.) Classical (data processing on a two-set partition).

\subsection{Proof of Lemma~\ref{lem:tail}}\label{sec:proof-lem-tail}

\begin{proof}
\proves{lem:tail}
\uses{def:normalised, lem:G}
Let $Q\in\Gk$, $a:=R(A)$ and $b:=Q(A)$. By the data processing inequality applied to the partition $\{A,A^c\}$,
\begin{equation*}
  \KL(R,Q)\ge\kl(a,b),\qquad\text{where}\quad\dt{s:kl}\kl(a,b):=a\log\frac ab+(1-a)\log\frac{1-a}{1-b}
\end{equation*}
is the Bernoulli divergence. Since $a\log a+(1-a)\log(1-a)\ge-\log2$ and $\log\frac1{1-b}\ge0$,
\begin{equation*}
  \kl(a,b)\ge a\log\frac1b-\log2 .
\end{equation*} Taking $b\le\sup_{\Gk}Q(A)$ and the infimum over $Q$
gives the first claim. The second is Markov's inequality with $\int y^4\,dQ<4\kappa^2$ (Lemma~\ref{lem:G}(a)).
\end{proof}

\subsection{Proof of Lemma~\ref{lem:escape}}\label{sec:proof-lem-escape}

\begin{proof}
\proves{lem:escape}
\uses{def:normalised, def:kinfty, lem:G, lem:lsc, lem:tail, lem:atom}
\emph{Idea.} Extract a subsequence along which the location, the scale and their ratio converge in the
extended line. Either the affine maps converge to a genuine affine map, or the normalised \term{w:law}{laws} spread
their mass to infinity, which Lemma~\ref{lem:tail} makes infinitely expensive, or they collapse to a Dirac mass,
which costs at least $\Kinfty(\kappa)$ by the largest atom.

\emph{Setting.} Let $c_-:=\inf_nc_n>-\infty$. Extract a subsequence along which $c_n\to c\in\R\cup\{+\infty\}$,
$s_n\to s\in[0,\infty]$ and $c_n/s_n\to r\in[-\infty,+\infty]$. Since $(P_n)$ converges weakly, it is tight (Prokhorov's
theorem). \emph{Goal.} Show that each of the four cases below gives (a) or (b), and that they exhaust all
possibilities.

\emph{Case 1: $c<\infty$ and $0<s<\infty$.} If $X_n\sim P_n$ and $X\sim P$ with $X_n\to X$ in \term{w:law}{law}, then
$(X_n-c_n)/s_n\to(X-c)/s$ in \term{w:law}{law} by Slutsky's lemma, so $R_n\to\tauc_{c,s}P$ weakly, and Lemma~\ref{lem:lsc} gives
$\liminf_n\KG(R_n)\ge\KG(\tauc_{c,s}P)$: this is (a).

\emph{Case 2: $c<\infty$ and $s=0$.} Let $\rho:=(1-P(\{c\}))/2$, positive because $P$ is not a Dirac mass. Since
$P([c-\eta,c+\eta])\downarrow P(\{c\})$ as $\eta\downarrow0$, fix $\eta>0$ with $P([c-\eta,c+\eta])<1-\rho$. Fix $M>0$. For $n$ large,
$|c_n-c|+Ms_n\le\eta$, so
\begin{equation*}
  R_n([-M,M])=P_n([c_n-Ms_n,c_n+Ms_n])\le P_n([c-\eta,c+\eta]),
\end{equation*}
whose $\limsup$ is at most $P([c-\eta,c+\eta])<1-\rho$ by the portmanteau theorem. So $R_n([-M,M]^c)\ge\rho$ for $n$
large, and Lemma~\ref{lem:tail} gives
\begin{equation*}
  \liminf_n\KG(R_n)\ge\rho\log\frac{M^4}{4\kappa^2}-\log2 .
\end{equation*}
Since $M$ is arbitrary, $\liminf_n\KG(R_n)=\infty$: this is (b).

\emph{Case 3: $c=\infty$ and $r=\infty$.} By tightness, fix $M'$ with $P_n([-M',M'])\ge\frac12$ for every $n$. For $n$ large,
$c_n\ge2M'$, and then for $x\in[-M',M']$,
\begin{equation*}
  \frac{x-c_n}{s_n}\le\frac{M'-c_n}{s_n}\le-\frac{c_n}{2s_n}\longrightarrow-\infty .
\end{equation*}
So for every $M$, $R_n((-\infty,-M])\ge\frac12$ for $n$ large, and Lemma~\ref{lem:tail} gives $\KG(R_n)\to\infty$: this is (b).

\emph{Case 4: $s=\infty$ and $r<\infty$.} Since $c_n/s_n\ge c_-/s_n\to0$, $r\ge0$. If $X_n\sim P_n$, tightness gives
$X_n/s_n\to0$ in probability, so $R_n$, the \term{w:law}{law} of $X_n/s_n-c_n/s_n$, converges weakly to $\delta_{-r}$. By
Lemma~\ref{lem:lsc} (\emph{lower semicontinuity of $\KG$}), then by Lemma~\ref{lem:atom} (\emph{the largest atom}) and $\Gk\subset\Hk$ (as in the proof of
Proposition~\ref{prop:kinfty}(iii)),
\begin{equation*}
  \liminf_n\KG(R_n)\ge\KG(\delta_{-r})\ge\log\frac1{\qk}=\Kinfty(\kappa):
\end{equation*}
this is (b).

\emph{Exhaustiveness.} If $c<\infty$, then $s\in(0,\infty)$ (Case 1), $s=0$ (Case 2), or $s=\infty$, and then
$r=\lim c_n/s_n=0$ (Case 4). If $c=\infty$, then $r=\infty$ (Case 3) or $r<\infty$, and then $s_n\ge c_n/(r+1)$ for $n$
large, so $s=\infty$ (Case 4).
\end{proof}

\subsection{Proof of Proposition~\ref{prop:regularity}}\label{sec:proof-prop-regularity}

\begin{proof}
\proves{prop:regularity}
\uses{def:kinf, def:kinfty, def:class, lem:repr, lem:strict, lem:escape, lem:mixture, lem:lsc, lem:G}
\emph{(i) Lower semicontinuity.} Let $\ell:=\liminf_n\Kc(P_n,\mu_n)$; pass to a subsequence realising it and assume
$\ell<\infty$. By Lemmas~\ref{lem:repr} and \ref{lem:strict}(a),
\begin{equation*}
  \Kc(P_n,\mu_n)=\Kge(P_n,\mu_n)=\inf_{c\ge\mu_n,\ s>0}\KG(\tauc_{c,s}P_n),
\end{equation*}
so there are $c_n\ge\mu_n$ and $s_n>0$ with $\KG(\tauc_{c_n,s_n}P_n)\le\Kc(P_n,\mu_n)+1/n$. Since $(\mu_n)$ converges,
$\inf_nc_n>-\infty$, and Lemma~\ref{lem:escape} applies. In its case (a), $c:=\lim c_n\ge\mu$ and
\begin{equation*}
  \ell\ge\KG(\tauc_{c,s}P)\ge\Kge(P,\mu)=\Kc(P,\mu),
\end{equation*}
by Lemmas~\ref{lem:repr} and \ref{lem:strict}(a). In its case (b), $\ell\ge\Kinfty(\kappa)\ge\Kc(P,\mu)$ by
Lemma~\ref{lem:mixture}, since $P$ has a finite fourth moment.

\emph{(ii) Continuity and limit.} The function $\Kc(P,\cdot)$ is non-decreasing and right-continuous
(Lemma~\ref{lem:strict}(a)). If $\mu_n\uparrow\mu$, then $\lim_n\Kc(P,\mu_n)\le\Kc(P,\mu)$ by monotonicity and
$\liminf_n\Kc(P,\mu_n)\ge\Kc(P,\mu)$ by (i) with $P_n:=P$; so $\Kc(P,\cdot)$ is also left-continuous. For the limit, let
$\mu_n\to\infty$ and $\ell:=\liminf_n\Kc(P,\mu_n)$; as in (i), pass first to a subsequence realising $\ell$ and assume
$\ell<\infty$. Take $c_n\ge\mu_n$ and $s_n>0$ with $\KG(\tauc_{c_n,s_n}P)\le\Kc(P,\mu_n)+1/n$, which exist by
Lemmas~\ref{lem:repr} and \ref{lem:strict}(a). Lemma~\ref{lem:escape} applies with $P_n:=P$, and its case (a) is impossible,
since $c_n\to\infty$ along every subsequence; so case (b) holds along a further subsequence, $\ell\ge\Kinfty(\kappa)$, and
Lemma~\ref{lem:mixture} gives the converse inequality $\Kc(P,\mu_n)\le\Kinfty(\kappa)$.

\emph{(iii) Positivity.} Let $\mu>\mean(P)$ and suppose $\Kc(P,\mu)=0$. By Lemmas~\ref{lem:repr} and \ref{lem:strict}(a), there
are $c_n\ge\mu$ and $s_n>0$ with $\KG(\tauc_{c_n,s_n}P)\to0$. Case (b) of Lemma~\ref{lem:escape} is impossible since
$\Kinfty(\kappa)>0$. In case (a), $\KG(\tauc_{c,s}P)=0$ with $c\ge\mu$; by Lemma~\ref{lem:lsc} the infimum is attained, so
$\tauc_{c,s}P\in\Gk$, and Lemma~\ref{lem:G}(b) gives $\mean(P)=c\ge\mu>\mean(P)$, a contradiction.
\end{proof}

\subsection{Proof of Corollary~\ref{cor:largegap}}\label{sec:proof-cor-largegap}

\begin{proof}
\proves{cor:largegap}
\uses{def:kinf, def:kinfty, def:class, lem:G, lem:strict, prop:regularity}
The \term{w:law}{law} $P$ is the image of $P_0$ under $x\mapsto m+\sigma x$, so Lemma~\ref{lem:G}(d) with $(c,s,\mu)=(m,\sigma,\Delta/\sigma)$ gives
$\Kc(P,m+\Delta)=\Kc(P_0,\Delta/\sigma)$. For $\Delta\le0$, $P$ itself is an alternative of mean $m\ge m+\Delta$, so
$\Kc(P,m+\Delta)=\Kge(P,m+\Delta)=0$ by Lemma~\ref{lem:strict}(a). The \term{w:law}{law} $P$ has a finite fourth moment and is not
a Dirac mass, since $\sigma>0$; so Proposition~\ref{prop:regularity}(ii)--(iii) gives the continuity, the positivity for
$\Delta>0$ and the limit, and Lemma~\ref{lem:strict}(a) the monotonicity.
\end{proof}

% ----------------------------------------------------------------
\section{\texorpdfstring{Proof of Lemma~\ref{lem:index}}{Proof of Lemma 5.1}}\label{sec:proof-index}

\begin{proof}[Proof of Lemma~\ref{lem:index}]
\proves{lem:index}
\uses{def:algorithm, def:kinf, def:kinfty, def:dualset, def:normalised, lem:repr, lem:strict, lem:finite, lem:lambda, prop:kinfty}
\emph{Setting.} Fix $n\ge1$. For $x=(x_1,\dots,x_n)\in\R^n$ let $\hat P_x:=\frac1n\sum_i\delta_{x_i}$ and
\begin{equation*}
  F(x,\lambda):=\frac1n\sum_{i\le n}\log\hl_\lambda(x_i)\in[-\infty,\infty),\qquad D(x):=\max_{\lambda\in\Lam}F(x,\lambda)=\KG(\hat P_x),
\end{equation*}
the last equality by Lemma~\ref{lem:finite} (\emph{duality for finitely supported laws}; group equal $x_i$ into one atom). \emph{Goal.} Show that $D$ is
continuous, then deduce (i); (ii) and (iii) follow from Proposition~\ref{prop:kinfty}.

\emph{Step 1: $D$ is upper semicontinuous.} $F$ is jointly upper semicontinuous on $\R^n\times\Lam$ (a finite sum
of $\log$, upper semicontinuous on $[0,\infty)$, composed with the continuous $(\lambda,y)\mapsto\hl_\lambda(y)$). If $x_k\to x$
and $\lambda_k$ maximises $F(x_k,\cdot)$, a subsequence realising $\limsup_kD(x_k)$ has $\lambda_k\to\bar\lambda\in\Lam$ (compactness,
Lemma~\ref{lem:lambda}(a)), and $\limsup_kD(x_k)\le F(x,\bar\lambda)\le D(x)$.

\emph{Step 2: $D$ is lower semicontinuous.} For $\lambda\in\mathrm{int}\,\Lam$, $\hl_\lambda>0$ on $\R$ (Lemma~\ref{lem:lambda}(c)), so
$F(\cdot,\lambda)$ is continuous and finite. Moreover $D(x)=\sup_{\lambda\in\mathrm{int}\,\Lam}F(x,\lambda)$: if $\lambda^\star$ maximises $F(x,\cdot)$
(with $F(x,\lambda^\star)\ge F(x,0)=0$) and $\lambda^0\in\mathrm{int}\,\Lam$, then $\lambda_t:=(1-t)\lambda^\star+t\lambda^0\in\mathrm{int}\,\Lam$ for
$t\in(0,1]$, and by concavity of $\log$ and affinity of $\hl$ in $\lambda$,
\begin{equation*}
  F(x,\lambda_t)\ge(1-t)F(x,\lambda^\star)+tF(x,\lambda^0)\longrightarrow F(x,\lambda^\star)\qquad(t\to0).
\end{equation*}
A supremum of continuous functions is lower semicontinuous. So $D$ is continuous on $\R^n$, with values in
$[0,\infty)$.

\emph{Step 3: the information function.} By Lemma~\ref{lem:repr}, since $\tauc_{c,s}\hat P_x=\hat P_{(x-c)/s}$,
\begin{equation*}
  \Kc(\hat P_x,\mu)=\inf_{c>\mu,\ s>0}D\Bigl(\frac{x-c}{s}\Bigr)=\inf_{(r,s)\in(0,\infty)^2}D\Bigl(\frac{x-\mu-r}{s}\Bigr).
\end{equation*}
Each map $(x,\mu)\mapsto D((x-\mu-r)/s)$ is continuous, and by continuity the infimum over $(0,\infty)^2$ equals the
infimum over the countable set $(\mathbb Q\cap(0,\infty))^2$. So $(x,\mu)\mapsto\Kc(\hat P_x,\mu)$ is upper semicontinuous on
$\R^n\times\R$, hence Borel, which proves the first claim of (i).

\emph{Step 4: the index.} At the end of round $t$, with $n:=\Np_a(t)$ and $x$ the observed rewards of arm $a$, the set
$S:=\{\mu:n\Kc(\hat P_x,\mu)\le\thr(t,n)\}$ is an interval unbounded below, because $\Kc(\hat P_x,\cdot)$ is non-decreasing
(Lemma~\ref{lem:strict}(a)). For $\mu\in\mathbb Q$ let $\psi_\mu:=\mu$ if $\mu\in S$ and $\psi_\mu:=-\infty$ otherwise, a Borel function of
$x$ by Step 3. Since $S$ is an interval unbounded below,
\begin{equation*}
  \Ix_a(t)=\sup S=\sup_{\mu\in\mathbb Q}\psi_\mu ,
\end{equation*}
a countable supremum of Borel functions of $x$. The number $n$ and the past choices are measurable
functions of the history, so $\Ix_a(t)$ is a Borel function of the history and $\Arm_{t+1}$ is a measurable choice
(with the independent randomness of the tie rule). This proves (i).

\emph{Step 5: (ii) and (iii).} An empirical \term{w:law}{law} has a finite fourth moment, so $\Kc(\hat P_x,\mu)\le\Kinfty(\kappa)$ for
every $\mu$ by Proposition~\ref{prop:kinfty}(ii) (\emph{the value at infinity bounds $\Kc$}). If $n\Kinfty(\kappa)\le\thr(t,n)$, then $S=\R$ and $\Ix_a(t)=+\infty$: this is
(ii). If all the $x_i$ are equal to $y$, then $\hat P_x=\delta_y$ and $\Kc(\delta_y,\mu)=\Kinfty(\kappa)$ for every $\mu$
(Proposition~\ref{prop:kinfty}(iii)); so $S=\R$ when $n\Kinfty(\kappa)\le\thr(t,n)$ and $S=\emptyset$ otherwise, which gives (iii).
\end{proof}

% ----------------------------------------------------------------
\section{\texorpdfstring{Proof of Theorem~\ref{thm:main} and Corollary~\ref{cor:exact}}{Proof of Theorem 6.1 and Corollary 6.2}}\label{sec:proof-main}

Throughout this appendix, $\kappa>1$, $\nu\in\Hk^\alpha$, $a$ is a \term{w:subopt}{suboptimal} arm, $a^\star$ an arm with $\marm_{a^\star}=\mstar$,
and $\varepsilon\in(0,\Ka)$. Lemmas~\ref{lem:decomp} and~\ref{lem:termB} hold for every $\delta>0$; only
Lemma~\ref{lem:termC} uses the particular $(\delta,J,\eta)$ given by Theorem~\ref{thm:deviation}(ii) applied to
$(P,\mu,\varepsilon):=(\nu_a,\mstar,\varepsilon)$. We will often use the fundamental inequality of index policies: for every
$s\ge1$, since $\Np_a(t)=s$ and $\Arm_{t+1}=a$ imply $\Np_a(t+1)=s+1$,
\begin{equation}\label{eq:fundamental}
  \sum_{t\ge1}\ind\bigl(\Arm_{t+1}=a,\ \Np_a(t)=s\bigr)\le1 .
\end{equation}

To prove Theorem~\ref{thm:main} we use the following decomposition.

\begin{lemma}[Decomposition]\label{lem:decomp}
\uses{def:algorithm, def:bandit, def:kinf}
For every $\delta>0$, every $T\ge3$ and every run of \algo{},
\begin{equation*}
  \Np_a(T)\ \le\ 1+\underbrace{\sum_{t=\Karms}^{T-1}\ind\bigl(\Ix_{a^\star}(t)<\mstar\bigr)}_{A}
  +\underbrace{\sum_{s=1}^{T}\ind\bigl(s\,\Kc(\Pe_{a,s},\mstar-\delta)\le\thr(T,s)\bigr)}_{B}.
\end{equation*}
The term $A$ does not depend on $\delta$; the term $B$ does.
\end{lemma}

(See the proof in Subsection~\ref{sec:proof-lem-decomp}.)

From this decomposition we use the following lemmas to control each term.

\begin{lemma}[The optimal arm]\label{lem:termA}
\uses{def:algorithm, def:class, def:kinf, thm:deviation}
$\E[A]\le\frac32$.
\end{lemma}

(See the proof in Subsection~\ref{sec:proof-lem-termA}.) This is the classical argument for the optimal arm
of KL-UCB-type indices, as in \citet[Section~3.3]{AJK21}; the new ingredient is the upper deviation inequality
of Theorem~\ref{thm:deviation}(i) for the class $\Hk$.

\begin{lemma}[Counting the \term{w:subopt}{suboptimal} pulls]\label{lem:termB}
\uses{def:algorithm, def:kinf, lem:decomp}
For every $\delta>0$, every $T\ge3$ and every run, with $B$ as in Lemma~\ref{lem:decomp} for this $\delta$,
\begin{equation*}
  B\le\nT(\varepsilon)+\underbrace{\sum_{s\ge1}\ind\bigl(\Kc(\Pe_{a,s},\mstar-\delta)\le\Ka-\varepsilon\bigr)}_{C},
\end{equation*}
and, with $\kep=\Ka-\varepsilon$ as in Theorem~\ref{thm:main},
\begin{equation*}
  \nT(\varepsilon)\le\frac{\LT+2\log\bigl(1+2(\LT+2\log(4/\kep))/\kep\bigr)}{\kep},\qquad
  \limsup_{T\to\infty}\frac{\nT(\varepsilon)}{\log T}\le\frac1{\kep}.
\end{equation*}
\end{lemma}

(See the proof in Subsection~\ref{sec:proof-lem-termB}.) Classical: no contribution here.

\begin{lemma}[The \term{w:subopt}{suboptimal} arm]\label{lem:termC}
\uses{def:class, def:kinf, thm:deviation, lem:termB}
Let $(\delta,J,\eta)$ be given by Theorem~\ref{thm:deviation}(ii) applied to $(P,\mu,\varepsilon):=(\nu_a,\mstar,\varepsilon)$, and let
$C$ be as in Lemma~\ref{lem:termB} for this $\delta$. Then
$\E[C]\le\dt{s:Ca}\Ca(\varepsilon):=\dfrac{2J}{e^{2\eta^2}-1}<\infty$.
\end{lemma}

(See the proof in Subsection~\ref{sec:proof-lem-termC}.) The argument is classical, but the lower deviation
inequality it rests on, Theorem~\ref{thm:deviation}(ii), is new for this class and is where the escape of mass
is controlled.

From these lemmas the proof of the theorem is short.

\begin{proof}[Proof of Theorem~\ref{thm:main}]
\proves{thm:main}
\uses{def:bandit, def:class, def:kinf, def:kinfty, def:algorithm, lem:decomp, lem:termA, lem:termB, lem:termC, prop:kinfty, prop:regularity, thm:deviation}
\emph{The constant $\Ka$.} By Proposition~\ref{prop:kinfty}(ii) (\emph{the value at infinity bounds $\Kc$}) and Proposition~\ref{prop:regularity}(iii) (\emph{positivity}) (the \term{w:law}{law}
$\nu_a\in\Hk$ has a finite fourth moment and positive variance, and $\mstar>\marm_a$), $0<\Ka\le\Kinfty(\kappa)$, and
$\Kinfty(\kappa)\le\log2<1$ since $\qk>\frac12$; so $\kep\in(0,1)$.

\emph{(i).} Let $(\delta,J,\eta)$ be given by Theorem~\ref{thm:deviation}(ii) at $(\nu_a,\mstar,\varepsilon)$. The constant
$\Ca(\varepsilon)$ of Lemma~\ref{lem:termC} depends only on $(\delta,J,\eta)$, hence, by the statement of
Theorem~\ref{thm:deviation}(ii), only on $\nu_a$, $\mstar$, $\varepsilon$ and $\kappa$, and not on $T$. Taking expectations in
Lemma~\ref{lem:decomp} with this $\delta$, then using Lemmas~\ref{lem:termA}, \ref{lem:termB} (with the same $\delta$) and
\ref{lem:termC}, we get for every $T\ge3$
\begin{align*}
  \E[\Np_a(T)]&\le1+\E[A]+\E[B]\\
  &\le1+\frac32+\nT(\varepsilon)+\E[C]\\
  &\le\nT(\varepsilon)+\frac52+\Ca(\varepsilon),
\end{align*}
and the bound on $\nT(\varepsilon)$ of Lemma~\ref{lem:termB} gives \eqref{eq:main-finite}.

\emph{(ii).} Dividing (i) by $\log T$ and using Lemma~\ref{lem:termB},
$\limsup_T\E[\Np_a(T)]/\log T\le1/(\Ka-\varepsilon)$ for every $\varepsilon\in(0,\Ka)$; letting $\varepsilon\to0$ gives the first claim. Summing
$\gap_a$ times it over the \term{w:subopt}{suboptimal} arms gives the bound on $\Reg_T$, since $\Reg_T=\sum_a\gap_a\E[\Np_a(T)]$.
\end{proof}

\begin{proof}[Proof of Corollary~\ref{cor:exact}]
\proves{cor:exact}
\uses{def:bandit, def:class, def:kinf, def:kinfty, def:algorithm, cite:LS-16.2, thm:main, prop:kinfty, prop:regularity, lem:strict, lem:expansion}
\emph{(i).} For every \term{w:subopt}{suboptimal} arm, $0<\Kc(\nu_a,\mstar)\le\Kinfty(\kappa)$ (Propositions~\ref{prop:regularity}(iii) and
\ref{prop:kinfty}(ii)), so $\Cst(\nu)$ is a finite sum of finite positive numbers, positive since some arm is
\term{w:subopt}{suboptimal}. The \term{w:law}{laws} of $\Hk$ have finite means, and $\Kinf(\nu_a,\mstar;\Hk)=\Kc(\nu_a,\mstar)$, so
Theorem~\ref{cite:LS-16.2} with $\mathcal D:=\Hk$ gives the lower bound.

\emph{(ii).} By Theorem~\ref{thm:main}(i), for every $\nu\in\Hk^\alpha$, $\Reg_T(\nu)=O(\log T)=o(T^p)$ for every $p>0$: \algo{}
is \term{w:consistent}{consistent} over $\Hk^\alpha$. Then (i) and Theorem~\ref{thm:main}(ii) give
\begin{equation*}
  \Cst(\nu)\le\liminf_{T\to\infty}\frac{\Reg_T}{\log T}\le\limsup_{T\to\infty}\frac{\Reg_T}{\log T}\le\Cst(\nu).
\end{equation*}

\emph{(iii).} $\Kc(\nu_a,\mstar)\le\Kinfty(\kappa)$ and $1/\Kinfty(\kappa)>\kappa$ (Lemma~\ref{lem:expansion}).

\emph{(iv).} The identity is Lemma~\ref{lem:strict}(b). For $\nu\in(\Hlt)^\alpha$, the constant of
Theorem~\ref{cite:LS-16.2} with $\mathcal D:=\Hlt$ is therefore $\Cst(\nu)$. Since $(\Hlt)^\alpha\subset\Hk^\alpha$,
Theorem~\ref{thm:main} applies to $\nu$, and \algo{} is \term{w:consistent}{consistent} over $(\Hlt)^\alpha$. Theorem~\ref{cite:LS-16.2}
(\emph{the lower bound for consistent algorithms}) with $\mathcal D:=\Hlt$ gives $\liminf_T\Reg_T/\log T\ge\Cst(\nu)$ for every algorithm
\term{w:consistent}{consistent} over $(\Hlt)^\alpha$, and Theorem~\ref{thm:main}(ii) gives $\limsup_T\Reg_T/\log T\le\Cst(\nu)$ for \algo{}.
\end{proof}

\subsection{Proof of Lemma~\ref{lem:decomp}}\label{sec:proof-lem-decomp}

\begin{proof}
\proves{lem:decomp}
\uses{def:algorithm, def:bandit, def:kinf, lem:strict}
If $T<\Karms+1$, then $\Np_a(T)\le1$. Otherwise each arm is pulled once in rounds $1,\dots,\Karms$, so
\begin{align*}
  \Np_a(T)&=1+\sum_{t=\Karms}^{T-1}\ind(\Arm_{t+1}=a)\\
  &\le1+\sum_{t=\Karms}^{T-1}\ind\bigl(\Ix_{a^\star}(t)<\mstar\bigr)+\sum_{t=\Karms}^{T-1}\ind\bigl(\Ix_{a^\star}(t)\ge \mstar,\ \Arm_{t+1}=a\bigr)\\
  &\le1+A+\sum_{t=\Karms}^{T-1}\ind\bigl(\Ix_a(t)\ge \mstar,\ \Arm_{t+1}=a\bigr),
\end{align*}
since $\Arm_{t+1}=a$ implies $\Ix_a(t)\ge\Ix_{a^\star}(t)$ (whatever the tie rule). Let us look at the last sum. If
$\Ix_a(t)\ge \mstar$, then $\Ix_a(t)>\mstar-\delta$, so some $\mu''>\mstar-\delta$ satisfies
$\Np_a(t)\Kc(\Pe_a(t),\mu'')\le\thr(t,\Np_a(t))$; since $\Kc(\Pe_a(t),\cdot)$ is non-decreasing (Lemma~\ref{lem:strict}(a)) and
$\thr(\cdot,n)$ is non-decreasing with $t\le T$,
\begin{equation*}
  \Np_a(t)\,\Kc\bigl(\Pe_a(t),\mstar-\delta\bigr)\le\thr\bigl(T,\Np_a(t)\bigr).
\end{equation*}
For $t\in\{\Karms,\dots,T-1\}$, $\Np_a(t)\in\{1,\dots,T\}$ and $\Pe_a(t)=\Pe_{a,\Np_a(t)}$ (stack model). So, splitting according to the
value $s$ of $\Np_a(t)$ and using \eqref{eq:fundamental},
\begin{align*}
  \sum_{t=\Karms}^{T-1}\ind\bigl(\Ix_a(t)\ge \mstar,\ \Arm_{t+1}=a\bigr)
  &\le\sum_{t=\Karms}^{T-1}\sum_{s=1}^{T}\ind\bigl(\Arm_{t+1}=a,\ \Np_a(t)=s\bigr)\,\ind\bigl(s\,\Kc(\Pe_{a,s},\mstar-\delta)\le\thr(T,s)\bigr)\\
  &\le\sum_{s=1}^{T}\ind\bigl(s\,\Kc(\Pe_{a,s},\mstar-\delta)\le\thr(T,s)\bigr)=B .
\end{align*}
\end{proof}

\subsection{Proof of Lemma~\ref{lem:termA}}\label{sec:proof-lem-termA}

\begin{proof}
\proves{lem:termA}
\uses{def:algorithm, def:class, def:kinf, thm:deviation, lem:strict}
Let $t\ge\Karms$ and suppose $\Ix_{a^\star}(t)<\mstar$. Pick $\mu\in(\Ix_{a^\star}(t),\mstar)$ (any $\mu<\mstar$ if $\Ix_{a^\star}(t)=-\infty$). Then $\mu$ is
not in the set defining $\Ix_{a^\star}(t)$, and $\Kc(\Pe_{a^\star}(t),\cdot)$ is non-decreasing (Lemma~\ref{lem:strict}(a)), so
\begin{align*}
  \ind\bigl(\Ix_{a^\star}(t)<\mstar\bigr)
  &\le\ind\bigl(\Np_{a^\star}(t)\Kc(\Pe_{a^\star}(t),\mu)>\thr(t,\Np_{a^\star}(t))\bigr)\\
  &\le\ind\bigl(\Np_{a^\star}(t)\Kc(\Pe_{a^\star}(t),\mstar)>\thr(t,\Np_{a^\star}(t))\bigr)\\
  &\le\ind\bigl(\exists s\ge1:\ s\,\Kc(\Pe_{a^\star,s},\mstar)\ge x_t+1+2\log(s+1)\bigr),
\end{align*}
with $x_t:=\log(t\vee3)+2\log\log(t\vee3)\ge0$. The rewards $\Xr_{a^\star,1},\Xr_{a^\star,2},\dots$ are independent with \term{w:law}{law}
$\nu_{a^\star}\in\Hk$ of mean $\mstar$, so Theorem~\ref{thm:deviation}(i) (\emph{the anytime upper deviation inequality}) gives
\begin{equation*}
  \Prob\bigl(\Ix_{a^\star}(t)<\mstar\bigr)\le e^{-x_t}=\frac1{(t\vee3)\log^2(t\vee3)} .
\end{equation*}
Finally, since $t\mapsto1/(t\log^2t)$ is decreasing on $[3,\infty)$ and $\Karms\ge2$,
\begin{equation*}
  \E[A]\le\frac1{3\log^23}+\sum_{t\ge3}\frac1{t\log^2t}\le\frac2{3\log^23}+\int_3^\infty\frac{dt}{t\log^2t}=\frac2{3\log^23}+\frac1{\log3}\le1.47\le\frac32 .
\end{equation*}
\end{proof}

\subsection{Proof of Lemma~\ref{lem:termB}}\label{sec:proof-lem-termB}

\begin{proof}
\proves{lem:termB}
\uses{def:algorithm, def:kinf}
\emph{The split.} For each $s$, we split according to the event $\{\Kc(\Pe_{a,s},\mstar-\delta)>\Ka-\varepsilon\}$:
\begin{align*}
  \ind\bigl(s\Kc(\Pe_{a,s},\mstar-\delta)\le\thr(T,s)\bigr)
  &\le\ind\bigl(s\Kc(\Pe_{a,s},\mstar-\delta)\le\thr(T,s),\ \Kc(\Pe_{a,s},\mstar-\delta)>\Ka-\varepsilon\bigr)\\
  &\qquad+\ind\bigl(\Kc(\Pe_{a,s},\mstar-\delta)\le\Ka-\varepsilon\bigr)\\
  &\le\ind\bigl(s(\Ka-\varepsilon)\le\thr(T,s)\bigr)+\ind\bigl(\Kc(\Pe_{a,s},\mstar-\delta)\le\Ka-\varepsilon\bigr).
\end{align*}
Every $s$ with $s(\Ka-\varepsilon)\le\thr(T,s)$ satisfies $s\le\nT(\varepsilon)$, so summing over $s\in\{1,\dots,T\}$ gives
$B\le\nT(\varepsilon)+C$.

\emph{The bound on $\nT(\varepsilon)$.} Let $k:=\kep\in(0,1)$ and $n\ge0$ with $nk\le\thr(T,n)=\LT+2\log(1+n)$ (for $T\ge3$,
$t\vee3=T$). By concavity of $\log$, $\log y\le\log c+(y-c)/c$ for $y,c>0$; with $y:=1+n$ and $c:=4/k$,
\begin{equation*}
  2\log(1+n)\le2\log\frac4k+\frac{k(1+n)}2-2 .
\end{equation*}
So $nk\le\LT+2\log(4/k)+k/2-2+nk/2$, and since $k/2-2<0$,
\begin{equation*}
  n\le\bar n:=\frac{2(\LT+2\log(4/k))}{k}.
\end{equation*}
Using this in $nk\le\LT+2\log(1+n)$ gives $nk\le\LT+2\log(1+\bar n)$, which is the bound on $\nT(\varepsilon)$. Finally,
$\LT/\log T\to1$ and $\log(1+\bar n)/\log T\to0$ as $T\to\infty$, since $\bar n=O(\log T)$; this gives the $\limsup$.
\end{proof}

\subsection{Proof of Lemma~\ref{lem:termC}}\label{sec:proof-lem-termC}

\begin{proof}
\proves{lem:termC}
\uses{def:class, def:kinf, thm:deviation}
The rewards $\Xr_{a,1},\Xr_{a,2},\dots$ are independent with \term{w:law}{law} $\nu_a\in\Hk$, and $\Pe_{a,s}$ is their empirical \term{w:law}{law} after $s$
samples. By the choice of $(\delta,J,\eta)$ (Theorem~\ref{thm:deviation}(ii), \emph{the lower deviation inequality}, at $(\nu_a,\mstar,\varepsilon)$),
\begin{equation*}
  \E[C]=\sum_{s\ge1}\Prob\bigl(\Kc(\Pe_{a,s},\mstar-\delta)\le\Kc(\nu_a,\mstar)-\varepsilon\bigr)\le\sum_{s\ge1}2Je^{-2s\eta^2}=\frac{2J}{e^{2\eta^2}-1}.
\end{equation*}
\end{proof}

% ----------------------------------------------------------------
\section{\texorpdfstring{Proof of Theorem~\ref{thm:deviation} and Corollary~\ref{cor:confseq}}{Proof of Theorem 7.1 and Corollary 7.2}}\label{sec:proof-deviation}

To prove Theorem~\ref{thm:deviation}(i) we use the following mixture bound.

\begin{lemma}[Mixture over the dual set]\label{lem:mixbound}
\uses{def:dualset, lem:lambda}
Let $\pi$ be the uniform \term{w:law}{law} on $\Lam$ (Lebesgue measure restricted to $\Lam$, normalised), $n\ge1$,
$y_1,\dots,y_n\in\R$ and $M(\lambda):=\prod_{i\le n}\hl_\lambda(y_i)$. Then
\begin{equation*}
  \int_\Lam M\,d\pi\ \ge\ \frac{1}{e(n+1)^2}\,\max_{\lambda\in\Lam}M(\lambda).
\end{equation*}
\end{lemma}

(See the proof in Subsection~\ref{sec:proof-lem-mixbound}.) This is the bound of \citet[Lemma~E.1 of the arXiv
version]{AKJ21} on the maximum of a sum of exp-concave functions, restated by \citet[Lemma~11]{AJK21}, with the
same constants; \citet{AKJ21} present it as the regret bound of the continuous exponentially weighted average
predictor, shown for universal portfolios by Blum and Kalai and in general by Hazan, Agarwal and Kale. We give
the short direct proof for our affine integrands.

\begin{proof}[Proof of Theorem~\ref{thm:deviation}]
\proves{thm:deviation}
\uses{def:class, def:kinf, def:laws, def:normalised, def:dualset, lem:G, lem:repr, lem:strict, lem:finite, lem:lambda, lem:mixbound, lem:index, prop:regularity}
\emph{(i) Idea.} Fix the location and the scale of the alternatives at the true mean and standard deviation:
this bounds $n\Kc(\hat P_n,\mean(P))$ by a maximum over the fixed set $\Lam$ of sums of logarithms of variables of
mean at most $1$, and a mixture over $\Lam$ is a supermartingale.

\emph{Step 1: the supermartingale.} Let $m:=\mean(P)$, $\sigma:=\sd(P)>0$ and $Y_i:=(X_i-m)/\sigma$. Then $\E Y_i=0$ and
$\E\phik(Y_i)=\kurt(P)-2\kappa+\kappa\le0$, so for every $\lambda\in\Lam$
\begin{equation*}
  \E\,\hl_\lambda(Y_i)=1+\lambda_2\E\phik(Y_i)\le1 .
\end{equation*}
Let $W_0:=1$ and $W_n:=\int_\Lam\prod_{i\le n}\hl_\lambda(Y_i)\,\pi(d\lambda)$, with $\pi$ as in Lemma~\ref{lem:mixbound} ($\Lam$ has positive
Lebesgue measure by Lemma~\ref{lem:lambda}(b)). The integrand is non-negative and jointly measurable, so by
Tonelli's theorem and the independence of $Y_n$ from $Y_1,\dots,Y_{n-1}$,
\begin{equation*}
  \E[W_n\mid Y_1,\dots,Y_{n-1}]=\int_\Lam\prod_{i<n}\hl_\lambda(Y_i)\,\E\bigl[\hl_\lambda(Y_n)\bigr]\,\pi(d\lambda)\le W_{n-1}.
\end{equation*}
So $(W_n)$ is a non-negative supermartingale with $W_0=1$, and the maximal inequality \citep[Theorem~3.9]{LS} gives
\begin{equation*}
  \Prob\bigl(\exists n\ge1:\ W_n\ge e^x\bigr)\le e^{-x}.
\end{equation*}

\emph{Step 2: the link with $\Kc$.} The \term{w:law}{law} $\tauc_{m,\sigma}\hat P_n$ is the empirical \term{w:law}{law} of $Y_1,\dots,Y_n$, and $c=m$ is admissible
in the non-strict representation. By Lemma~\ref{lem:strict}(a) (\emph{strict and non-strict versions}), Lemma~\ref{lem:repr} (\emph{representation}), Lemma~\ref{lem:finite} (\emph{duality for finitely supported laws}) and Lemma~\ref{lem:mixbound},
\begin{align*}
  n\,\Kc(\hat P_n,m)&=n\,\Kge(\hat P_n,m)\\
  &\le n\,\KG(\tauc_{m,\sigma}\hat P_n)\\
  &=\max_{\lambda\in\Lam}\sum_{i\le n}\log\hl_\lambda(Y_i)\\
  &\le\log W_n+1+2\log(n+1).
\end{align*}
Hence the event of \eqref{eq:upper-dev} is contained in $\{\exists n:W_n\ge e^x\}$, which proves \eqref{eq:upper-dev}. For
the consequence, let $n\ge1$ and $x\in\R$, and $x':=x-1-2\log(n+1)$. If $x'<0$ the bound $e(n+1)^2e^{-x}=e^{-x'}>1$ is trivial;
otherwise \eqref{eq:upper-dev} with $x'$ gives $\Prob(n\Kc(\hat P_n,m)\ge x)\le e^{-x'}=e(n+1)^2e^{-x}$.

\emph{(ii) Idea.} Lower semicontinuity at $(P,\mu)$ gives a weak neighbourhood of $P$, described by finitely many
test functions, on which $\Kc(\cdot,\mu-\delta)$ stays above $\Kc(P,\mu)-\varepsilon$; Hoeffding's inequality controls the
empirical means of the test functions.

\emph{Step 3: a neighbourhood.} If $\Kc(P,\mu)-\varepsilon<0$, the event of \eqref{eq:lower-dev} is empty and any $(\delta,J,\eta)$ works.
Otherwise, $P\in\Hk$ has a finite fourth moment and is not a Dirac mass ($\sd(P)>0$), so $\Kc$ is lower
semicontinuous at $(P,\mu)$ by Proposition~\ref{prop:regularity}(i) (\emph{lower semicontinuity away from Dirac masses}). Since $\Pspace\times\R$ is metrisable, lower
semicontinuity along sequences gives an open neighbourhood $\mathcal N$ of $(P,\mu)$ on which
$\Kc>\Kc(P,\mu)-\varepsilon$ (otherwise a sequence of points of $\{\Kc\le\Kc(P,\mu)-\varepsilon\}$ would converge to $(P,\mu)$). The sets
$\{R:|\int f_j\,dR-\int f_j\,dP|<\eta,\ j\le J\}$, for bounded continuous $f_j$ and $\eta>0$, form a base of neighbourhoods of
$P$. Dropping the constant $f_j$, replacing each other $f_j$ by $(f_j-\inf f_j)/(\sup f_j-\inf f_j)$ and $\eta$ by
$\min_j\eta/(\sup f_j-\inf f_j)$ only shrinks such a set, so we may take $f_j$ with values in $[0,1]$ and a single $\eta$. So there are
$\delta>0$, $J$, continuous $f_1,\dots,f_J:\R\to[0,1]$ and $\eta>0$ with
\begin{equation*}
  \Bigl\{R:\Bigl|\int f_j\,dR-\int f_j\,dP\Bigr|<\eta\ \text{for all } j\le J\Bigr\}\times[\mu-\delta,\mu+\delta]\subset\mathcal N .
\end{equation*}

\emph{Step 4: Hoeffding.} Since $(\hat P_n,\mu-\delta)\in\mathcal N$ as soon as all the $J$ empirical means are $\eta$-close to their
expectations,
\begin{align*}
  \Prob\bigl(\Kc(\hat P_n,\mu-\delta)\le\Kc(P,\mu)-\varepsilon\bigr)
  &\le\Prob\Bigl(\exists j\le J:\ \Bigl|\frac1n\sum_{i\le n}f_j(X_i)-\E f_j(X_1)\Bigr|\ge\eta\Bigr)\\
  &\le\sum_{j\le J}2e^{-2n\eta^2}=2Je^{-2n\eta^2},
\end{align*}
by Hoeffding's inequality for variables in $[0,1]$. The events are measurable by Lemma~\ref{lem:index}(i).
\end{proof}

\begin{proof}[Proof of Corollary~\ref{cor:confseq}]
\proves{cor:confseq}
\uses{def:class, def:kinf, thm:deviation}
The \term{w:law}{law} of $-X_1$ belongs to $\Hk$ (the kurtosis is invariant under $x\mapsto-x$) and has mean $-\mean(P)$. Let $E^+$ and
$E^-$ be the events of \eqref{eq:upper-dev} for $X$ and for $-X$; each has probability at most $e^{-x}$. Outside
$E^+\cup E^-$, for every $n$, $n\Kc(\hat P_n,\mean(P))<b_n$ and $n\Kc(\hat P^-_n,-\mean(P))<b_n$, so $\mean(P)$ belongs to the set defining
$U^+_n$ and $-\mean(P)$ to the set defining $U^-_n$; hence $-U^-_n\le\mean(P)\le U^+_n$.
\end{proof}

\subsection{Proof of Lemma~\ref{lem:mixbound}}\label{sec:proof-lem-mixbound}

\begin{proof}
\proves{lem:mixbound}
\uses{def:dualset, lem:lambda}
$M$ is continuous on the compact set $\Lam$ (Lemma~\ref{lem:lambda}(a)); let $\lambda^\star$ maximise it, $\theta:=1/(n+1)$ and
$S:=\lambda^\star+\theta(\Lam-\lambda^\star)$. By convexity $S\subset\Lam$, and $\pi(S)=\theta^2$ since $S$ is a homothetic copy of $\Lam$ in
$\R^2$. For $\lambda=\lambda^\star+\theta(\lambda'-\lambda^\star)\in S$, $\lambda'\in\Lam$, the affinity of $\hl$ in $\lambda$ and $\hl_{\lambda'}\ge0$ give
\begin{equation*}
  \hl_\lambda(y)=(1-\theta)\hl_{\lambda^\star}(y)+\theta\hl_{\lambda'}(y)\ge(1-\theta)\hl_{\lambda^\star}(y)\qquad(y\in\R),
\end{equation*}
so $M(\lambda)\ge(1-\theta)^nM(\lambda^\star)$ on $S$. Hence
\begin{equation*}
  \int_\Lam M\,d\pi\ge\pi(S)(1-\theta)^nM(\lambda^\star)=\frac1{(n+1)^2}\Bigl(1+\frac1n\Bigr)^{-n}M(\lambda^\star)\ge\frac{M(\lambda^\star)}{e(n+1)^2}.
\end{equation*}
\end{proof}

% ----------------------------------------------------------------
\section{Sources of the techniques}\label{sec:attribution}

For each technique of the proofs, where it is used and the earliest source we found.

\begin{center}
\small
\begin{tabular}{@{}p{0.47\textwidth}p{0.47\textwidth}@{}}
\toprule
Technique (where) & Earliest source found \\
\midrule
Lower bound for \term{w:consistent}{consistent} algorithms with $\Kinf$ (Corollary~\ref{cor:exact}) & Burnetas and Katehakis, in the form of \citet[Theorem~16.2]{LS} \\
KL-UCB index with $\Kinf$ of a non-parametric class, threshold $\log t+2\log\log t+2\log(1+n)+1$, regret decomposition (Algorithm~\ref{alg:kucb}, Appendix~\ref{sec:proof-main}) & \citet[Algorithm~1, Theorem~1, Section~3.3]{AJK21}; for an abstract class \citet[Theorem~3.1]{PA26}; KL-UCB: \citet{CGMMS13} \\
Outer minimisation over a scalar parameter that makes a non-convex $\Kinf$ problem convex (Lemmas~\ref{lem:G}, \ref{lem:repr}) & \citet[Theorem~3.4]{AKJ21} for CVaR (see also \citealp[Chapter~6]{A22}); for centred moments, over the location, \citet[Section~6.10]{A22} and, with a regret theorem under a known lower bound on the means, \citet[Appendix~D.5 and Corollary~6]{BSH23}; location and scale as nuisance parameters of moment equations: empirical likelihood (see also \citealp[Sections~1.1 and 3.5]{O01}). The inequality $\kurt\le\kappa$ as a linear constraint with a free scale: source not known \\
Fixing the outer parameter at its true value to bound the empirical $\Kinf$ from above (Theorem~\ref{thm:deviation}(i)) & \citet[Appendix~6.D.1, equation~(6.55)]{A22} (the value-at-risk for CVaR) \\
Dual formula of $\Kinf$ under moment constraints, of the form $\max\E\log(1+\lambda\cdot\gam)$ (Lemmas~\ref{lem:weak}, \ref{lem:finite}) & \citet[Theorem~8]{HT10} for bounded supports; \citet[Lemma~7]{AJK21} for a moment bound \\
Direct construction of the optimal \term{w:law}{law} from a maximiser of the dual (Lemma~\ref{lem:finite}) & standard first-order (Karush--Kuhn--Tucker) argument for a concave maximisation over a convex set; the optimal \term{w:law}{law} is read off from the multipliers, as in empirical likelihood (see also \citealp[Section~3.14]{O01}) \\
Extension of the dual formula from finitely supported \term{w:law}{laws} to every \term{w:law}{law} by approximation (Lemma~\ref{lem:alllaws}) & the proof of \citet[Theorem~8]{HT10} (finite supports, then approximation) for bounded supports; the variant with Fatou's lemma on an unbounded support is ours \\
Mixture of supermartingales over the dual set, universal-portfolio bound $d\log(n+1)+1$, maximal inequality (Theorem~\ref{thm:deviation}(i), Lemma~\ref{lem:mixbound}) & in bandits, \citet[Lemma~E.1 of the arXiv version]{AKJ21}, restated by \citet[Lemma~11, used for Proposition~5]{AJK21}; \citet{AKJ21} credit the bound to Blum and Kalai (universal portfolios) and to Hazan, Agarwal and Kale (in general) \\
Mixing an alternative with a far atom (``outlier'') (Lemma~\ref{lem:mixture}) & \citet[Appendix~A]{L17}, with a Bernoulli outlier; the sharp mass $1-\qk$: source not known \\
Kurtosis of two-point \term{w:law}{laws} and the largest atom of a \term{w:law}{law} of kurtosis at most $\kappa$ (Lemma~\ref{lem:atom}) & elementary, of the type of the maximal mass at a point under moment constraints (Chebyshev--Markov--Krein); the statement: source not known \\
Mixing to make a constraint strict (Lemma~\ref{lem:strict}) & classical; source not known \\
Compactness modulo the affine group with the boundary value $\Kinfty(\kappa)$ (Lemma~\ref{lem:escape}, Proposition~\ref{prop:regularity}) & source not known \\
Lower semicontinuity turned into an exponential lower deviation by finitely many test functions and Hoeffding's inequality (Theorem~\ref{thm:deviation}(ii)) & source not known (a standard topological argument); Hoeffding's inequality \\
Joint lower semicontinuity of $\KL$, Prokhorov, Skorokhod, Varadarajan (Lemmas~\ref{lem:lsc}, \ref{lem:alllaws}, \ref{lem:escape}) & textbook facts \\
\bottomrule
\end{tabular}
\end{center}

% ----------------------------------------------------------------
\section{Notation}\label{sec:notation}

Every symbol is a link to its definition.

\begin{center}
\small
\begin{tabular}{@{}p{0.2\textwidth}p{0.55\textwidth}p{0.17\textwidth}@{}}
\toprule
Symbol & Meaning & Defined in \\
\midrule
$\Pspace$ & \term{w:law}{laws} on $\R$, weak topology & Section~\ref{sec:setting} \\
$\mean(P)$, $\sd^2(P)$, $\kurt(P)$ & mean, variance, kurtosis & Definition~\ref{def:laws} \\
$\KL$, $\kl$ & Kullback--Leibler divergence; its Bernoulli version & Definition~\ref{def:laws}; Lemma~\ref{lem:tail} \\
$\alpha$, $\Karms$ & set of arms, number of arms & Definition~\ref{def:bandit} \\
$\nu$, $\marm_a$, $\sda_a$, $\mstar$, $\gap_a$ & \term{w:instance}{instance}, means, standard deviations, best mean, gaps & Definition~\ref{def:bandit} \\
$\Xr_{a,j}$, $\Arm_t$ & $j$-th reward of arm $a$ (stack model); arm pulled at round $t$ & Definition~\ref{def:bandit} \\
$\Np_a(t)$, $\Reg_T$ & pulls of $a$ in rounds $1,\dots,t$; regret & Definition~\ref{def:bandit} \\
$\Pe_{a,s}$, $\Pe_a(t)$ & empirical \term{w:law}{law} of the first $s$ rewards of $a$; at round $t$ & Definition~\ref{def:bandit} \\
$\Kinf(P,\mu;\mathcal D)$ & information function of a model (strict) & Definition~\ref{def:kinfD} \\
$\Hk$, $\Hlt$ & \term{w:law}{laws} of kurtosis at most $\kappa$; less than $\kappa$ & Definition~\ref{def:class} \\
$\Kc(P,\mu)$, $\Kge(P,\mu)$ & $\Kinf(P,\mu;\Hk)$; its non-strict version & Definition~\ref{def:kinf} \\
$\phik$ & $y^4-2\kappa y^2+\kappa$ & Definition~\ref{def:normalised} \\
$\Gk$, $\KG$ & normalised class; its information function & Definition~\ref{def:normalised} \\
$\tauc_{c,s}$ & $x\mapsto(x-c)/s$, and image \term{w:law}{laws} & Definition~\ref{def:normalised} \\
$\hl_\lambda$, $\gam(y)$ & $1+\lambda_1y+\lambda_2\phik(y)$; $(y,\phik(y))$ & Definition~\ref{def:dualset}, Appendix~\ref{sec:proof-dual} \\
$\Lam$, $\Lw$ & dual set; its half-width & Definition~\ref{def:dualset} \\
$\qk$, $\Kinfty(\kappa)$ & largest atom; value at infinity $\log(1/\qk)$ & Definition~\ref{def:kinfty} \\
$\Cst(\nu)$ & first-order constant $\sum_a\gap_a/\Kc(\nu_a,\mstar)$ & \eqref{eq:cstar} \\
$\thr(t,n)$ & threshold & \eqref{eq:threshold} \\
$\Ix_a(t)$ & index of \algo{} & \eqref{eq:index} \\
$\Ka$, $\LT$, $\nT(\varepsilon)$, $\kep$ & $\Kc(\nu_a,\mstar)$; $\log T+2\log\log T+1$; largest $n$ with $n(\Ka-\varepsilon)\le\thr(T,n)$; $\Ka-\varepsilon$ & \eqref{eq:Ka}, Theorem~\ref{thm:main} \\
$\Ca(\varepsilon)$, $(\delta,J,\eta)$ & existential constant; parameters of the lower deviation & Lemma~\ref{lem:termC}, Theorem~\ref{thm:deviation} \\
\bottomrule
\end{tabular}
\end{center}

\bibliographystyle{plainnat}
\bibliography{refs,../../papers/library}

\end{document}
